\documentclass[11pt,reqno]{amsart}
\usepackage[T1]{fontenc}
\usepackage[utf8]{inputenc}
\usepackage{lmodern}
\usepackage{amsmath,amssymb,mathrsfs}
\usepackage[margin=30mm]{geometry}
\usepackage{microtype}
\usepackage{xurl}
\usepackage{enumitem}
\usepackage{tikz-cd}
\usepackage{xcolor}
\usepackage{hyperref}
\hypersetup{colorlinks=true,linkcolor=blue!50!black,citecolor=blue!50!black,urlcolor=blue!50!black,
pdftitle={Motivic Adams isomorphism and tom Dieck splitting in characteristic zero},pdfauthor={Mengzhudong Feng}}
\setlength{\emergencystretch}{2em}
\setlist[enumerate]{label=(\roman*),leftmargin=2.5em}
\newtheorem{theorem}{Theorem}[section]
\newtheorem{proposition}[theorem]{Proposition}
\newtheorem{lemma}[theorem]{Lemma}
\newtheorem{corollary}[theorem]{Corollary}
\theoremstyle{definition}
\newtheorem{definition}[theorem]{Definition}
\newtheorem{example}[theorem]{Example}
\theoremstyle{remark}
\newtheorem{remark}[theorem]{Remark}
\theoremstyle{plain}
\newtheorem*{maintheorem}{Theorem A}
\newtheorem*{layertheorem}{Theorem B}
\newtheorem*{charptheorem}{Theorem D}
\newtheorem*{splittingtheorem}{Theorem C}
\numberwithin{equation}{section}
\newcommand{\SH}{\mathrm{SH}}
\newcommand{\Sp}{\mathrm{Sp}}
\newcommand{\one}{\mathbf 1}
\newcommand{\A}{\mathbb A}
\newcommand{\Pj}{\mathbb P}
\newcommand{\Gm}{\mathbb G_m}
\newcommand{\Sm}{\mathrm{Sm}}
\newcommand{\lis}{\mathrm{lis}}
\newcommand{\free}{\mathrm{free}}
\newcommand{\et}{\mathrm{\acute et}}
\newcommand{\Hom}{\underline{\operatorname{Hom}}}
\DeclareMathOperator{\Map}{Map}
\DeclareMathOperator{\Fun}{Fun}
\DeclareMathOperator{\Nat}{Nat}
\DeclareMathOperator{\Ind}{Ind}
\DeclareMathOperator{\Lan}{Lan}
\DeclareMathOperator{\Spec}{Spec}
\DeclareMathOperator{\Gal}{Gal}
\DeclareMathOperator{\Th}{Th}
\DeclareMathOperator{\cofib}{cofib}
\DeclareMathOperator{\fib}{fib}
\DeclareMathOperator*{\colim}{colim}
\DeclareMathOperator*{\limm}{lim}
\newcommand{\Nm}{\operatorname{Nm}}
\newcommand{\incl}{\iota}
\newcommand{\cont}{\mathrm{cont}}
\newcommand{\C}{\mathcal C}
\newcommand{\D}{\mathcal D}
\newcommand{\F}{\mathcal F}
\newcommand{\X}{\mathcal X}
\newcommand{\Y}{\mathcal Y}
\newcommand{\E}{\mathcal E}
\newcommand{\B}{\mathcal B}
\newcommand{\smashto}{\longrightarrow}
\newcommand{\todo}[1]{\textcolor{red}{[TODO: #1]}}


\title[Motivic Adams Isomorphism]{Motivic Adams isomorphism and tom Dieck splitting in characteristic zero}
\author{Mengzhudong Feng}
\email{research@timfeng.com}
\date{Version 3, October 28, 2026}

\begin{document}
\begin{abstract}
We prove a motivic Adams isomorphism for smooth affine
linearly reductive groups over fields of characteristic zero.
The diagonal of the classifying stack defines an invertible twist.
We calculate the twist for finite groups and split tori, and show why
it need not be a constant sphere for a split reductive group.
We show that geometric fixed points identify each isotropy layer with the lisse
motivic category of its Weyl group.
Furthermore, over an algebraically closed field, a normalizer transfer and a
isotropy filtration argument give a twisted tom Dieck splitting
for split spectra, including suspension spectra.
\end{abstract}
\maketitle
\setcounter{tocdepth}{1}
\tableofcontents
\section{Introduction}\label{sec:introduction}

Equivariant homotopy theory studies an object together with its group action.
Two fundamental constructions in equivariant homotopy theory are orbits and fixed points.
An Adams isomorphism relates these constructions on the objects with free actions.
Whereas a tom Dieck splitting goes further, it identifies the fixed points of a
suspension spectrum with a sum of contributions from subgroups.
The motivic versions consider objects in equivariant algebraic geometry, 
and hence account for torsors and for the geometry of
the group schemes.

For finite groups of order invertible on the base, Gepner and Heller
proved the motivic Adams isomorphism and the tom Dieck splitting theorem
\cite[Theorems~6.33 and 7.4]{GH}. 
The main tools used in their paper include isotropy filtrations, duality, 
and ambidexterity for finite \'etale maps. 

Following their foundational work, the present paper extends their result to 
smooth affine linearly reductive groups in characteristic zero. 
These include connected reductive groups and
their finite extensions. Nonsplit groups are allowed.

For local systems of spectra, the continuous approximation to cohomology
is homology twisted by the Spivak--Klein dualizing spectrum
\cite[Theorem~I.4.1 and Definition~I.4.2]{NikolausScholze}.
Our diagonal twist and continuous Adams map have the same formal shape.
The proofs use the motivic duality and descent arguments.
For earlier topological constructions, see
\cite[Theorem~D and Definition~1.3]{KleinDualizing} and
\cite[Definition~5.2.2]{RognesDualizable}.

There are two motivic categories in the paper.
The lisse category \(\SH(BG_S)\) describes spectra on the classifying
stack of \(G\)-torsors over \(S\), which is the motivic analogue of the Borel/homotopy category.
The genuine category \(\SH^G(S)\) also records isotropy and inverts
representation spheres.
Its free subcategory identifies with \(\SH(BG_S)\).
This identification (and its companions) plays a core role in the paper.
Namely, we use them to pass from a lisse statement to a genuine statement.  

\subsection{Lisse Adams and its twist}

Let \(k\) be a field of characteristic zero, let \(G/k\) be smooth,
affine, and linearly reductive.
Let \(S\) be a quasi-compact,
quasi-separated \(k\)-scheme of finite type with a trivial \(G\)-action.
Write
\[
 B=BG_S=[S/G],\qquad f:B\longrightarrow S.
\]
In the lisse category \( \SH(BG_S) \), smooth pullback has adjoints
\[
 f_\sharp\dashv f^*\dashv f_*.
\]
We call \(f_\sharp\) homotopy orbits and \(f_*\) homotopy fixed points.
Ordinary direct image need not preserve colimits.
We therefore use its continuous approximation \(f_*^{\cont}\).
This is the colimit-preserving extension of \(f_*\) from compact objects.
See Definition~\ref{def:continuous} for the precise construction.

Let \(\Delta:B\to B\times_S B\) be the diagonal, with projections
\(q_1,q_2\). The diagonal determines the coefficient
\begin{equation}\label{eq:intro-mu}
 \mu_f=q_{1*}\Delta_\sharp\one_B.
\end{equation}
The counit \(q_1^*q_{1*}\Delta_\sharp\one\to\Delta_\sharp\one\)
then defines a norm. Its explicit construction is in
Section~\ref{sec:norm}.

\begin{maintheorem}[Motivic Adams]
The object \(\mu_f\) is invertible. The diagonal norm induces an equivalence
\[
 \Nm_f:f_\sharp(M\otimes\mu_f)
       \xrightarrow{\ \simeq\ }f_*^{\cont}M
 \qquad(M\in\SH(BG_S)).
\]
If \(M\) is compact, its composite with \(f_*^{\cont}M\to f_*M\)
is an equivalence as well.
If \(S\) is smooth and quasiprojective over \(k\), there is a compatible
equivalence
\[
 f_*^{\cont}M\simeq(i_!M)^G,
\]
where \(i_!:\SH(BG_S)\to\SH^G(S)\) is the free embedding and
\((-)^G\) denotes genuine categorical fixed points.
\end{maintheorem}

The three assertions are proved in
Theorem~\ref{thm:new-mu-invertible}, Theorem~\ref{thm:lisse-adams},
and Corollary~\ref{cor:genuine-adams-repaired}.
Invertibility comes from a duality argument.
For smooth maps \(l:C\to E\) and \(p:E\to C\) with \(pl=\mathrm{id}\),
suppose that \(L=l_\sharp\one\) is dualizable.
Then \(K=p_*L\) satisfies
\[
 L\simeq DL\otimes p^*K,
 \qquad K^{-1}\simeq p_\sharp DL.
\]
Applying this to the diagonal gives the claimed invertibility.
This duality equivalence will also help identify the
specified norm map \(\Nm_f\).

The proof of motivic Adams uses a theorem of Lucchini Arteche
\cite[Corollary~18]{LA}.
It gives a finite subgroup \(H\subset G\) such that every
\(G\)-torsor over every extension field of \(k\) comes from a \(H\)-torsor.
We explain why \(BH_S\to BG_S\) then has Nisnevich local sections.
This gives a conservative pullback and a collection of compact generators in \(\SH(BG_S)\).
The finite group motivic Adams theorem applies to these generators.
The general motivic Adams then follows from the finite group case 
by a calculation with the graph of \(BH_S\to BG_S\).

For a finite \'etale group, \(\mu_f=\one\).
For a split torus of rank \(r\), \(\mu_f=S^{r,r}\).
For a split connected reductive group of dimension \(n\) and rank \(r\),
the pullback of \(\mu_f\) to the trivial torsor (i.e. forgetting the action) is \(S^{n,(n+r)/2}\).
Note that this does not mean that \(\mu_f\) itself is a sphere with trivial action.
For example, take the frame torsor of \(\mathcal O\oplus\mathcal O(1)\)
on \(\Pj^1_{\mathbb C}\). Pullback to this torsor gives
\[
 \mu_{GL_2}|_{\Pj^1}
   \simeq S^{2,2}\otimes\Th(\mathcal O(1)),
\]
which is not \(S^{4,3}\).
Section~\ref{sec:new-twist} proves these calculations and also gives
an example for a nonsplit torus.

\subsection{Isotropy layers}

Now assume that \(S/k\) is smooth and quasiprojective.
A geometric stabilizer of a smooth \(G\)-scheme need not be reductive.
In characteristic zero it has a Levi subgroup, unique up to conjugacy
by its unipotent radical.
We filter \(\SH^G(S)\) by these embedded Levi classes.
It is crucial that this is compatible with equivariant affine bundle invariance.

Fix a linearly reductive \(k\)-subgroup \(H\subset G\), and put
\[
 N_H=N_G(H),\qquad W_H=N_H/H.
\]
Let \(e_{\le H}\) be the idempotent for linearly reductive classes
subconjugate to \(H\), and let \(e_{<H}\) remove the class of \(H\).
The object
\[
 e_H=e_{\le H}\otimes\cofib(e_{<H}\to\one)
\]
selects the corresponding layer of \(H\).

The geometric fixed-point functor \(\Phi^H\) first restricts to \(N_H\), then
takes geometric \(H\)-fixed points. Restricting \(\Phi^H\) to the
lisse \(W_H\)-category \(\SH(BW_H)\) gives a functor \(\Phi_{\lis}^H\).

\begin{layertheorem}[Single layer identification]
There is an equivalence
\[
 \Phi_{\lis}^H:e_H\otimes\SH^G(S)
             \xrightarrow{\ \simeq\ }\SH(BW_{H,S}).
\]
An inverse sends \(M\) to
\[
 e_H\otimes\operatorname{Ind}_{N_H}^G
       \operatorname{inf}_{W_H}^{N_H}i_{W_H!}M.
\]
If \(k\) is algebraically closed, these layers form an exhaustive
filtration indexed by conjugacy classes of linearly reductive subgroups.
\end{layertheorem}

This is Theorem~\ref{thm:isotropy-layer} and
Proposition~\ref{prop:family-generation}.
For an affine smooth scheme with linearly reductive classes bounded by \(H\),
Luna's slice theorem identifies the maximal isotropy locus with
\(G\times_{N_H}X^H\).
Its complement and punctured normal bundle have strictly smaller isotropy.
Homotopy purity therefore identifies the \(H\)-layer its motive with the \(H\)-layer
of the motive of \(X\).
A Jouanolou replacement extends this argument to the compact generators
of \(\SH^G(S)\).

An exhaustive filtration need not split.
It also does not by itself describe the gluing functors that reconstruct
the entire category from its layers.
While Theorem~B only concerns the single layers,
we also prove an additional splitting statement of the entire category if the required
sections are available.

\subsection{The tom Dieck splitting}

For a genuine spectrum \(X\), categorical fixed points map naturally to
geometric fixed points.
We call \(X\) \emph{\(G\)-split} if, for every \(H\), the map
\[
 (\operatorname{res}_{N_H}^G X)^H
       \longrightarrow \Phi_{N_H}^H\operatorname{res}_{N_H}^G X
\]
has a specified section in the genuine \(W_H\)-category.
Suspension spectra have canonical such sections, obtained from the
unstable map \(X^H \to X\).

\begin{splittingtheorem}[motivic tom Dieck splitting]
Assume that \(k\) is algebraically closed of characteristic zero.
For a \(G\)-split spectrum \(X\), there is an equivalence
\[
 \bigoplus_{(H)}
 \bigl(\Phi_{\lis}^H X\otimes\mu_{W_H}\bigr)_{hW_H}
       \xrightarrow{\ \simeq\ } X^G.
\]
Here \(M_{hW_H}=(f_{W_H})_\sharp M\), where
\(f_{W_H}:BW_{H,S}\to S\).
The sum ranges over conjugacy classes of linearly reductive subgroups.
The equivalence is natural in maps preserving the chosen sections.
In particular, it gives a canonical splitting for suspension spectra.
\end{splittingtheorem}

See Theorem~\ref{thm:td-general} and Corollary~\ref{cor:td-suspensions}.
The proof has two separate steps.
First, a normalizer transfer turns the chosen fixed point sections into
sections of the actual maps from a supported piece to its layer.
Second, a general theorem for well-founded filtrations assembles these
sections into a direct sum.
The Adams isomorphism identifies each summand with the displayed
twisted homotopy orbits.
The transfer is required only on the single layer. We do not assume an
equivalence between induction and coinduction on every object.

For split diagonalizable \(G\), the same splitting holds over every
characteristic zero field. In that case the geometric subgroups descend
through the character lattice, and every normalizer equals \(G\).
For a general nonsplit group, a full arithmetic indexing theorem would
also need descent data for geometric subgroup classes.

\subsection{Positive characteristic and the scope of the paper}

The section-duality argument has no characteristic restriction.
In characteristic \(p>0\), a smooth affine linearly reductive group is
an extension of a finite etale group by a torus.
An invariant toric compactification gives the following result.

\begin{charptheorem}[Invertibility]
Let \(k\) have characteristic \(p>0\), let \(G/k\) be smooth affine and
linearly reductive, and let \(S/k\) be smooth and quasiprojective.
Then \(\mu_f\) is invertible in \(\SH(BG_S)\), with inverse
\(q_{1\sharp}D(\Delta_\sharp\one)\).
\end{charptheorem}

Appendix~\ref{subsec:new-charp-twist} proves Theorem~D.
It establishes the geometric hypothesis for twist invertibility.
It does not establish a general positive-characteristic Adams
isomorphism or isotropy splitting.
The characteristic-zero proofs use resolution, finite torsor detection,
and Levi subgroups at separate points.
These uses should not be conflated with invertibility of the twist.

In a sequel, we will treat the Adams isomorphism and tom Dieck splitting
over arbitrary fields and in positive characteristic.
The extension will use fixed points indexed by gerbes,
following the approach of Bachmann
\cite[Section~1.2.2]{bac22}.
We will group geometric stabilizer classes into Galois orbits and
describe the corresponding isotropy layers over Weyl gerbes.
These gerbes retain the descent data even when a geometric conjugacy
class has no subgroup representative over its base field.


\subsection*{Organization and conventions}

Section~\ref{sec:formalism} gives the categorical foundations and
the section lemma.
Section~\ref{sec:new-twist} constructs and calculates the twist.
Sections~\ref{sec:reduction} and~\ref{sec:norm} prove the Adams theorem;
Section~\ref{sec:genuine-comparison} gives its genuine form.
Sections~\ref{sec:isotropy} and~\ref{sec:one-layer} construct the
isotropy filtration and its layers.
Sections~\ref{sec:splitting} and~\ref{sec:td-splitting} prove the splitting.
Appendix~\ref{sec:mixed-bc} records the mixed base-change consequence.

The ground field has characteristic zero unless stated otherwise.
Groups are smooth affine of finite type; \(G\) is linearly reductive.
This means that taking invariants is exact on its representations.
The lisse base is quasi-compact and quasi-separated of finite type.
Whenever genuine categories occur, the base is smooth quasiprojective
with trivial group action.
Fiber products of stacks are 2-fiber products.
We use right torsors, and write \(P\times^G V\) with the convention
\((pg,v)\sim(p,gv)\).
Categories, adjunctions, and mapping objects are understood in the
stable \(\infty\)-categorical sense.

The foundational operations are those of Hoyois \cite{Hoyois} and
the lisse extensions developed in
\cite{KhanRavi,KhanRaviStacks,ChowdhuryDAngelo,KhanLisse}.

Here is a summary of the notation used throughout the paper.
\begin{center}
\renewcommand{\arraystretch}{1.16}
\begin{tabular}{@{}p{.33\textwidth}p{.62\textwidth}@{}}
\(\SH(\X)\) & Lisse motivic spectra on a stack; ordinary motivic spectra if \(\X\) is a scheme.\\
\(\D_G=\SH^G(S)\) & Genuine equivariant motivic spectra.\\
\(\Map(M,N),\ \Hom(M,N)\) & Mapping spectrum and internal Hom object.\\
\(DM=\Hom(M,\one)\) & Tensor dual; used as a dual when \(M\) is dualizable.\\
\(\Th(V),\ S^V\) & Thom object of a bundle or representation, respectively.\\
\(S^{a,b}\) & \((S^{1,0})^{\otimes(a-b)}\otimes(\Gm,1)^{\otimes b}\).\\
\(a_\sharp,\ a_*\) & Smooth left adjoint and ordinary right adjoint to \(a^*\).\\
\(a_*^{\cont}\) & Colimit-preserving extension of \(a_*\) from compact objects.\\
\(i_!,i^*,i_*\) & Free embedding, restriction, and cofree embedding.\\
\(X^H,\ \Phi^H X\) & Categorical and geometric fixed points of a spectrum.\\
\(e_\F,\ \widetilde e_\F\) & Family colocalization object and \(\cofib(e_\F\to\one)\).\\
\(e_{\le H},e_{<H},e_H\) & Family, strict lower family, and pure layer.\\
\end{tabular}
\end{center}

For a scheme or unstable motivic space, \(X^H\) instead denotes its
fixed scheme or unstable fixed point space.
The distinction will be clear from the context.
For a smooth representable map \(a\), our purity convention is
\[
 a_\sharp \simeq a_!(-\otimes\Th(T_a)).
\]
In particular, a trivial rank-\(d\) bundle has Thom object \(S^{2d,d}\).

\section*{Acknowledgements}
This paper is the first of a series that revises the author's PhD thesis,
completed at the University of Illinois at Urbana--Champaign under
the supervision of Professor Jeremiah Heller.
The series reorganizes the material, simplifies several arguments,
and corrects some errors and fills gaps in the original proofs.
The author thanks Professor Jeremiah Heller for guidance and support during the
doctoral work.

Large language model (LLM) tools were used to assist with exposition,
organization, literature searches, and the development and checking
of mathematical arguments.
The author takes responsibility for the mathematical claims, proofs,
references, and final text.


\section{Motivic categories and adjunctions}\label{sec:formalism}

Throughout, we use \emph{motive} in the sense of stable motivic
homotopy theory. For a smooth morphism of schemes \(p:Y\to S\),
the motive of \(Y\) over \(S\) is
\[
 p_\sharp\one_Y\simeq\Sigma^\infty Y_+\in\SH(S).
\]
This is its stable motivic homotopy type.
For a smooth map \(a:\Y\to\X\) in the lisse formalism, we
likewise call \(a_\sharp\one_\Y\in\SH(\X)\) its relative
motive.


\subsection{The lisse category}
For an algebraic stack \(\X\), its lisse motivic category is the limit
\begin{equation}\label{eq:lisse}
 \SH(\X)=
\limm_{(U\to\X)\in\operatorname{Lis}(\X)^{\mathrm{op}}}\SH(U).
\end{equation}
Here \(\operatorname{Lis}(\X)\) has smooth finite type schemes over \(\X\) as objects
 and all morphisms over \(\X\) as arrows.
The connecting morphisms are pullback functors.
An object of \(\SH(\X)\) consists of motivic spectra on smooth schemes, 
with compatible pullback identifications.
This construction agrees with ordinary \(\SH\) on schemes.
We use the comparison and descent results of Khan and Ravi
\cite[Section~1 and Lemma~1.5]{KhanRavi}.

For quotient stacks (more generally, algebraic stacks admitting Nisnevich-local sections),
it also has a presentation by the smooth scheme site.
Indeed, one imposes NL descent and affine-bundle invariance, point the category, and invert \(\Pj^1\).
This is \cite[Theorem~3.21]{ChowdhuryDAngelo}.

For the smooth maps in this paper, including \(BG_S\to S\), we use
\begin{equation}\label{eq:smooth-adjoints}
 a_\sharp\dashv a^*\dashv a_*.
\end{equation}
Pullback is symmetric monoidal.
The smooth projection formula is
\begin{equation}\label{eq:smooth-projection}
 a_\sharp(M\otimes a^*N)\simeq a_\sharp M\otimes N.
\end{equation}
The nonrepresentable smooth case of the above formula is covered by Chowdhury--D'Angelo
\cite[Theorem~4.26]{ChowdhuryDAngelo}.
For representable maps we also use proper pushforward, localization, smooth purity, and proper excision.
These properties can be checked on scheme charts.
We use the Nisnevich-lisse formalism in \cite[Corollary~7.2.4]{KhanLisse} and
the scheme results in \cite[Theorems~2.35, 2.45, and 2.48]{Khan}.

For clarity, consider a Cartesian square
\begin{equation}\label{eq:bc-square}
\begin{array}{ccc}
 Y'&\xrightarrow{b'}&Y\\
 {\scriptstyle a'}\downarrow&&\downarrow{\scriptstyle a}\\
 X'&\xrightarrow{b}&X .
\end{array}
\end{equation}
If \(a\) is smooth, left base change gives
\[
 b^*a_\sharp\simeq a'_\sharp b'^*.
\]
If \(b\) is smooth, smooth base change for direct image gives
\[
 b^*a_*\simeq a'_*b'^*.
\]
These are the canonical exchange maps.
They follow from \cite[Proposition~C.1(iii)]{ChowdhuryDAngelo}.

\subsection{The genuine category}
The genuine category \(\SH^G(S)\) is built on smooth equivariantly
quasiprojective \(G\)-schemes over \(S\).
One imposes equivariant Nisnevich excision and invariance under torsors
under equivariant vector bundles (i.e., equivariant affine bundles), 
then invert equivariant Thom spheres. See \cite[Sections~3 and 6]{Hoyois}.
An equivariant Nisnevich distinguished square consists of an invariant
open immersion and an equivariant etale map that is an isomorphism
over its closed complement.
The category has compact generators given by smooth suspension spectra
and representation sphere desuspensions.
It has the representable six operations used below.
We will construct its relation to the lisse category in
Section~\ref{sec:genuine-comparison}.

Restriction along a subgroup takes compact generators to compact objects.
Inflation along a quotient does also.
Their right adjoints therefore preserve filtered colimits.
They are exact, so they preserve all colimits.
This observation applies to coinduction and to genuine categorical
fixed points. Note that it does not claim continuity of lisse \(f_*\).

\subsection{Compact generation in the lisse category}
An object \(C\) is compact if \(\Map(C,-)\) preserves filtered colimits.
A set of objects generates a stable category if vanishing of maps from
all their shifts detects the zero object.
For a scheme \(U\), the motivic category \(\SH(U)\) is compactly generated by
\[
 \Sigma^{m,n}v_\sharp\one_V,\qquad v:V\to U\text{ smooth}.
\]
For the lisse motivic category \(\SH(\X)\), take a smooth scheme chart \(j:U\to\X\). 
Then the objects \(j_\sharp C\), with \(C\) ranging over the above generators, are compact.
Indeed, \(j^*\) preserves colimits, so adjunction gives
\[
 \Map(j_\sharp C,\colim_\alpha M_\alpha)
 \simeq\colim_\alpha\Map(C,j^*M_\alpha)
\]
for every filtered diagram.
Pullbacks to all smooth charts are jointly conservative by \eqref{eq:lisse}.
The objects \(j_\sharp C\) therefore generate \(\SH(\X)\).
In particular, this shows compact generation for the quotient stacks used below.

One may instead use a single smooth atlas with Nisnevich local sections.
For \(BG_S\), choose a faithful representation \(G\hookrightarrow GL_m\).
The map
\[
 S\times_k(GL_m/G)\longrightarrow BG_S
\]
admits a Nisnevich local section.
Its pullback along a \(G\)-torsor \(P\to V\) is the induced \(GL_m\)-torsor.
Such a torsor is Zariski locally trivial.
Descent along this atlas is precisely the form of \cite[Lemma~1.5]{KhanRavi} that we need.

\begin{lemma}\label{lem:compact-pullback}
Let \(a:\Y\to\X\) be smooth and representable by schemes.
Then \(a^*\) preserves compact objects.
The same is true of \(a_\sharp\).
If \(a^*\) is conservative and \(C_i\) are compact generators of \(\SH(\Y)\),
then \(a_\sharp C_i\) are compact generators of \(\SH(\X)\).
\end{lemma}
\begin{proof}
Pull a compact generator (smooth scheme) on \(\X\) back along \(a\).
Smooth base change shows that the pullback is the smooth pushforward of a compact object on \(\Y\).
Therefore \(a^*\) preserves the compact generators above and hence all compact objects.
This proves the first assertion.
The second follows from \(a_\sharp\dashv a^*\) and preservation of filtered colimits by \(a^*\).
For the last assertion, suppose that every \(\Map(a_\sharp C_i,M)\), including all shifts, vanishes.
Adjunction implies \(a^*M=0\).
Conservativity then gives \(M=0\).
\end{proof}

The representability assumption in the first assertion matters.
We do not claim that the pullback \(f^* : \SH(S)\to\SH(BG_S)\) preserves compact objects.

\subsection{Continuous approximation}
\begin{definition}\label{def:continuous}
Let \(T:\C\to\D\) be an exact functor between presentable stable categories, where \(\C\) is compactly generated.
For the inclusion \(u:\C^\omega\hookrightarrow\C\), define
\[
 T^{\cont}=\Lan_u(T|_{\C^\omega}).
\]
The counit of left Kan extension is a comparison \(T^{\cont}\to T\).
\end{definition}

\begin{lemma}\label{lem:continuous}
The functor \(T^{\cont}\) preserves colimits.
The comparison is an equivalence on compact objects.
If \(Q:\C\to\D\) preserves colimits, then
\[
 \Nat(Q,T^{\cont})\simeq\Nat(Q,T).
\]
\end{lemma}
\begin{proof}
The restriction of \(T\) to \(\C^\omega\) preserves finite colimits.
The equivalence \(\C\simeq\Ind(\C^\omega)\) extends it to a functor preserving filtered colimits and finite colimits.
This extension preserves all colimits and is the displayed left Kan extension.
Since \(u\) is fully faithful, this extension agrees with \(T\) on compact objects.
Both sides of the displayed equivalence are identified with
\(\Nat(Q|_{\C^\omega},T|_{\C^\omega})\) by the universal property of left Kan extension.
\end{proof}

We apply this to \(T=f_*\).
Then \(f_*^{\cont}\) is determined by ordinary direct image on compact objects.
A natural transformation between colimit preserving exact functors is an equivalence
if and only if it is an equivalence on compact generators.

\subsection{A duality lemma}
A dual of an object \(M\) consists of an object \(DM\), an evaluation
\(DM\otimes M\to\one\), and a coevaluation
\(\one\to M\otimes DM\), satisfying the two triangle identities.
An object with such data is \emph{dualizable}.
An object is \emph{invertible} if it has a tensor inverse.
Every invertible object is dualizable, but dualizability alone is weaker.
The next lemma will be used with smooth pushforward and pullback.

Let
\[
 F:\C\rightleftarrows\D:U,\qquad U\dashv R\dashv V,
\]
so that \(F\dashv U\dashv R\dashv V\).
Assume \(U\) is symmetric monoidal and assume the following projection formula holds.
\[
 F(A\otimes UB)\simeq FA\otimes B.
\]
The adjunction gives
\begin{equation}\label{eq:closed-adjunction}
 \Hom(FA,B)\simeq R\Hom(A,UB).
\end{equation}
This can be verified by mapping an arbitrary object into both sides.

\begin{lemma}\label{lem:abstract-duality}
If \(A\) and \(FA\) are dualizable, then
\[
 FA\otimes B\simeq R(A\otimes VB)
 \qquad(B\in\D).
\]
\end{lemma}
\begin{proof}
For \(Z\in\D\), duality, \eqref{eq:closed-adjunction}, and adjunction give
\begin{align*}
 \Map(Z,FA\otimes B)
 &\simeq\Map(\Hom(FA,Z),B)\\
 &\simeq\Map(R\Hom(A,UZ),B)\\
 &\simeq\Map(\Hom(A,UZ),VB)\\
 &\simeq\Map(UZ,A\otimes VB)\\
 &\simeq\Map(Z,R(A\otimes VB)).
\end{align*}
These identifications are natural in \(Z\), which proves the assertion.
\end{proof}

\begin{lemma}[Duality and invertibility for a section]\label{lem:section}
Let \(l:C\to E\) and \(p:E\to C\) be smooth maps with \(pl=\mathrm{id}_C\).
Suppose
\[
 L=l_\sharp\one_C
\]
is dualizable.
Put \(K=p_*L\).
There is an equivalence \(b:L\simeq l_*K\) such that \(p_*b=\mathrm{id}_K\).
If \(\epsilon:p^*K\to L\) is the counit, and if \(\eta_l\) is the unit 
for the adjunction \(l^* \dashv l_*\), then
\begin{equation}\label{eq:section-unit}
 b\epsilon=\eta_l:p^*K\longrightarrow l_*K.
\end{equation}
The object \(K\) is invertible, with
\begin{equation}\label{eq:new-section-inverse}
 K^{-1}\simeq p_\sharp DL.
\end{equation}
Moreover, \(l_*\) has a right adjoint \(l^\sharp\), and
\(l^\sharp\one_E\simeq K\).
For every \(N\in\SH(C)\), tensoring \(b\) gives
\begin{equation}\label{eq:section-coefficients}
 b_N:l_\sharp N\xrightarrow{\ \sim\ }l_*(K\otimes N).
\end{equation}
Under these identifications, \(b_N(\epsilon\otimes p^*N)\) is \(\eta_l\) at \(p^*(K\otimes N)\).
\end{lemma}
\begin{proof}
Projection and \eqref{eq:closed-adjunction} give
\begin{equation}\label{eq:section-direct-image}
 l_\sharp N=L\otimes p^*N,\qquad
 l_*N=\Hom(L,p^*N)=DL\otimes p^*N.
\end{equation}
Thus \(l_*\) preserves colimits and has a right adjoint \(l^\sharp\).
Apply Lemma~\ref{lem:abstract-duality} to \(A=\one_C\), \(B=\one_E\), gives
\[
 L\simeq l_*l^\sharp\one_E.
\]
Apply \(p_*\).
Since \(p_*l_*=\mathrm{id}\), this identifies \(l^\sharp\one_E\) with \(K\).

\eqref{eq:section-unit} correspond to \(\mathrm{id}_K\) under \(p^*\dashv p_*\).
For the right side, this follows from the unit of the composite adjunction associated with \(pl=\mathrm{id}\).

The direct image formula in \eqref{eq:section-direct-image} rewrites \(b\) as
\begin{equation}\label{eq:new-section-self-duality}
 L\simeq DL\otimes p^*K.
\end{equation}
Apply \(p_\sharp\) and its projection formula. Since
\(p_\sharp l_\sharp=\mathrm{id}\), this gives
\[
 \one_C\simeq(p_\sharp DL)\otimes K.
\]
The monoidal structure is symmetric, so \(p_\sharp DL\) is a
two-sided inverse for \(K\). In particular, \(K\) is dualizable.
Finally, tensor \(b\) with \(p^*N\) and use \eqref{eq:section-direct-image}.
The unit in this formula is the tensor of the unit in \eqref{eq:section-unit}.
This follows by taking its \(l^*\dashv l_*\) adjoint.
\end{proof}

\begin{remark}
The functor \(l^\sharp\) is the right adjoint of \(l_*\).
It is not the exceptional inverse image \(l^!\).
The proof establishes its existence in the particular situation where it is used.
\end{remark}


\begin{remark}\label{rem:new-section-sign}
Write \(\sigma_l=l^\sharp\one_E\).
The lemma gives \(\sigma_l=K\), so
\[
 DL\simeq L\otimes p^*(\sigma_l)^{-1}.
\]
\end{remark}


\section{Dualizing twists}
\label{sec:new-twist}\label{sec:duality}

The twist in Theorem A is defined by the diagonal of a classifying stack.
Its invertibility follows if the diagonal motive is dualizable.
We first establish this geometric input, then calculate the twist explicitly in
several examples.

\subsection{Equivariant compactification and duality}

We next prove the geometric input for the section lemma.
The proof keeps the Thom twists and permits singular boundaries.
It applies in the genuine equivariant category of Hoyois \cite{Hoyois}
and in the lisse category of quotient stacks.
In the latter case every map in the argument is representable.

\begin{proposition}\label{prop:geometric-duality}
Let \(k\) have characteristic zero and let \(J\) be a smooth affine linearly reductive \(k\)-group.
Let \(U\) be a smooth \(J\)-quasiprojective \(k\)-scheme.
Then \(\Sigma_J^\infty U_+\) is dualizable in \(\SH^J(k)\).
For \(u:[U/J]\to BJ\), the object \(u_\sharp\one\) is dualizable in \(\SH(BJ)\).
\end{proposition}
\begin{proof}
The two proofs use the same operations.
We give the argument once.

First let \(q:Y\to\Spec k\) be smooth proper and equivariant, and let \(C\) be dualizable over \(Y\).
Closed adjunction and the proper projection formula give
\[
 \Hom(q_\sharp C,M)
 \simeq q_*\Hom(C,q^*M)
 \simeq q_*DC\otimes M.
\]
This identifies a dual of \(q_\sharp C\).
Purity gives \(q_*C=q_\sharp(C\otimes\Th(-T_q))\), so \(q_*C\) is also dualizable.
For the genuine category these properties are
\cite[Theorems~6.9 and 6.18]{Hoyois}.

Let \(z\) be the structure map of \(Z\). 
We claim that \(z_*\Th(V)\) is dualizable whenever \(Z\) is proper and \(J\)-quasiprojective
and \(V\) is an equivariant vector bundle on \(Z\).
Nilpotent invariance reduces to \(Z\) reduced.
We induct on \(\dim Z\).
The zero-dimensional case is smooth because \(k\) is perfect.

Use functorial resolution in characteristic zero
\cite[Theorem~1.2.1]{Temkin} to obtain \(\rho:Z'\to Z\).
The resolution is an isomorphism over the regular locus and \(Z'\) is smooth.
Functoriality for smooth morphisms lifts the \(J\)-action to \(Z'\).
Indeed, apply it to the action and projection maps \(J\times Z\to Z\).
Compatibility with the group law follows on \(J\times J\times Z\).

Let \(C\) be the reduced singular locus and \(C'=Z'\times_Z C\).
Both \(C\) and \((C')_{\mathrm{red}}\) have smaller dimension.
For \(C'\), this holds because each component of \(Z'\) dominates a component of \(Z\),
and the intersection of each component with \(C'\) is a proper closed subset 
(e.g., the generic point of each component is not in the intersection).

Proper excision gives an exact sequence
\begin{equation}\label{eq:proper-excision}
\begin{split}
 z_*\Th(V)\longrightarrow{}&
 z'_*\Th(\rho^*V)\oplus c_*\Th(V|_C)\\
 \longrightarrow{}&c'_*\Th(V|_{C'}).
\end{split}
\end{equation}
Here the structure maps to the base are denoted by \(z,z',c,c'\).
In the genuine category, this is equivariant proper excision \cite[Corollary~6.25]{Hoyois}.
In the lisse category it follows on scheme charts from \cite[Theorem~2.48]{Khan}
and proper base change.
The first summand is dualizable by the smooth proper case.
The \(c\) and \(c'\) terms are dualizable by induction.
Dualizable objects are closed under finite fibers, so the claim follows.

Choose an equivariant projective compactification of \(U\).
Resolve it without changing the smooth open \(U\).
We obtain \(j:U\hookrightarrow\overline U\), with \(\overline U\) smooth projective,
and an invariant closed complement \(i:D\hookrightarrow\overline U\).
Let \(q:\overline U\to\Spec k\) and \(d=qi\).
Apply localization to \(\Th(T_q)\), then apply the functor \(q_*=q_!\).
Since \(j^*\Th(T_q)=\Th(T_u)\), the result is
\begin{equation}\label{eq:open-duality}
 u_\sharp\one\longrightarrow q_*\Th(T_q)\longrightarrow d_*i^*\Th(T_q).
\end{equation}
The middle term is dualizable by smooth proper duality.
The last term is dualizable by the claim, applied to the restricted vector bundle.
The first term is therefore dualizable.
\end{proof}

For the applications below \(U\) is affine, so the compactification has an elementary construction.
A finite set of generators of \(k[U]\) lies in a finite-dimensional \(J\)-stable subspace.
It gives an equivariant closed immersion into a representation.
Take the closure in the corresponding projective representation after adding a trivial coordinate.
Resolution leaves \(U\) unchanged.
For an equivariantly quasiprojective \(U\), its equivariant projective immersion supplies the same construction.

\begin{corollary}\label{cor:graph-dual}
Let \(H\subset G\) be a finite subgroup.
For \(r:BH_S\to BG_S\), the object \(r_\sharp\one\) is dualizable.
The same holds for the motive of the graph
\[
 l:BH_S\longrightarrow BH_S\times_S BG_S.
\]
\end{corollary}
\begin{proof}
In characteristic zero the variety \(G/H\) is smooth affine, and \([G/H\,/G]=BH\).
Thus Proposition~\ref{prop:geometric-duality} gives the first claim over \(k\).
The graph over \(k\) is represented by the affine \(H\times G\)-variety \(G\),
with action \((a,g)x=gxa^{-1}\).
The proposition gives the second claim.
Smooth left base change identifies both motives over \(S\) with pullbacks of these motives over \(k\).
Symmetric monoidal pullback preserves duality.
\end{proof}


\begin{corollary}\label{cor:absolute-compact-duality}
Every compact object of \(\SH^G(k)\) is dualizable.
For every linearly reductive subgroup \(K\subset G\), the genuine
and lisse motives of \(G/K\) are dualizable.
These duality data remain valid after pullback to \(S\).
\end{corollary}
\begin{proof}
The compact generators over \(k\) are (representation desuspensions
of) suspension spectra of smooth schemes.
Proposition~\ref{prop:geometric-duality} makes them dualizable.
Compact objects are their thick closure, and dualizable objects are
closed under finite cofibers and retracts.
The homogeneous space \(G/K\) is smooth affine for reductive \(K\),
so the proposition applies to it as well.
Pullback to \(S\) preserves duality.
\end{proof}

Note that this assertion about all compact objects concerns the absolute base
\(\Spec k\). We do not claim it for arbitrary relative compact
objects over \(S\).



\subsection{Application to classifying stacks}
Let \(G\) be a smooth affine group over \(k\), and put \(B=BG_S\), \(f:B\to S\).
Write \(\Delta:B\to B\times_S B\) for the diagonal and \(q_1,q_2\) for the projections.
Define
\begin{equation}\label{eq:new-mu-definition}
 \mu_f=q_{1*}\Delta_\sharp\one_B.
\end{equation}
Both \(\Delta\) and \(q_1\) are smooth, and \(q_1\Delta=\mathrm{id}_B\).

The passage from duality to an invertible object is related to
\cite[Remark~1.4.9]{RaksitTate}.
For stably dualizable topological groups, compare
\cite[Theorems~3.1.4 and~3.3.4]{RognesDualizable}.

\begin{theorem}\label{thm:new-mu-invertible}
If \(\Delta_\sharp\one_B\) is dualizable, then \(\mu_f\) is invertible, with
\begin{equation}\label{eq:new-mu-inverse}
 \mu_f^{-1}\simeq q_{1\sharp}D(\Delta_\sharp\one_B).
\end{equation}
In particular, this holds for every smooth affine linearly reductive group over a field of characteristic zero, 
under the base conventions of this paper.
\end{theorem}
\begin{proof}
Apply Lemma~\ref{lem:section} with \(l=\Delta\), \(p=q_1\).
For the final assertion, apply Proposition~\ref{prop:geometric-duality} to the smooth affine \(G \times G\)-scheme
\(G\), with action \((g,h)\cdot x=gxh^{-1}\).
Since \([G / (G \times G)] \simeq BG\), \(\Delta_{k\sharp} \one_{BG_S} \) is dualizable.
Pullback preserves dualizable objects, so the same conclusion holds over \(S\).
\end{proof}

\begin{remark}\label{rem:new-char-p-invertibility}
The formal theorem has no characteristic restriction.
Theorem~\ref{thm:new-charp-invertibility} proves its geometric hypothesis integrally for smooth
linearly reductive groups in positive characteristic.
It uses an invariant toric compactification, with no restriction on
the torus's splitting degree.
\end{remark}

\begin{corollary}[The torsor formula]\label{cor:new-torsor-duality}
Let \(j:U\to B\) be smooth and let
\[
 l=(\mathrm{id},j):U\to U\times_S B,\qquad p:U\times_S B\to U.
\]
Under Theorem~\ref{thm:new-mu-invertible}'s hypotheses,
\[
 \sigma_l=j^*\mu_f,\qquad
 D(l_\sharp\one_U)\simeq
 l_\sharp\one_U\otimes p^*(j^*\mu_f^{-1}).
\]
In particular, \(j^*\mu_f\) and \(\sigma_l\) are invertible.
\end{corollary}
\begin{proof}
The graph is a base change of \(\Delta\), so its motive is dualizable.
Smooth base change gives
\[
 p_*l_\sharp\one_U
 \simeq j^*q_{1*}\Delta_\sharp\one_B=j^*\mu_f.
\]
Apply Lemma~\ref{lem:section} and Remark~\ref{rem:new-section-sign}.
\end{proof}

\subsection{Cases with a constant sphere}
\begin{proposition}\label{prop:new-finite-torus-twists}
If \(G\) is finite \'etale, then \(\mu_f=\one_B\).
If \(G=T=\Gm^r\) is a split torus, then
\[
 \mu_f\simeq S^{r,r}_B.
\]
The torus formula is global on \(BT_S\), and hence holds on every torsor chart.
\end{proposition}
\begin{proof}
For a finite \'etale group, the diagonal is finite \'etale.
Ambidexterity gives \(\Delta_\sharp=\Delta_*\), and therefore
\(\mu_f=q_{1*}\Delta_*\one=\one\).

For the torus, first consider a line bundle \(V\) on \(Y\).
Let \(A\) be the relative suspension spectrum of its complement of zero.
Purity gives \(A\to\one_Y\to\Th(V)\).
Dualizing and rotating gives
\[
 DA\simeq A\otimes\Th(-V)\otimes S^{1,0}.
\]
The projection formula gives
\[A\otimes\Th(-V) \simeq p_{\sharp}(p^*\Th(-V)) \simeq p_{\sharp}S^{-2, -1} 
\simeq A\otimes S^{-2,-1}.\]
The second equivalence follows from the fact that the tautological section trivializes \(V\)
on the complement of 0.
Thus \(DA\simeq A\otimes S^{-1,-1}\).

We describe the diagonal torsor explicitly.
Let \(\chi_i:T\to\Gm\) be the coordinate characters, and let
\(\mathcal U_i\) be their associated line bundles on \(BT\).
On \(X=BT\times_S BT\), set
\[
 \mathcal L_i=q_1^*\mathcal U_i^\vee\otimes q_2^*\mathcal U_i.
\]
There is an equivalence over \(X\)
\begin{equation}\label{eq:torus-diagonal-lines}
 (BT\xrightarrow{\Delta}X)\simeq
 \bigl(\mathcal L_1^\times\times_X\cdots\times_X
             \mathcal L_r^\times\longrightarrow X\bigr),
\end{equation}
where the superscript \(\times\) means the complement of the zero section.

Use the homogeneous space formula
\(BH\simeq[(G/H)/G]\) of
\cite[Corollary~1.16(a)]{IllusieZheng}, translated from right to left
actions, with \(G=T\times T\) and \(H=\Delta T\).
The quotient homomorphism
\[G \to T, \quad (t_1,t_2)\mapsto t_2t_1^{-1}.\]
has coordinate characters \(\chi_i(t_2)\chi_i(t_1)^{-1}\).
Their associated lines are exactly \(\mathcal L_i\).
Thus the cited general formula specializes to
\eqref{eq:torus-diagonal-lines}.

Smooth base change and the projection formula identify the motive
\(L=\Delta_\sharp\one\) with the tensor product of the motives of
these \(r\) complements of zero.
It follows that \(DL\simeq L\otimes S^{-r,-r}\).
Adjunction gives \(DL=\Delta_*\one\).
Applying \(q_{1*}\) and the projection formula for the sphere coefficient yields
\[
 \one=q_{1*}\Delta_*\one
 \simeq\mu_f\otimes S^{-r,-r}.
\]
\end{proof}

The finite \'etale and split-torus calculations also determine
the twist when the identity component is a central split torus.

\begin{lemma}[A central split torus as identity component]
\label{lem:explicit-central-torus-twist}
Let \(K\) be a smooth affine linearly reductive group over a
field \(k\) of characteristic zero. Suppose that
\(T=K^\circ\) is central in \(K\), and choose an isomorphism
\(T\simeq\Gm^d\).
Then
\[
 \mu_K\simeq S^{d,d}_{BK_S}.
\]
The extension \(1\to T\to K\to K/T\to1\) need not split.
\end{lemma}

\begin{proof}
We will calculate the dual of the diagonal motive and then
apply the first projection's direct image to identify
\(\mu_K\).

Put \(B=BK_S\), let
\(\Delta_K:B\to B\times_S B\) be the diagonal, and write
\[
 \mathcal X=B\times_S B,
 \qquad
 q_1:\mathcal X\to B,
 \qquad
 L=\Delta_{K\sharp}\one_B.
\]
Our first goal is to prove
\[
 DL\simeq L\otimes S^{-d,-d}_{\mathcal X}.
\]

Let \(F=K/T\), which is finite \'etale, and form
\[
 \mathcal Z
 =\mathcal X\times_{BF_S\times_S BF_S}BF_S,
\]
where the map from \(BF_S\) is its diagonal.
The diagonal of \(B\) factors as
\[
 B\xrightarrow{t}\mathcal Z
  \xrightarrow{a}\mathcal X.
\]
The map \(a\) is finite \'etale, since it is a base change
of the diagonal of \(BF_S\).

We next identify \(t\) as a split-torus torsor.
An object of \(\mathcal Z\) consists of two \(K\)-torsors
together with an isomorphism of their associated \(F\)-torsors.
Lifts of this isomorphism to the \(K\)-torsors exist fppf
locally and form a torsor under the corresponding twist of
\(T\) by conjugation. Since \(T\) is central in \(K\), this
twist is canonically \(T_{\mathcal Z}\). Thus \(t\) is a
\(T_{\mathcal Z}\)-torsor.

The chosen coordinates on \(T\) express \(t\) as a fiber
product of complements of zero in \(d\) line bundles on
\(\mathcal Z\). Put \(A=t_\sharp\one_B\).
Smooth base change and the projection formula express \(A\)
as the tensor product of their relative motives.
The line-bundle calculation in the proof of
Proposition~\ref{prop:new-finite-torus-twists} therefore gives
\[
 DA\simeq A\otimes S^{-d,-d}_{\mathcal Z}.
\]
That calculation is compatible with pullback and
isomorphisms of line bundles, so it applies on the lisse
site of \(\mathcal Z\).

Since \(L=a_\sharp A\), closed adjunction and finite
\'etale ambidexterity give
\[
 \begin{aligned}
 DL
 &\simeq a_*DA\\
 &\simeq a_\sharp DA\\
 &\simeq a_\sharp
       \bigl(A\otimes a^*S^{-d,-d}_{\mathcal X}\bigr)\\
 &\simeq L\otimes S^{-d,-d}_{\mathcal X}.
 \end{aligned}
\]
Here finite \'etale ambidexterity is the case of
\cite[Theorem~6.9]{Hoyois} with zero relative tangent bundle,
applied on smooth scheme charts; the last equivalence is
the smooth projection formula.

It remains to recover the twist.
Closed adjunction gives \(DL\simeq\Delta_{K*}\one_B\).
Using \(q_1\Delta_K=\mathrm{id}_B\), the definition of
\(\mu_K\), and the projection formula for the invertible
sphere coefficient, we obtain
\[
 \begin{aligned}
 \one_B
 &\simeq q_{1*}\Delta_{K*}\one_B\\
 &\simeq q_{1*}DL\\
 &\simeq q_{1*}L\otimes S^{-d,-d}_B\\
 &\simeq \mu_K\otimes S^{-d,-d}_B.
 \end{aligned}
\]
Tensoring with \(S^{d,d}_B\) proves the assertion.
\end{proof}


Let \(B_0=T\ltimes U\) be a split connected solvable group, with \(T=\Gm^r\) and \(U\) split unipotent.
The example needed below is a Borel subgroup of a split reductive group.
All operations for \(BB_0\) in the next two results are lisse operations.
We use smooth adjunctions, affine-bundle invariance on charts, and
representable ambidexterity.
Note that we do not need a genuine motivic category for the non-linearly reductive
group \(B_0\).

\begin{lemma}\label{lem:new-borel-twist}
For \(B_0\) as above, \(\mu_{BB_0\to S}=S^{r,r}_{BB_0}\).
\end{lemma}
\begin{proof}
Compare its diagonal motive with the pullback of the diagonal motive of \(BT\).
On a chart classifying two \(B_0\)-torsors, the comparison is induced by
\[
 \operatorname{Isom}_{B_0}(P_1,P_2)
 \longrightarrow
 \operatorname{Isom}_T(P_1/U,P_2/U).
\]
Its fibers are torsors under twists of \(U\).
Split-solvable torsors are Zariski locally trivial.
To see this, first trivialize their torus quotient.
Then apply a unipotent filtration and the vanishing of first cohomology of quasi-coherent sheaves on affine schemes.
After this local trivialization, the displayed comparison is an iterated affine-bundle map.
Homotopy invariance identifies its source and target motives.

The split-torus calculation gives
\[
 D(\Delta_{B_0\sharp}\one)
 \simeq\Delta_{B_0\sharp}\one\otimes S^{-r,-r}.
\]
Apply the first projection's direct image, as in Proposition~\ref{prop:new-finite-torus-twists}.
\end{proof}

\subsection{A proper restriction formula}
\begin{proposition}\label{prop:new-proper-restriction}
Let \(C\) and \(B\) be smooth stacks over \(S\) with smooth diagonals.
Let \(r:C\to B\) be representable, smooth, and proper over \(S\).
Suppose the diagonal motives of \(C\) and \(B\) are dualizable.
For their respective twists,
\begin{equation}\label{eq:new-proper-restriction}
 r^*\mu_B\simeq\mu_C\otimes\Th(T_r).
\end{equation}
\end{proposition}
\begin{proof}
Set
\[
 a=\mathrm{id}_C\times r:C\times_S C\to C\times_S B,\qquad
 l=a\Delta_C,
\]
and let \(p:C\times_S B\to C\) be projection.
Write \(L_C=\Delta_{C\sharp}\one\) and \(L_l=a_\sharp L_C\).
Adjunction and smooth proper ambidexterity give
\[
 DL_l=a_*DL_C
 \simeq a_\sharp\bigl(DL_C\otimes\Th(-T_a)\bigr).
\]
The second equivalence is \(a_*=a_\sharp(\Th(-T_a)\otimes-)\) by \cite[Theorem~6.9]{Hoyois}.
By section duality,
\[
 DL_C=L_C\otimes\operatorname{pr}_1^*\mu_C^{-1},
 \qquad T_a=\operatorname{pr}_2^*T_r.
\]
Smooth projection formula along \(\Delta_C\) identifies
\[
 L_C\otimes\operatorname{pr}_1^*N
 \simeq\Delta_{C\sharp}N
 \simeq L_C\otimes\operatorname{pr}_2^*N
\]
for every \(N\) on \(C\).
Move the tangent coefficient to the first projection factor to obtain
\[
 DL_l\simeq a_{\sharp}(L_C \otimes pr_1^*(\mu_C^{-1} \otimes \Th(-T_r)) )
\simeq L_l\otimes p^*(\mu_C^{-1}\otimes\Th(-T_r)).
\]
where \(p\) is the projection \(p: C \times_S B \to C.\)

Apply \(p_*\).
On the left, \(DL_l=l_*\one\), so \(p_*DL_l=\one\).
The right hand side tangent coefficient is invertible, and hence satisfies the ordinary projection formula.
Smooth base change along \(r\) gives 
\[p_*L_l =
p_* l_{\sharp} r^* \one_B 
\simeq r^* q_{1*} \Delta_{B\sharp} \one_B
= r^*\mu_B .\]
Thus
\[
 \one\simeq r^*\mu_B\otimes\mu_C^{-1}\otimes\Th(-T_r),
\]
which proves the claim.
\end{proof}

\begin{corollary}[A split reductive group]\label{cor:new-split-reductive}
Let \(G\) be split connected reductive, \(n=\dim G\), and \(r=\operatorname{rank}G\).
Choose a split Borel \(B_0\), and write \(b:BB_0\to BG\).
Then
\begin{equation}\label{eq:new-borel-restriction}
 b^*\mu_G\simeq S^{r,r}\otimes\Th(T_b).
\end{equation}
For a \(G\)-torsor on \(X\) with a \(B_0\) reduction \(Q\), 
let \(u: X \to BB_0\) be the map that classifies the reduction \(Q\).
Let \(E\) be its associated vector bundle with fiber 
\(\mathfrak g/\mathfrak b_0\), i.e., \(E = Q \times^{B_0} (\mathfrak{g}/\mathfrak{b}_0)\).
Then the pullback of \(\mu_G\) on \(X\) is \(S^{r,r}\otimes\Th(E)\).
In particular, at the trivial torsor \(e:S\to BG\),
\[
 e^*\mu_G\simeq S^{n,(n+r)/2}.
\]
\end{corollary}
\begin{proof}
The map \(b\) is smooth proper with fiber \(G/B_0\).
Its tangent bundle is associated with \(\mathfrak g/\mathfrak b_0\).
Apply Proposition~\ref{prop:new-proper-restriction} and Lemma~\ref{lem:new-borel-twist}.
The rank of this tangent bundle is \((n-r)/2\).
At the trivial torsor it is a constant vector space.
\end{proof}

This is saying that when ignoring the group action, the twist \(\mu_G\) is a constant sphere.
The following example shows the group action on \(\mu_G\) itself can be nontrivial. 

\begin{example}[A split group with a nontrivial twist]\label{ex:new-gl2-twist}
Take \(k=\mathbb C\), \(X=\Pj^1\). The \(GL_2\) torsor of frames of
\(V=\mathcal O_X\oplus\mathcal O_X(1)\) is 
\[
\operatorname{Fr}(V)
=
\underline{\operatorname{Isom}}_X
\bigl(\mathcal O_X^{\oplus 2},V\bigr).
\]
The line \(\mathcal O_X\subset V\) gives a reduction to the upper triangular Borel \(B_0\).

Let \(Q\subset\operatorname{Fr}(V)\) be the \(B_0\)-torsor of
frames carrying the first coordinate line to
\(\mathcal O_X\subset V\).
The associated relative tangent line bundle is
\[
 E=Q\times^{B_0}(\mathfrak{gl}_2/\mathfrak b_0)
 \simeq \operatorname{Hom}(\mathcal O_X,V/\mathcal O_X)
 \simeq \mathcal O_X(1).
\]
The first isomorphism identifies tangent vectors to the flag
with first-order deformations of the line.
Proposition~\ref{prop:new-proper-restriction} and
Lemma~\ref{lem:new-borel-twist} therefore give
\[
 \mu_{GL_2}|_X
 \simeq S^{2,2}\otimes\Th_X(\mathcal O_X(1)).
\]
This is not \(S^{4,3}\).

Indeed, an equivalence \(\Th_X(\mathcal O_X(1))\simeq S^{2,1}\) would, 
after smooth pushforward along \(X\to k\), 
identify the two absolute Thom spectra.
With \(\mathbb F_2\)-coefficients, each has a unique nonzero class in bidegree \((2,1)\).
Voevodsky's Thom formula states
\[
 \operatorname{Sq}^2(t_L)=c_1(L)t_L
\]
for a line bundle \(L\). See \cite[Corollary~14.4]{VoevodskyPower} with \(\ell=2\).
For \(\mathcal O(1)\), the right side is nonzero by the Thom isomorphism
\cite[Proposition~4.3]{VoevodskyPower} and
\(\operatorname{CH}^1(\Pj^1)/2=\mathbb F_2\).
For the trivial line \(c_1\) is zero.
This gives a contradiction.
\end{example}

\subsection{Spheres for tori}
\begin{proposition}\label{prop:new-commutative-twist}
Suppose \(G\) is commutative and its diagonal motive is dualizable.
Let \(e:S\to BG_S\) be the trivial-torsor section and put
\[
 \omega_G=f_*e_\sharp\one_S.
\]
Then \(\omega_G\) is invertible in \(\SH(S)\), and \(\mu_f\simeq f^*\omega_G\).
\end{proposition}
\begin{proof}
This is the translation argument for a group object. Compare
\cite[Proposition~1.5.3]{RaksitTate}.
Commutativity makes \((g,h)\mapsto(g,g^{-1}h)\) an automorphism of \(G\times G\).
This automorphism preserves \(q_1\) and takes \(\Delta\) to \(\mathrm{id}\times e\).
Smooth base change therefore gives
\[
 \mu_f\simeq q_{1*}\operatorname{pr}_2^*(e_\sharp\one)
 \simeq f^*f_*(e_\sharp\one).
\]
The second equivalence can be checked on every smooth scheme chart \(U\to BG\).
There it is ordinary smooth base change along \(U\to S\).
Such chart pullbacks are jointly conservative in the lisse category.
Finally, \(\omega_G=e^*\mu_f\), so it is invertible.
\end{proof}

For a torus \(T\), \(\omega_T\) is an arithmetic invertible sphere.
Here \emph{sphere} means an element of the Picard group, not a bidegree constant sphere.
For split \(T\) of rank \(r\), it is \(S^{r,r}\).
The formula can fail for nonsplit tori.

\begin{example}\label{ex:new-real-torus}
Let \(T=R^1_{\mathbb C/\mathbb R}\Gm\).
If \(\omega_T=S^{1,1}\), Corollary~\ref{cor:new-torsor-duality} at the trivial torsor would imply
\[
 D\Sigma^\infty T_+
 \simeq\Sigma^\infty T_+\otimes S^{-1,-1}.
\]
Real points realization is the composite of the strong symmetric monoidal
\(C_2\)-equivariant Betti realization of \cite[Proposition~4.8]{HellerOrmsbyGalois}
with geometric fixed points.
It sends \(S^{1,1}\) to \(S^0\).
It sends \(\Sigma^\infty T_+\) to
\(\one\oplus\Sigma\one\), whose dual is
\(\one\oplus\Sigma^{-1}\one\).
Their integral homology differs, which contradicts the displayed equivalence.
\end{example}

The exact Picard class of \(\omega_T\) for an arbitrary torus is a further question.
We do not need it in this paper. 

\section{Finite detection and the finite Adams theorem}\label{sec:reduction}

The next result extends the finite-reduction theorem of
Chernousov--Gille--Reichstein over algebraically closed fields
\cite[Corollary~1.4]{CGR08} to the generality needed here:
the group may be disconnected, and the characteristic-zero
base field need not be algebraically closed. 
It reduces an arbitrary torsor to a finite subgroup,
up to passing to a Nisnevich neighborhood on its base.
It will give the compact generators on which we test the Adams norm.

\subsection{A finite subgroup detecting field torsors}

\begin{theorem}[Lucchini Arteche]\label{thm:finite-subgroup}
There is a finite étale \(k\)-subgroup \(H\subset G\) such that
\[
 H^1(F,H)\longrightarrow H^1(F,G)
\]
is surjective for every field extension \(F/k\).
\end{theorem}
\begin{proof}
This is the characteristic-zero case of \cite[Corollary~18]{LA}.
Its hypothesis is that the identity component of \(G\) is reductive.
It allows disconnected and nonsplit groups.
Our hypotheses imply this condition.
A finite group scheme over a characteristic-zero field is étale.
We do not need the exact upper bound on order of \(H\).
\end{proof}

Here \(H^1(F,G)\) denotes the pointed set of isomorphism classes of fppf \(G\)-torsors over \(F\).
The theorem says that every \(G\)-torsor over every field extension \(F/k\) admits a reduction 
of structure group to the same finite subgroup \(H\), defined over \(k\).
It is well known that, for a scheme \(V\) of finite Krull dimension, the pullback functors
\(\SH(V)\to\SH(\kappa(v))\), \(v\in V\), form a jointly conservative family
\cite[Proposition~B.3]{BachmannHoyoisNorms}.
Applied on smooth scheme charts, this fact and the theorem already imply that pullback 
along \(BH_S\to BG_S\) is conservative.
For completeness, we construct explicit Nisnevich local sections of this map by using smoothness 
to extend the reductions over residue fields to Nisnevich neighborhoods (Proposition~\ref{prop:nis-reduction}).

\begin{lemma}\label{lem:torsor-reduction}
Let \(P\to V\) be a right \(G\)-torsor.
The pullback of \(BH\to BG\) along \(V\to BG\) is \(P/H\).
Sections of \(P/H\to V\) are equivalent to \(H\)-reductions of \(P\).
The scheme \(P/H\) is smooth over \(V\).
\end{lemma}
\begin{proof}
The torsor \(P\) is affine over \(V\).
The free finite quotient exists, and \(P\to P/H\) is an \(H\)-torsor
\cite[Tag~07S7]{Stacks}.
For an \(H\)-torsor \(Q\to V\) and an isomorphism \(Q\times^H G\simeq P\),
the map \(Q\to P/H, q \mapsto (q, 1)\) descends to a section over \(V\) by faithfully flat descent.
Conversely, pull \(P\to P/H\) back along a section.
The resulting \(H\)-torsor \(Q\) satisfies \(Q\times^H G\simeq P\).
This can be checked after trivializing \(Q\).
These constructions commute with base change and are inverse.

After a faithfully flat cover trivializing \(P\), the quotient is \(G/H\times V\).
Smoothness descends, so \(P/H\to V\) is smooth.
\end{proof}

\begin{lemma}\label{lem:nisnevich-slice}
Let \(Y\to V\) be smooth.
If the fiber at \(v\in V\) has a \(\kappa(v)\)-point \(y\),
there is an étale neighborhood \((V',v')\to(V,v)\) with
\(\kappa(v')=\kappa(v)\) and a \(V\)-map \((V',v')\to (Y,y)\).
\end{lemma}
\begin{proof}
This is the étale slice statement \cite[Tag~057G]{Stacks}.
We include the argument here for completeness.
Choose étale coordinates \(\phi: W\to\A^n_{V_0}\) near \(y\), 
where \(V_0\) is an open neighborhood of \(v\).
Such coordinates exist by \cite[Tag~039Q]{Stacks}.
Write the coordinates of \(\phi(y)\) by \((\bar a_1,\ldots,\bar a_n)\in\kappa(v)^n.\)
Lift these coordinates from \(\kappa(v)\) to \(\mathcal O_{V,v}\).
After shrinking \(V_0\), these lifts are regular functions on \(V_0\).
They define a section \(a:V_0\to\A^n_{V_0}\).
Set
\[
 V'=W\times_{\A^n_{V_0},a}V_0.
\]
The map \(V'\to V_0\) is étale.
The compatible pair \((y,v)\) gives a point \(v'\) with residue field \(\kappa(v)\).
Projection to \(W\) gives the required map to \(Y\).
\end{proof}

\begin{proposition}\label{prop:nis-reduction}
Choose \(H\) as in Theorem~\ref{thm:finite-subgroup}.
Then \(r:BH_S\to BG_S\) has Nisnevich local sections after every scheme base change.
In particular, \(r^*\) is conservative.
\end{proposition}
\begin{proof}
For \(V\to BG_S\), Lemma~\ref{lem:torsor-reduction} identifies the pullback with \(P/H\to V\).
For every \(v\in V\), Theorem~\ref{thm:finite-subgroup} reduces the torsor \(P_v\) to \(H\).
Lemma~\ref{lem:torsor-reduction} gives a point in \((P/H)_v(\kappa(v))\).
Lemma~\ref{lem:nisnevich-slice} extends it over a pointed Nisnevich neighborhood.
Choose one such neighborhood for each \(v\).
They form a Nisnevich cover on which \(P/H\to V\) has sections.
Nisnevich descent on scheme charts now shows that \(r^*\) detects equivalences.
\end{proof}

\begin{proposition}\label{prop:tensor-detect}
Let \(A=r_\sharp\one\).
Both \(A\otimes-\) and \(DA\otimes-\) are conservative.
Moreover,
\[
 DA\otimes M\simeq r_*r^*M.
\]
\end{proposition}
\begin{proof}
Smooth base change identifies \(r^*A\) with the motive of
\(BH_S\times_{BG_S}BH_S\to BH_S\).
The diagonal section exhibits \(\one\) as a retract of \(r^*A\).
Thus \(A\otimes M=0\) implies \(r^*A \otimes r^*M = 0\).
So \(r^*M=0\), and then \(M=0\).
This proves conservativity.

By Corollary~\ref{cor:graph-dual}, \(A\) is dualizable.
A duality triangle makes \(A\) a retract of \(A\otimes DA\otimes A\).
Consequently \(DA\otimes M=0\) implies \(A\otimes M=0\), proving the second conservativity assertion.
Finally adjunction gives
\[
 DA\otimes M=\Hom(A,M)\simeq r_*r^*M.
\]
\end{proof}


\subsection{The finite group input}\label{sec:finite}

We explain exactly how the finite group theorem enters the proof.
Recall that Gepner--Heller's Adams isomorphism theorem has a genuine fixed point target.
We first compare that target with lisse direct image on compact objects.

\subsubsection*{Free and cofree objects}
Suppose that presentable stable categories have adjunctions
\[
 i_!:\C\rightleftarrows\D:i^*,\qquad i^*:\D\rightleftarrows\C:i_*,
\]
with \(i_!\) and \(i_*\) fully faithful.
Assume \(\D\) is presentably symmetric monoidal and
\[
 i_!i^*(-)\simeq E\otimes-.
\]
The essential image of \(i_!\) is then a tensor ideal.
We call its objects free.
We call objects in the image of \(i_*\) cofree.
These terms describe two different embeddings of \(\C\).

\begin{lemma}\label{lem:compact-free-cofree}
Suppose \(E = i_! \one\) is a filtered colimit of free objects \(E_\alpha\)
that are dualizable in \(\D\).
Then for every compact \(N\in\C\), the canonical comparison
\[
 i_!N\longrightarrow i_*N
\]
is an equivalence.
\end{lemma}
\begin{proof}
Set \(M=i_!N\).
The functor \(i^*\) preserves colimits, since it is also a left adjoint.
Thus \(M\) is compact in \(\D\).
We have
\[
 M\simeq E\otimes M\simeq\colim_\alpha(E_\alpha\otimes M).
\]
Compactness implies that the identity of \(M\) factors through one term, up to homotopy.
Hence \(M\) is a retract of \(E_\alpha\otimes M\) for some \(\alpha\).

Let \(\mathcal K=\ker(i^*)\).
The essential image of \(i_*\) is the right orthogonal \(\mathcal K^\perp\).
To see this description directly, the unit \(Z\to i_*i^*Z\) has cofiber in \(\mathcal K\).
Every object in the image of \(i_*\) is right orthogonal to \(\mathcal K\), by adjunction.
Conversely, if \(Z\) is in \(\mathcal K^\perp\), its cofiber belongs to both \(\mathcal K\) and \(\mathcal K^\perp\), so is zero. 
This shows that \(Z \simeq i_*i^*Z\) is in the image of \(i_*\). 

The dual \(DE_\alpha\) is free.
Indeed, a duality triangle makes it a retract of
\(DE_\alpha\otimes E_\alpha\otimes DE_\alpha\), which is free.
For \(Z\in\mathcal K\), we have \(E\otimes Z=0\) by the projection formula, and therefore
\[
 DE_\alpha\otimes Z
 \simeq DE_\alpha\otimes E\otimes Z=0.
\]
It follows that
\[
 \Map(Z,E_\alpha\otimes M)
 \simeq\Map(DE_\alpha\otimes Z,M)=0.
\]
Thus the object \(E_\alpha\otimes M\) is cofree.
Its retract \(M\) is cofree as well.
The equivalence \(M\to i_*i^*M=i_*N\) is precisely the canonical comparison.
\end{proof}

In particular, this argument does not say that \(i_!\) preserves duals.
The free embedding need not preserve the tensor unit.
The hypothesis concerns the dualizable objects \(E_\alpha\), together with compactness of \(N\).

\subsubsection*{Constant finite groups}
Let \(H\) be a finite constant group.
Its order is invertible over our characteristic zero base.
Gepner and Heller work with equivariantly quasiprojective schemes over a common qcqs base.
We first remain in that setting.

Their free category identifies with \(\SH(BH_S)\).
At the level of charts, a scheme \(V\to BH_S\) is an \(H\)-torsor \(P\to V\).
The total space \(P\) is a free \(H\)-scheme.
Conversely, a free equivariant map \(P \to S\) gives such a torsor after taking its quotient.
These constructions identify Nisnevich descent and affine bundle invariance on the free site.
The comparison of stabilizations is \cite[Theorem~5.2]{GH}.
Thus the quotient functor \(\pi_!\) in that paper is \(h_\sharp\), where \(h:BH_S\to S\).
Adjunction gives
\[
 h^*=i^*\pi^*,\qquad h_*=\pi_*i_*.
\]

By \cite[Proposition~3.11]{GH}, the free colocalization is tensoring with the universal free spectrum \(E\).
By \cite[Corollary~6.29]{GH}, \(E\) is a sequential colimit of genuinely dualizable free spectra.
The same assertions hold for the family
\(\{J:J\cap N=1\}\) attached to a normal subgroup \(N\) of a finite constant group.
Lemma~\ref{lem:compact-free-cofree} therefore applies.

The Adams map of \cite[Definition~6.4 and Theorem~6.33]{GH} is an equivalence
\[
 \tau:\pi_!\xrightarrow{\ \sim\ }\pi_*i_!.
\]
For compact \(N\in\SH(BH_S)\), compose it with the canonical free-to-cofree map:
\begin{equation}\label{eq:finite-ordinary}
 h_\sharp N\xrightarrow{\ \sim\ }\pi_*i_!N
 \xrightarrow{\ \sim\ }\pi_*i_*N=h_*N.
\end{equation}
This composite is the diagonal norm.
Here is the map comparison.
Take its \(h^*\dashv h_*\) adjoint.
The comparison \(i_!N\to i_*N\) corresponds to \(\mathrm{id}_N\) under \(i^*\dashv i_*\).
The remaining adjoint is the restriction of the diagonal transformation in \cite[Definition~6.4]{GH}.
On the free site its finite index ambidexterity is finite étale ambidexterity.
Its unit and counit are therefore the unit and counit for the diagonal of \(BH_S\).
These are exactly the maps in the norm construction in Section~\ref{sec:norm}.

\subsubsection*{Finite etale groups}
Now let \(H/k\) be finite étale, without assuming it constant.
Choose a finite Galois extension \(L/k\) splitting \(H\).
Put
\[
 H_0=H(L),\qquad \Gamma=\Gal(L/k),\qquad K=H_0\rtimes\Gamma.
\]
The constant group \(K\) acts on \(S_L\) through \(\Gamma\).
The normal subgroup \(H_0\) acts trivially on the base.
There are equivalences
\begin{equation}\label{eq:galois-presentations}
 [S_L/K]\simeq BH_S,\qquad [S_L/\Gamma]\simeq S.
\end{equation}
The action of \(\Gamma\) on \(H_0\) is the descent datum defining \(H\).
The splitting \(\Gamma\to K\) supplies the neutral section in the first equivalence.

These constructions preserve equivariant morphisms, Nisnevich squares, and affine-bundle torsors.
For stabilization, \cite[Theorem~5.2]{GH} requires only \(H_0\)-trivial spheres on the source.
Such a bundle on \(S_L\) is a \(\Gamma\)-equivariant bundle.
It descends to a vector bundle on \(S\).
Its Thom object is invertible in ordinary motivic stabilization.
On the target, \cite[Remark~5.3]{GH} identifies genuine \(\Gamma\)-spectra on the free base \(S_L\) with \(\SH(S)\).
Consequently the normal-subgroup form of the finite constant-group theorem transports to \eqref{eq:finite-ordinary} for \(H\).
It transports the adjunctions and hence the norm maps as well.

\begin{proposition}\label{prop:finite-adams}
Let \(H/k\) be finite étale and \(h:BH_S\to S\).
For compact \(N\in\SH(BH_S)\), the diagonal norm
\[
 \nu_h:h_\sharp N\longrightarrow h_*N
\]
is an equivalence.
It extends to an equivalence \(h_\sharp\simeq h_*^{\cont}\).
The ordinary norm is compatible with smooth base change.
\end{proposition}
\begin{proof}
The preceding argument proves the compact assertion on the admissible quasiprojective bases.
For a general \(S\), take a finite affine open cover.
Restriction to these opens preserves compact objects by Lemma~\ref{lem:compact-pullback}.
Smooth base change identifies the restricted ordinary norm with the norm just proved.
Nisnevich descent proves the assertion over \(S\).
This also explains the asserted base change compatibility: all maps are built from the same diagonal, projection, and adjunction maps.
Finally, both \(h_\sharp\) and \(h_*^{\cont}\) preserve colimits.
The equivalence on compact objects extends by Lemma~\ref{lem:continuous}.
\end{proof}



\section{The diagonal norm and the Adams isomorphism}\label{sec:norm}

\subsection{Construction of the map}
Let \(f:B=BG_S\to S\).
Let \(q_1,q_2:B\times_S B\to B\) be the first and second projections.
Recall \(\mu_f=q_{1*}\Delta_\sharp\one_B\).
The adjoint of the ordinary norm
\[
 \nu_f:f_\sharp(M\otimes\mu_f)\longrightarrow f_*M
\]
is the following map:
\begin{equation}\label{eq:norm-adjoint}
\begin{split}
 f^*f_\sharp(M\otimes\mu_f)
 &\simeq q_{2\sharp}q_1^*(M\otimes\mu_f)\\
 &\longrightarrow q_{2\sharp}
       (q_1^*M\otimes\Delta_\sharp\one_B)\\
 &\simeq q_{2\sharp}\Delta_\sharp M
 \simeq M.
\end{split}
\end{equation}
The arrow is induced by the counit
\(q_1^*q_{1*}\Delta_\sharp\one\to\Delta_\sharp\one\).
The first equivalence is smooth base change.
The next equivalence is the smooth projection formula.
This constructs a natural transformation on all coefficients.
Lemma~\ref{lem:continuous} gives its unique lift
\[
 \Nm_f:f_\sharp(-\otimes\mu_f)\longrightarrow f_*^{\cont}.
\]

For a finite étale group \(H\), the diagonal of \(BH_S\) is finite étale.
Finite étale ambidexterity gives
\[
 \mu_h=q_{1*}\Delta_\sharp\one
 \simeq q_{1*}\Delta_*\one
 =\one.
\]
Thus Proposition~\ref{prop:finite-adams} applies to exactly this construction.

The norm follows the construction in
\cite[Notation~1.3.1 and Construction~1.3.5]{RaksitTate},
with the adjustment between smooth and exceptional pushforwards.

\subsection{The graph calculation}
Choose \(H\subset G\) by Theorem~\ref{thm:finite-subgroup}, and write
\[
 C=BH_S,\quad B=BG_S,\quad r:C\to B,\quad h=fr.
\]
Form
\begin{equation}\label{eq:graph-diagram}
 E=C\times_S B,\qquad
 p:E\to C,\quad q:E\to B,\quad
 l=(\mathrm{id},r):C\to E.
\end{equation}
Thus \(pl=\mathrm{id}\) and \(ql=r\).
Set
\[
 L=l_\sharp\one_C,\qquad K=p_*L.
\]
The graph is a pullback of \(\Delta\).
Smooth base change identifies
\begin{equation}\label{eq:K-mu}
 K\simeq r^*\mu_f.
\end{equation}
By Corollary~\ref{cor:graph-dual}, \(L\) is dualizable.
Lemma~\ref{lem:section} provides a dualizable \(K\) and an equivalence
\[
 b:L\xrightarrow{\ \sim\ }l_*K,\qquad p_*b=\mathrm{id}_K.
\]
If \(\epsilon:p^*K\to L\) is the counit, then \(b\epsilon=\eta_l\).

For \(N\in\SH(C)\), put
\[
 L_N=l_\sharp N=L\otimes p^*N,\qquad
 \epsilon_N=\epsilon\otimes p^*N.
\]
The same lemma gives
\begin{equation}\label{eq:bN-unit}
 b_N:L_N\xrightarrow{\ \sim\ }l_*(K\otimes N),
 \qquad b_N\epsilon_N=\eta_l.
\end{equation}
The unit on the right has source \(p^*(K\otimes N)\).
The formula holds for every \(N\); there is no dualizability assumption on \(N\).

\begin{lemma}\label{lem:graph-finite-norm}
For compact \(N\in\SH(C)\), the map
\begin{equation}\label{eq:wN}
\begin{split}
 w_N:r_\sharp N=q_\sharp L_N
 &\xrightarrow{\nu_q(L_N)}q_*L_N\\
 &\xrightarrow{q_*b_N}q_*l_*(K\otimes N)
 =r_*(K\otimes N)
\end{split}
\end{equation}
is an equivalence.
\end{lemma}
\begin{proof}
The map \(q\) is the base change of \(h:BH_S\to S\).
Its twist is the unit.
We can test \(\nu_q(L_N)\) on a smooth scheme atlas \(U\to B\) with Nisnevich local sections.
Let \(P\to U\) be its \(G\)-torsor.
Its pullback along \(r\) has source \(P/H\).
The maps from \(P/H\) to \(BH_U\) and \(BH_S\) are smooth and representable. Denote them by \(t\) and \(z\).
The map \(t\) is the base change of \(l\).
The map \(z\) is the base change of the atlas \(U\to B\).
The coefficient on this chart is
\[
 (L_N)|_{BH_U}=t_\sharp z^*N.
\]
It is compact by Lemma~\ref{lem:compact-pullback}.
Proposition~\ref{prop:finite-adams} therefore proves that the norm is an equivalence on this chart.
Smooth base change and conservativity of the atlas prove it globally.
The second arrow in \eqref{eq:wN} is an equivalence by \eqref{eq:bN-unit}.
\end{proof}

\subsection{Compatibility of norms}
We now verify the specific natural transformation.
By smooth projection and \eqref{eq:K-mu},
\[
 f_\sharp(r_\sharp N\otimes\mu_f)
 \simeq h_\sharp(K\otimes N).
\]
Under this identification, the \(f^*\dashv f_*\) adjoint of \(\nu_f(r_\sharp N)\) is
\begin{equation}\label{eq:induced-norm-adjoint}
 q_\sharp p^*(K\otimes N)
 \xrightarrow{q_\sharp\epsilon_N}
 q_\sharp L_N=r_\sharp N.
\end{equation}
To obtain this formula, pull the diagonal square in \eqref{eq:norm-adjoint} back along \(r\) in its first factor.
More precisely, use \(a=r\times\mathrm{id}:C\times_S B\to B\times_S B\).
Smooth base change identifies \(q_1^*r_\sharp N\) with \(a_\sharp p^*N\).
The diagonal becomes \(l\).
The counit for \(q_1\) becomes the counit for \(p\).
The projection formula moves the coefficient \(N\) to \(p^*N\).
The final composite is \(ql=r\).
These four identifications give \eqref{eq:induced-norm-adjoint}.

We postcompose \eqref{eq:induced-norm-adjoint} with \(w_N\).
Put \(Y=K\otimes N\). Recall that
\[
 \epsilon_N:
 p^*Y\longrightarrow L_N
\]
is the tensor product of the counit
\(\epsilon:p^*K\to L\) with the identity of \(p^*N\).

First, the normalization in \eqref{eq:bN-unit} gives the
commutative triangle
\[
\begin{tikzcd}[column sep=large,row sep=large]
 p^*Y
   \arrow[r,"\epsilon_N"]
   \arrow[dr,"\eta_l"']
 &
 L_N
   \arrow[d,"b_N","\sim"']
 \\
 & l_*Y .
\end{tikzcd}
\]

Next, \(q\) is the base change of \(h:BH_S\to S\).
Since \(H\) is finite étale, the twist of \(h\), and hence
that of \(q\), is canonically the unit.
The ordinary norm therefore has the form
\[
 \nu_q:q_\sharp\longrightarrow q_*.
\]
Naturality of this transformation at \(\epsilon_N\)
gives the commutative square
\[
\begin{tikzcd}[column sep=huge,row sep=large]
 q_\sharp p^*Y
   \arrow[r,"q_\sharp\epsilon_N"]
   \arrow[d,"\nu_q(p^*Y)"']
 &
 q_\sharp L_N
   \arrow[d,"\nu_q(L_N)"]
 \\
 q_*p^*Y
   \arrow[r,"q_*\epsilon_N"']
 &
 q_*L_N .
\end{tikzcd}
\]

By \eqref{eq:wN}, the map \(w_N\) is the composite
\[
\begin{aligned}
 r_\sharp N=q_\sharp L_N
 &\xrightarrow{\nu_q(L_N)}q_*L_N\\
 &\xrightarrow{q_*b_N}q_*l_*Y=r_*Y.
\end{aligned}
\]
This composite is defined for every \(N\).
Compactness is needed only for the assertion in
Lemma~\ref{lem:graph-finite-norm} that it is an equivalence.

Postcomposing the preceding naturality square with
\(q_*b_N\), and using \(b_N\epsilon_N=\eta_l\), gives
\[
\begin{tikzcd}[column sep=huge,row sep=large]
 q_\sharp p^*Y
   \arrow[r,"q_\sharp\epsilon_N"]
   \arrow[d,"\nu_q(p^*Y)"']
 &
 r_\sharp N
   \arrow[d,"w_N"]
 \\
 q_*p^*Y
   \arrow[r,"q_*\eta_l"']
 &
 r_*Y .
\end{tikzcd}
\]
Equivalently, we obtain
\begin{align}
 w_N\circ q_\sharp\epsilon_N
 &=q_*b_N\circ\nu_q(L_N)\circ q_\sharp\epsilon_N
   \notag\\
 &=q_*(b_N\circ\epsilon_N)\circ
       \nu_q\bigl(p^*(K\otimes N)\bigr)
   \notag\\
 &=q_*\eta_l\circ
       \nu_q\bigl(p^*(K\otimes N)\bigr).
 \label{eq:norm-pasting}
\end{align}
Since \(h_*=f_*r_*\), the
\(f^*\dashv f_*\) adjoint of \(\nu_h(Y)\) is
\[
 f^*h_\sharp Y
 \xrightarrow{f^*\nu_h(Y)}f^*h_*Y
 \xrightarrow{\epsilon^f_{r_*Y}}r_*Y,
\]
where \(\epsilon^f\) is the counit.
Write the canonical smooth base-change equivalences as
\[
 \alpha_Y:f^*h_\sharp Y\xrightarrow{\sim}q_\sharp p^*Y,
 \qquad
 \beta_Y:f^*h_*Y\xrightarrow{\sim}q_*p^*Y.
\]
The diagonal construction of the norm commutes with base change, so
\[
 \beta_Y\circ f^*\nu_h(Y)
 =\nu_q(p^*Y)\circ\alpha_Y.
\]

Consequently,
\[
 \epsilon^f_{r_*Y}\circ f^*\nu_h(Y)
 =
 q_*\eta_l\circ\nu_q(p^*Y)\circ\alpha_Y.
\]
After identifying the source by \(\alpha_Y\), this is exactly
the last expression in \eqref{eq:norm-pasting}.
Note that \(w_N\circ q_\sharp\epsilon_N\) is the the
\(f^*\dashv f_*\) adjoint of \(f_*(w_N)\circ\nu_f(r_\sharp N)\). 
So by passing to the adjoint of \eqref{eq:norm-pasting}, we obtain
\begin{equation}\label{eq:norm-transitivity}
 f_*(w_N)\circ\nu_f(r_\sharp N)=\nu_h(K\otimes N).
\end{equation}
This says that 
\begin{quote}
\itshape
The norm for \(G\), evaluated on an object induced from \(H\), agrees with the norm for \(H\)
up to the comparison \(w_N\)
\end{quote}

\begin{theorem}[Lisse Adams isomorphism]\label{thm:lisse-adams}
The norm
\[
 \Nm_f:f_\sharp(-\otimes\mu_f)\xrightarrow{\ \sim\ }f_*^{\cont}
\]
is an equivalence.
\end{theorem}
\begin{proof}
If \(N\in\SH(C)\) is compact, so is \(K\otimes N\).
Indeed, \(K\) is dualizable, and tensoring with its dual preserves filtered colimits.
Proposition~\ref{prop:finite-adams} makes the right side of \eqref{eq:norm-transitivity} an equivalence.
Lemma~\ref{lem:graph-finite-norm} makes \(f_*(w_N)\) an equivalence.
Hence \(\nu_f(r_\sharp N)\) is an equivalence.

Choose compact generators \(N_i\) of \(\SH(C)\).
Proposition~\ref{prop:nis-reduction} and Lemma~\ref{lem:compact-pullback}
show that \(r_\sharp N_i\) are compact generators of \(\SH(B)\).
On them the ordinary and continuous direct images agree.
Both functors in the theorem preserve colimits.
The equivalence on these generators therefore proves the theorem.
\end{proof}

\begin{corollary}\label{cor:compact-adams}
If \(M\) is compact, the ordinary norm
\[
 f_\sharp(M\otimes\mu_f)\xrightarrow{\ \sim\ }f_*M
\]
is an equivalence.
In particular this holds for \(M=j_\sharp\one_U\), where \(j:U\to BG_S\) is a smooth scheme chart.
\end{corollary}
\begin{proof}
The comparison \(f_*^{\cont}\to f_*\) is an equivalence on compact objects.
The last assertion follows from compactness of the scheme unit and \(j_\sharp\dashv j^*\).
\end{proof}

Define \(T_f\) as the cofiber of the ordinary norm.
Under the comparison above, this is the Tate pushforward of
\cite[Construction~1.3.8]{RaksitTate}.
Theorem~\ref{thm:lisse-adams} gives
\[
 T_f(M)\simeq
 \cofib\bigl(f_*^{\cont}M\longrightarrow f_*M\bigr).
\]
It vanishes on compact objects. For the genuine Tate construction for compact Lie groups, see
\cite{GreenleesMayTate}.

\section{Comparison with genuine fixed points}
\label{sec:genuine-comparison}

In this section, \(S\) is smooth and quasiprojective over \(k\).
The \(G\)-action on \(S\) is trivial.
We write
\[
 f:BG_S\longrightarrow S,\qquad
 \mathcal C_S=\SH(BG_S),\qquad
 \mathcal D_S=\SH^G(S).
\]
The smoothness assumption on \(S\) will enter the proof of
Lemma~\ref{lem:free-pullback}.
Note that the lisse Adams theorem itself does not need this assumption.

\subsection{The free and cofree embeddings}

We first describe the relation between the two categories.
This also specifies the meaning of \emph{free} in this section.
Let \(\mathcal A_S=\mathrm{Sm}^G_S\) be the admissible smooth equivariant site.
Let \(\mathcal A_S^{\mathrm{fr}}\) be its full subcategory of schemes with
free \(G\)-action.

A free action has a quotient algebraic space \(P/G\), and
\(P\to P/G\) is a \(G\)-torsor
\cite[Tag 06PH]{Stacks}.
If \(P\) is smooth over \(S\), then \(P/G\) is smooth over \(S\).
A smooth scheme chart \(X\to BG_S\) is therefore the same datum as
a \(G\)-torsor \(P\to X\) whose total space is smooth over \(S\).
Morphisms of charts correspond to equivariant maps of their total spaces.
One may allow algebraic space charts first and then use a Nisnevich cover
by schemes. See \cite[Example 12.23(iii)]{KhanRaviStacks}.
Conversely, smooth charts \(X\to BG_S\) with \(X\) affine form a Nisnevich basis. 
Their associated torsors \(P\to X\) have smooth equivariantly quasiprojective total spaces over \(S\), 
so they belong to the genuine category.

Under this identification, equivariant Nisnevich squares of free schemes
are Nisnevich squares on the quotient.
Equivariant vector bundles and their torsors descend to vector bundles
and their torsors on the quotient.
Thus the motivic category of the free site, stabilized by the ordinary
Tate sphere, presents \(\mathcal C_S\).
This agrees with the lisse extension convention by Nisnevich descent.
See \cite[Section 1, especially Lemma 1.5]{KhanRavi}.
An equivariant representation sphere restricts to the Thom object of
an associated vector bundle on each quotient chart.
These Thom objects are already invertible in the ordinary stable
motivic category, and their inverses are compatible with pullback.
Hence stabilization by all equivariant spheres does not change this
free category.

\begin{lemma}\label{lem:free-embeddings}
Restriction to the free site gives adjunctions
\[
 i_!:\mathcal C_S\rightleftarrows\mathcal D_S:i^*,
 \qquad
 i^*:\mathcal D_S\rightleftarrows\mathcal C_S:i_*.
\]
Both \(i_!\) and \(i_*\) are fully faithful.
There is an object \(E_S\), equipped with a map \(E_S\to\one\), such that
\[
 i_!i^*M\simeq E_S\otimes M.
\]
The image of \(i_!\) is a localizing tensor ideal.
The canonical map
\[
 \kappa_C:i_!C\longrightarrow i_*C
\]
is the map whose restriction under \(i^*\) is \(\mathrm{id}_C\).
\end{lemma}

\begin{proof}
The free subcategory is a sieve: if an equivariant map \(V\to U\)
has free target, then \(V\) is free.
Let \(F_S\) be its characteristic presheaf:
\[
 F_S(V)=
 \begin{cases}
 *&\text{if \(V\) is free},\\
 \varnothing&\text{otherwise}.
 \end{cases}
\]
At the presheaf level, left Kan extension from the sieve is extension
by the empty space.
It is fully faithful, and its composite with restriction is
\(F_S\times-\).

We check that restriction respects the relations defining motivic spaces.
Write \(r\) for restriction to the free site.
If \(Y\to X\) is an equivariant affine bundle and \(V\to X\) is a
free test scheme, then
\[
 Y\times_XV\longrightarrow V
\]
is an affine bundle between free schemes.
Its quotient is an ordinary affine bundle, so it becomes a
motivic equivalence on the free site.

Presheaf density assembles these maps into the restricted map
\(rh_Y\to rh_X\). More precisely,
\[
 rh_X\simeq\operatorname*{colim}_{V\to X,\;V\ {\rm free}}h_V,
 \qquad
 rh_Y\simeq\operatorname*{colim}_{V\to X,\;V\ {\rm free}}
 h_{Y\times_XV}.
\]
Since motivic localization preserves colimits, \(rh_Y\to rh_X\)
also becomes an equivalence.
The same argument applies to Nisnevich relations, because their
defining squares remain distinguished after pullback to free schemes.


Left extension and restriction commute with tensoring by representation
spheres and their restrictions.
Thus the adjunction and its identity unit pass through pointing and
inversion of these spheres.
This gives the first adjunction and its full faithfulness.
Since restriction preserves colimits, it also has a right adjoint.
This right adjoint is fully faithful. In fact, for \(A,B\in\mathcal C_S\),
the two adjunctions and full faithfulness of \(i_!\) give
\[
 \Map(A,i^*i_*B)
 \simeq\Map(i_!A,i_*B)
 \simeq\Map(i^*i_!A,B)
 \simeq\Map(A,B).
\]
Thus the counit \(i^*i_*B\to B\) is an equivalence.

More explicitly, put
\[
 E_S=\Sigma^\infty_G(L_{\mathrm{mot}}F_S)_+.
\]
Motivic localization preserves finite products
\cite[Proposition 3.15]{Hoyois}.
The identity \(F_S\times F_S=F_S\), followed by stabilization, gives
\(E_S\otimes E_S\simeq E_S\).
The presheaf formula for left Kan extension now gives
\(i_!i^*=E_S\otimes-\).
This proves the tensor ideal assertion.
Finally, the unit \(i_!C\to i_*i^*i_!C\) gives \(\kappa_C\).
Full faithfulness identifies its restriction with the identity of \(C\).
\end{proof}

\begin{remark}
Note that the presheaf \(F_S\) need not be invariant under equivariant affine
bundles. We use its motivic localization in the definition of \(E_S\).
In particular, the proof does not assume that an equivariant affine
bundle preserves strict freeness in both directions.
This distinction matters for groups with unipotent subgroups.
\end{remark}

\subsection{Dualizable approximations to the free object}

Let \(p:S\to\operatorname{Spec}k\) be the structural map.

\begin{lemma}\label{lem:free-pullback}
There is a canonical equivalence
\[
 p^*E_k\simeq E_S.
\]
Moreover, \(E_S\) is a filtered colimit of objects which are both free
and dualizable in \(\mathcal D_S\).
\end{lemma}

\begin{proof}
We first prove the corresponding statement for presheaves.
By the presheaf density formula,
\[
 F_k\simeq
 \colim_{U\in\mathcal A_k^{\mathrm{fr}}}h_U.
\]
Pullback sends \(h_U\) to \(h_{U\times_k S}\).
For \(V\in\mathcal A_S\), evaluation therefore gives
\[
 (p^*F_k)(V)\simeq
 \colim_{U\in\mathcal A_k^{\mathrm{fr}}}
 \operatorname{Hom}_{G,k}(V,U).
\]

If \(V\) is free, it belongs to \(\mathcal A_k^{\mathrm{fr}}\).
Here we use that \(S/k\) is smooth and quasiprojective.
The pair \((V,\mathrm{id}_V)\) is initial in the category \(\mathcal A_k^{\mathrm{fr}}\),
so the colimit is contractible.
If \(V\) is not free, that category is empty.
Indeed, an equivariant map to a free scheme forces every stabilizer
functor in the source to be trivial.
We have proved \(p^*F_k\simeq F_S\).
Pullback preserves the relations used in motivic localization.
It also commutes with suspension spectra.
This gives \(p^*E_k\simeq E_S\).

The same density formula writes \(E_k\) as a colimit of spectra
\(\Sigma^\infty_G U_+\), with \(U\) free, smooth, and equivariantly
quasiprojective over \(k\).
Each of these spectra is dualizable by Proposition~\ref{prop:geometric-duality}.
It is also free.

We explain how to express the given colimit as a filtered colimit
of free dualizable objects. Write the original diagram as
\[
 D:I\longrightarrow\mathcal D_S,
\]
where every \(D(i)\) is free and dualizable.

Let $\operatorname{Fin}(I)$ be the poset of finite simplicial
subsets of the underlying simplicial set of $I$.
Here finite means having finitely many nondegenerate simplices.
This poset is filtered, since finite unions remain finite.

Following the proof of
\cite[Proposition~9.2.1.1, \href{https://kerodon.net/tag/06DB}
{Tag~06DB}]{Kerodon}, which uses
\cite[Corollary~9.1.6.6, \href{https://kerodon.net/tag/03DJ}
{Tag~03DJ}]{Kerodon}, we obtain
\[
  \operatorname*{colim}_{i\in I} D(i)
  \simeq
  \operatorname*{colim}_{K\in\operatorname{Fin}(I)}
  \left(\operatorname*{colim}_{K} D|_K\right).
\]
Each inner colimit is a finite colimit of dualizable objects,
and hence is dualizable.
Each inner colimit is also free, since the free subcategory is
closed under finite colimits. Thus the original colimit is a
filtered colimit of free dualizable objects.

Finally, apply \(p^*\) to this presentation.
Pullback preserves colimits and dualizable objects.
It takes a free chart to a free chart.
Thus we obtain
\[
 E_S\simeq\colim_a E_a,
 \qquad E_a\text{ free and dualizable in }\mathcal D_S.
 \qedhere
\]
\end{proof}

\subsection{Compact objects and fixed points}

\begin{corollary}\label{cor:genuine-compact-free-cofree}
For every compact \(C\in\mathcal C_S\), the canonical comparison
\[
 \kappa_C:i_!C\longrightarrow i_*C
\]
is an equivalence.
\end{corollary}
\begin{proof}
Lemma~\ref{lem:free-embeddings} identifies the free colocalization with tensoring by \(E_S\).
Lemma~\ref{lem:free-pullback} supplies its filtered presentation by genuinely dualizable free objects.
Apply Lemma~\ref{lem:compact-free-cofree}.
\end{proof}

Let
\[
 \pi^*:\SH(S)\longrightarrow\SH^G(S)
\]
be inflation, with right adjoint \(\pi_*\).
On the free site, a map from a torsor \(P\) to a scheme with trivial
action is the same as a map from \(P/G\).
Thus \(f^*=i^*\pi^*\), and adjunction gives \(f_*=\pi_*i_*\).

\begin{lemma}\label{lem:genuine-fixed-continuous}
The genuine fixed point functor \(\pi_*\) preserves colimits.
\end{lemma}

\begin{proof}
The category \(\SH(S)\) is generated by compact objects
\(\Sigma^{a,b}\Sigma^\infty U_+\), for smooth quasiprojective
\(U\to S\).
Inflation sends this object to the corresponding genuine suspension
spectrum with trivial \(G\)-action, with the same Tate suspension.
It is compact by \cite[Proposition 6.4(3)]{Hoyois}.
Therefore \(\pi^*\) preserves compact objects.
Its right adjoint preserves filtered colimits.
It is exact, so it also preserves finite colimits.
Writing arbitrary coproducts as filtered colimits of finite coproducts
shows that it preserves coproducts.
An exact functor between stable categories which preserves coproducts
preserves all colimits.
\end{proof}

\begin{theorem}\label{thm:genuine-continuous-comparison}
There is a canonical equivalence
\[
 \theta:f_*^{\cont}\xrightarrow{\ \sim\ }\pi_*i_!.
\]
It commutes with the canonical maps from both functors to \(f_*\).
\end{theorem}

\begin{proof}
Both \(f_*^{\cont}\) and \(\pi_*i_!\) preserve colimits.
By Corollary~\ref{cor:genuine-compact-free-cofree}, the comparison
\[
 \pi_*i_!C\longrightarrow\pi_*i_*C=f_*C
\]
is an equivalence for every compact \(C\in\mathcal C_S\).
The same is true for \(f_*^{\cont} C\to f_*C\), by the definition of
the continuous approximation.
Both functors are therefore the left Kan extension of
\(f_*|_{\mathcal C_S^\omega}\).
The uniqueness of left Kan extension gives \(\theta\), including its
compatibility with the maps to \(f_*\).
\end{proof}

\begin{corollary}[Genuine Adams isomorphism]
\label{cor:genuine-adams-repaired}
The diagonal lisse norm has a canonical continuous lift
\[
 \mathrm{Nm}^{\mathrm{gen}}_f:
 f_\sharp(-\otimes\mu_f)\longrightarrow\pi_*i_!.
\]
This lift is an equivalence.
Its composite with \(\pi_*i_!\to f_*\) is the diagonal lisse norm.
\end{corollary}

\begin{proof}
On compact objects define the lift by composing the diagonal norm
with the inverse of the comparison in
Corollary~\ref{cor:genuine-compact-free-cofree}, after applying \(\pi_*\).
Both the source and the target of the desired lift preserve colimits.
Left Kan extension therefore gives a unique lift on all objects.
Theorem~\ref{thm:genuine-continuous-comparison} identifies it with
\(\theta\) composed with the equivalence of Theorem~\ref{thm:lisse-adams}.
It is consequently an equivalence.
The compatibility with the ordinary norm holds on compact objects
by construction and therefore holds everywhere by left Kan extension.
\end{proof}

The continuous/ordinary distinction has a topological counterpart in
the difference between genuine Adams and the norm to homotopy
fixed points; see
\cite[Theorem~5.2.5 and Remark~5.2.6]{RognesDualizable}.
A related formal approach uses the compactness locus of a geometric
functor \cite[Theorems~1.4 and~2.9, Section~3]{SandersAdams}.
We do not apply its rigid-compactness hypotheses to the lisse category.

\begin{corollary}[Free torsors for split reductive groups]
\label{cor:explicit-free-reductive-splitting}
Let \(G\) be split connected reductive of dimension \(n\)
and rank \(r\), and let \(S\) be smooth and quasiprojective
over \(k\), with trivial \(G\)-action.
Let \(p:Z\to S\) be smooth and quasiprojective, and let
\(P\to Z\) be a right \(G\)-torsor.
We regard \(P\) as a left \(G\)-scheme by
\(g\cdot p_0=p_0g^{-1}\).
Suppose that \(P\) has a reduction \(Q\) to a split Borel
subgroup \(B_0\subseteq G\), and put
\[
 E=Q\times^{B_0}(\mathfrak g/\mathfrak b_0).
\]
Then
\[
 (\Sigma_G^\infty P_+)^G
 \simeq S^{r,r}\otimes p_\sharp\Th_Z(E).
\]
For the trivial torsor this becomes
\[
 \bigl(\Sigma_G^\infty(Z\times_S G_S)_+\bigr)^G
 \simeq S^{n,(n+r)/2}\otimes\Sigma^\infty Z_+.
\]
\end{corollary}

\begin{proof}
We will apply the genuine Adams isomorphism to the free
suspension spectrum of \(P\), then compute the resulting
twist using the specified Borel reduction.

Let \(u:Z\to BG_S\) classify \(P\), and let
\(f:BG_S\to S\) be the structure map, so that \(fu=p\).
The free-site comparison of
Section~\ref{sec:genuine-comparison} identifies
\[
 \Sigma_G^\infty P_+\simeq i_!u_\sharp\one_Z.
\]
Here \(u\) is smooth: its base change along the standard
atlas \(S\to BG_S\) is the smooth morphism \(P\to S\).
Applying Corollary~\ref{cor:genuine-adams-repaired} and
the smooth projection formula gives
\[
 \begin{aligned}
 (\Sigma_G^\infty P_+)^G
 &\simeq f_\sharp(u_\sharp\one_Z\otimes\mu_G)\\
 &\simeq f_\sharp u_\sharp(u^*\mu_G)\\
 &\simeq p_\sharp(u^*\mu_G).
 \end{aligned}
\]
By Corollary~\ref{cor:new-split-reductive}, the chosen
Borel reduction \(Q\) identifies
\[
 u^*\mu_G
 \simeq S^{r,r}\otimes
 \Th_Z\bigl(Q\times^{B_0}(\mathfrak g/\mathfrak b_0)\bigr)
 = S^{r,r}\otimes\Th_Z(E).
\]
The projection formula for \(p_\sharp\) proves the
first assertion.

For the trivial torsor, choose the trivial Borel reduction.
Then \(E\) is a trivial vector bundle of rank \((n-r)/2\).
Consequently,
\[
 S^{r,r}\otimes\Th_Z(E)
 \simeq S^{n,(n+r)/2}_Z.
\]
Applying \(p_\sharp\) and using
\(p_\sharp\one_Z=\Sigma^\infty Z_+\) proves the second
assertion.
\end{proof}

\begin{example}[Classical split reductive groups]
With \(S=Z=\Spec k\),
Corollary~\ref{cor:explicit-free-reductive-splitting} gives
\[
 \begin{aligned}
 (\Sigma_{\mathrm{GL}_m}^\infty(\mathrm{GL}_m)_+)
       ^{\mathrm{GL}_m}
 &\simeq S^{m^2,m(m+1)/2},\\
 (\Sigma_{\mathrm{SL}_m}^\infty(\mathrm{SL}_m)_+)
       ^{\mathrm{SL}_m}
 &\simeq S^{m^2-1,(m^2+m-2)/2},\\
 (\Sigma_{\mathrm{Sp}_{2m}}^\infty(\mathrm{Sp}_{2m})_+)
       ^{\mathrm{Sp}_{2m}}
 &\simeq S^{m(2m+1),m(m+1)}.
 \end{aligned}
\]
Here \(m\geq1\) in the first and third formulas and
\(m\geq2\) in the second, and the groups act freely by
translation. 
\end{example}


\section{Families and geometric fixed points}
\label{sec:isotropy}

Throughout this section, \(k\) has characteristic zero.
The group \(G\) is smooth, affine, and linearly reductive.
The base \(S\) is a smooth quasiprojective \(k\)-scheme with trivial
\(G\)-action.
We write
\[
 \D_G=\SH^G(S).
\]
When an argument needs a scheme affine over \(k\), rather than merely
affine over \(S\), we make two Jouanolou replacements.
First choose an ordinary affine Jouanolou torsor \(S'\to S\).
Then apply equivariant Jouanolou to \(X\times_S S'\) over the affine
base \(S'\).
The resulting scheme is affine over \(k\).
Its map to \(X\) is a composite of equivariant affine bundles, so it
is a motivic equivalence over the original base \(S\).

\subsection{The stabilizer convention}

A stabilizer at a geometric point is a group over its algebraically
closed residue field.
It need not be defined over \(k\).
Moreover, its reductive quotient is not a specified subgroup of \(G\).
We therefore use an embedded Levi subgroup.

In characteristic zero, every affine algebraic group \(K\) has a
decomposition
\[
 K=R_u(K)\rtimes L.
\]
Here \(R_u(K)\) is its unipotent radical and \(L\) is linearly reductive.
Any two choices of \(L\) are conjugate by \(R_u(K)\).
These assertions include disconnected groups. See
\cite[Proposition 5.4.1]{ConradReductive}.
Thus the conjugacy class of \(L\) inside \(G\) is well defined.
We call it the \emph{Levi class} of the stabilizer.


A \emph{geometric family} \(\F\) specifies the allowed Levi classes.
For each algebraically closed extension \(\Omega/k\), it gives a set
\(\F(\Omega)\) of \(G_\Omega\)-conjugacy classes of closed linearly
reductive subgroups of \(G_\Omega\).
We require two compatibility conditions:
\begin{enumerate}
\item
If \([H]\in\F(\Omega)\) and \(K\subset H\) is a closed linearly
reductive subgroup, then \([K]\in\F(\Omega)\).
\item
For every \(k\)-embedding \(\Omega\hookrightarrow\Omega'\) of
algebraically closed fields,
\[
 [H]\in\F(\Omega)
 \quad\Longleftrightarrow\quad
 [H_{\Omega'}]\in\F(\Omega').
\]
Here \(H_{\Omega'}\) denotes extension of scalars.
Thus membership is independent of enlarging the field of coefficients.
\end{enumerate}

Note that (ii) implies Galois invariance.
For \(\sigma\in\operatorname{Gal}(\bar k/k)\), let \(\sigma H\)
be obtained by applying \(\sigma\) to the coefficients defining \(H\).
Then
\[
 [H]\in\F(\bar k)
 \quad\Longleftrightarrow\quad
 [\sigma H]\in\F(\bar k).
\]
The allowed classes may be permuted by Galois groups. They need not
individually have representatives over \(k\).

We write \(X\in\Sm^G_S[\F]\) if, for every geometric point
\(x:\operatorname{Spec}\Omega\to X\), the Levi subgroup
\(L_x\subset G_x\subset G_\Omega\) satisfies
\([L_x]\in\F(\Omega)\).
This is independent of the choice of \(L_x\), by Levi conjugacy.

\begin{lemma}\label{lem:family-motivic}
The presheaf
\[
 E\F(X)=
 \begin{cases}
 *,&X\in\Sm^G_S[\F],\\
 \varnothing,&\text{otherwise}
 \end{cases}
\]
is a motivic space.
\end{lemma}

\begin{proof}
We first recall why a linearly reductive subgroup \(R\subset K\)
is conjugate into any Levi subgroup \(L\subset K\).
Let \(q:K\to K/R_u(K)\) be the quotient map.
Since \(R\cap R_u(K)=1\), both \(R\) and
\((q|_L)^{-1}(q(R))\) are Levi subgroups of \(q^{-1}(q(R))\).
Levi conjugacy therefore gives the assertion.

Now let \(f:X\to Y\) be an equivariant map, and suppose
\(Y\in\Sm^G_S[\F]\).
For a geometric point \(x:\operatorname{Spec}\Omega\to X\),
put \(y=f(x)\).
Equivariance gives \(G_x\subseteq G_y\): every element fixing
\(x\) also fixes \(y\).
Choose Levi subgroups \(L_x\subset G_x\) and \(L_y\subset G_y\).
The preceding argument shows that \(L_x\) is conjugate to a
subgroup of \(L_y\).
Since \([L_y]\in\F(\Omega)\), closure under linearly reductive
subgroups and conjugacy gives \([L_x]\in\F(\Omega)\).
Thus \(X\in\Sm^G_S[\F]\).
In particular, the family condition passes from the target
to the source of an equivariant map.

Let \(p:X\to Y\) be a torsor under a \(G\)-equivariant vector
bundle. If \(Y\) satisfies the family condition, then so does \(X\):
stabilizers can only decrease along an equivariant map, and the
preceding subgroup argument applies.

Conversely, suppose that \(X\) satisfies the family condition.
Choose a geometric point \(y:\operatorname{Spec}\Omega\to Y\)
and a Levi subgroup \(L\subset G_y\).
The fiber \(X_y\) is an affine space with an affine-linear
\(L\)-action. Such an action has a fixed point
\cite[Corollary~A.8]{ReevesHyperspherical}.
We recall the short argument.

Choose an origin in \(X_y\), identifying it with a vector space \(V\).
Write the action as
\[
 \ell\cdot v=\rho(\ell)v+c(\ell).
\]
The identity
\(c(\ell m)=\rho(\ell)c(m)+c(\ell)\)
makes \(W=V\oplus\Omega\) an \(L\)-representation via
\[
 \ell\cdot(v,t)=(\rho(\ell)v+t\,c(\ell),t).
\]
Projection onto the second coordinate gives an exact sequence
\[
 0\longrightarrow V\longrightarrow W
   \longrightarrow\Omega\longrightarrow0,
\]
where the quotient has trivial \(L\)-action.
Linear reductivity gives an \(L\)-equivariant splitting
\cite[5.14 and 5.18]{MilneReductiveGroups}.
The image of \(1\) is therefore an invariant vector \((v_0,1)\).
Thus
\[
 \rho(\ell)v_0+c(\ell)=v_0,
\]
so \(v_0\) determines an \(L\)-fixed point \(x\in X_y(\Omega)\).

We have \(L\subset G_x\).
Choose a Levi subgroup \(M\subset G_x\).
Its class belongs to \(\mathcal F\), since \(X\) satisfies the
family condition.
By the preceding subgroup argument, \(L\) is conjugate into \(M\).
Closure under linearly reductive subgroups and conjugacy gives
\([L]\in\mathcal F\).
Since \(y\) was arbitrary, \(Y\) satisfies the family condition.
Hence \(E\mathcal F(Y)\to E\mathcal F(X)\) is an equivalence,
proving affine bundle invariance.

In an equivariant Nisnevich distinguished square, each geometric
point of the lower-right scheme lies either in the open part or
in the closed complement, over which the other map is an equivariant
isomorphism.
The family condition therefore holds on that scheme exactly when
it holds on the two schemes covering it.
Since \(E\F\) is subterminal, this is precisely Nisnevich excision.
Finally \(E\F(\varnothing)=*\).
\end{proof}

Define
\[
 e_\F=\Sigma^\infty_G(E\F)_+,
 \qquad
 \widetilde e_\F=\cofib(e_\F\to\one).
\]
The basepoint is adjoined in the motivic category.
A localizing tensor ideal is a full stable subcategory closed under
colimits and tensoring with arbitrary objects of \(\D_G\).
Let \(J_\F\subset\D_G\) be the localizing tensor ideal generated by
the suspension spectra of schemes in \(\Sm^G_S[\F]\).
The presheaf density formula gives
\[
 e_\F\otimes\D_G=J_\F,
 \qquad
 e_\F\otimes e_{\F'}\simeq e_{\F\cap\F'}.
\]
In particular, tensoring with \(e_\F\) is the colocalization onto
\(J_\F\).
Tensoring with \(\widetilde e_\F\) gives the complementary
localization.

\subsection{Families and dualizable approximations}

The following approximation will let us apply duality with arbitrary
relative coefficients. 
Write \(e_{\F,S}\) when the base of the family object needs to be shown.

\begin{lemma}\label{lem:new-family-approx}
Suppose \(k\) has characteristic zero and \(S/k\) is smooth and
quasiprojective. Then \(e_{\F,S}\) is a filtered colimit of dualizable
objects belonging to its family ideal.
If \(j:K\hookrightarrow G\) is an inclusion of reductive groups, the
approximations for a family of \(K\) can be chosen so that their
inductions \(j_\sharp C_\alpha\) and \(j_\sharp DC_\alpha\) are
dualizable as well.
\end{lemma}

\begin{proof}
Apply the argument of Lemma~\ref{lem:free-pullback}, replacing
free schemes by schemes satisfying the family condition.
Over \(k\), presheaf density gives a filtered presentation
\[
 e_{\F,k}\simeq\colim_\alpha C_\alpha,
\]
whose terms are finite colimits and retracts of allowed suspension
spectra. They belong to the family ideal and are dualizable by
Proposition~\ref{prop:geometric-duality}.

Let \(p:S\to\Spec k\).
Since the family condition is preserved by equivariant maps,
the same base-change argument gives
\(p^*e_{\F,k}\simeq e_{\F,S}\).
Here smoothness and quasiprojectivity of \(S/k\) ensure that
allowed schemes over \(S\) are also allowed over \(k\).
Pulling back the presentation above proves the first assertion:
\(p^*\) preserves colimits, dualizable objects, and the family
condition.

For the last statement, start with the same constructions for
\(K\).
Over \(k\), each stage and its dual are compact.
Induction preserves compact objects because its right adjoint,
restriction, preserves colimits.
The preceding absolute observation therefore proves dualizability of
both induced objects.
Induction and duality of dualizable objects commute with base change
to \(S\). This proves the claim.
\end{proof}


\subsection{Normal subgroups}

Let \(N\triangleleft G\) be a closed normal subgroup and put
\(Q=G/N\). Both groups are linearly reductive.

For every algebraically closed field extension \(\Omega/k\), write
\(G_\Omega=G\times_k\Omega\) and \(N_\Omega=N\times_k\Omega\).
Define
\[
 \F[N](\Omega)=
 \left\{
 [H]_{G_\Omega}\ \middle|\
 \begin{array}{l}
 H\subset G_\Omega\text{ is a closed linearly reductive subgroup},\\
 N_\Omega\not\subseteq H
 \end{array}
 \right\}.
\]
Here \([H]_{G_\Omega}\) denotes the conjugacy class of the
embedded subgroup \(H\).
Normality of \(N_\Omega\) makes this condition independent
of the representative of the conjugacy class.

These subsets are closed under passage to linearly reductive
subgroups. For every \(k\)-embedding of algebraically closed
fields \(\Omega\hookrightarrow\Omega'\), faithful flatness gives
\[
 [H]\in\F[N](\Omega)
 \quad\Longleftrightarrow\quad
 [H_{\Omega'}]\in\F[N](\Omega').
\]
We denote this geometric family by \(\F[N]\).

\begin{lemma}\label{lem:fixed-scheme-functor}
Fixed schemes induce a colimit-preserving symmetric monoidal functor
\[
 \Phi^N:\D_G\longrightarrow\D_Q
\]
with
\[
 \Phi^N(\Sigma^\infty_G X_+)
   \simeq\Sigma^\infty_Q(X^N)_+,
 \qquad
 \Phi^N(S^V)\simeq S^{V^N}.
\]
\end{lemma}

\begin{proof}
We first record the geometry of fixed schemes.
Since \(X\) is separated, \(X^N\) is a closed subscheme of \(X\).
It is smooth over \(S\) by
\cite[Theorem~4.3.6]{RomagnyFixedPoints}.
Normality of \(N\) gives a \(Q\)-action on \(X^N\).

We construct the functor first on presheaves.
Let \(\mathcal C_G\) and \(\mathcal C_Q\) be the smooth
equivariantly quasiprojective sites over \(S\).
Inflation defines \(I:\mathcal C_Q\to\mathcal C_G\).
Restriction along \(I\) gives
\[
 R:\mathcal P(\mathcal C_G)\longrightarrow
   \mathcal P(\mathcal C_Q),
 \qquad R(F)(T)=F(I(T)).
\]
It preserves colimits and finite products, since these are computed
objectwise.
For a \(Q\)-scheme \(Z\), write \(h_Z\) for its functor of points
on \(\mathcal C_Q\), whether or not \(Z\) belongs to this site.
Then
\[
 R(h_X)(T)
 =\operatorname{Hom}_G(I(T),X)
 =\operatorname{Hom}_Q(T,X^N),
\]
so \(R(h_X)=h_{X^N}\).

We next check affine-bundle invariance.
Let \(Y\to X\) be a \(G\)-equivariant torsor under a vector bundle
\(E\).
Linear reductivity makes
\[
 E_N:=\bigl(E|_{X^N}\bigr)^N
\]
a vector bundle.
On an affine open of \(X^N\), choose a section of the underlying
torsor.
Homogenizing its affine \(N\)-action, as in
Lemma~\ref{lem:family-motivic}, gives an extension of
\(\mathcal O\) by \(E\).
This extension splits \(N\)-equivariantly by
\cite[Lemma~2.17]{Hoyois}.
Thus invariant sections exist locally, and
\(Y^N\to X^N\) is a \(Q\)-equivariant torsor under \(E_N\).

We also check Nisnevich excision.
Fixed schemes preserve fiber products, invariant open immersions,
and the empty scheme.
An equivariant etale morphism remains etale after taking fixed
points.
The isomorphism over the closed complement is preserved as well.
Hence fixed schemes preserve Nisnevich distinguished squares.

These geometric checks imply that \(R\) respects the motivic
relations, even without assuming \(X^N\in\mathcal C_Q\).
Indeed, to prove that a map of presheaves is a motivic equivalence, 
it is enough to prove this after pulling it back to every representable presheaf over its target.
Pull each relation back along \(h_T\to h_{X^N}\), with
\(T\in\mathcal C_Q\).
The resulting affine bundle or distinguished square belongs to
\(\mathcal C_Q\).
Thus these pullback checks show that \(R\) sends the generating
motivic relations to motivic equivalences.

The universal property of motivic localization now gives a
colimit preserving functor on motivic spaces.
It preserves finite products, since both \(R\) and motivic
localization do so
\cite[Proposition~3.15]{Hoyois}.
After pointing, it is therefore symmetric monoidal.

Finally, we pass to spectra.
For a \(G\)-vector bundle \(V\) over \(S\), the functor sends
the cofiber defining its Thom sphere to
\[
 V^N/(V^N-0)=S^{V^N}.
\]
This sphere is invertible in \(\D_Q\).
The universal property of stabilization therefore gives
\[
 \Phi^N:\D_G\longrightarrow\D_Q.
\]
The calculation on representables gives the suspension formula,
and the calculation on Thom spheres gives the second formula.
\end{proof}

Write \(Y^N\) for unstable fixed points.
Presheaf inflation is left adjoint to restriction.
Since inflation preserves affine bundles and Nisnevich squares,
restriction preserves local presheaves.
Thus the adjunction descends to motivic spaces, with counit
\[
 \operatorname{inf}_Q^G(Y^N)\longrightarrow Y.
\]
Suspending and taking the stable adjoint gives
\[
 \Sigma_Q^\infty(Y^N)
 \longrightarrow (\Sigma_G^\infty Y)^N,
\]
where the target uses genuine fixed points, the right adjoint
to genuine inflation.


\begin{theorem}[Localized inflation]\label{thm:normal-inflation}
The functor
\[
 L_N\operatorname{inf}_Q^G:
 \D_Q\xrightarrow{\ \simeq\ }L_N\D_G
\]
is an equivalence of symmetric monoidal categories.
Its inverse is the restriction of \(\Phi^N\).
\end{theorem}

\begin{proof}
We prove the key comparison on a smooth \(G\)-scheme \(X\).
Put \(Z=X^N\) and \(U=X-Z\).
No stabilizer in \(U\) contains \(N\).
This is equivalent to saying that its Levi class belongs to
\(\F[N]\).
Thus \(L_N\Sigma^\infty_G U_+=0\).

We now prove that the inclusion \(j:Z=X^N\hookrightarrow X\)
becomes an equivalence after applying \(L_N\).

Let \(V=N_{Z/X}\). Since \(X\) and \(Z\) are smooth over \(S\),
there is an \(N\)-equivariant exact sequence
\[
 0\longrightarrow T_{Z/S}
 \longrightarrow T_{X/S}|_Z
 \longrightarrow V
 \longrightarrow 0.
\]
The group \(N\) acts trivially on \(Z\).
Explicit computations on tangent vectors identifies
\(T_{Z/S}\) with \((T_{X/S}|_Z)^N\). 
See also \cite[Theorem~13.1 and \S13.12]{MilneAG}.
Taking invariants is exact because \(N\) is linearly reductive.
Thus \(V^N=0\). So
\(L_N\Sigma^\infty_G(V-Z)_+=0\).

Equivariant homotopy purity
\cite[Theorem~3.23]{Hoyois} identifies
\[
 \Sigma^\infty_G(X/U)
 \simeq \Th_Z(V)
 =\Sigma^\infty_G\bigl(V/(V-Z)\bigr).
\]
The cofiber sequences for \(U\subset X\) and \(V-Z\subset V\)
therefore give equivalences
\[
 L_N\Sigma^\infty_G X_+
   \xrightarrow{\sim} L_N\Th_Z(V),
 \qquad
 L_N\Sigma^\infty_G V_+
   \xrightarrow{\sim} L_N\Th_Z(V),
\]
because both complements are killed by \(L_N\).

Let \(s:Z\to V\) be the zero section.
The purity comparison identifies the composite
\[
 \Sigma^\infty_G Z_+
 \xrightarrow{j_+}\Sigma^\infty_G X_+
 \longrightarrow\Sigma^\infty_G(X/U)
\]
with the composite
\[
 \Sigma^\infty_G Z_+
 \xrightarrow{s_+}\Sigma^\infty_G V_+
 \longrightarrow\Th_Z(V).
\]
After applying \(L_N\), the second composite is an equivalence.
The quotient map from \(X_+\) is also an equivalence after \(L_N\).
Hence \(L_Nj_+\) is an equivalence:
\begin{equation}\label{eq:normal-inclusion}
 L_N\Sigma^\infty_G(X^N)_+
   \xrightarrow{\ \simeq\ }L_N\Sigma^\infty_G X_+.
\end{equation}

We next extend the comparison to representation desuspensions.
The inclusion $A^N\hookrightarrow A$ induces an equivalence
\[
  L_NS^{A^N}\xrightarrow{\simeq}L_NS^A.
\]
Taking tensor inverses in the localized category gives an equivalence
\[
  L_NS^{-A^N}\xrightarrow{\simeq}L_NS^{-A}.
\]
Here and below, representations and spectra with a $Q$-action are
implicitly inflated to $G$.

Let $M=S^{-A}\otimes\Sigma^\infty_GX_+$, where $X$ is a smooth
$G$-scheme over $S$. Since $\Phi^N$ is symmetric monoidal,
\[
  L_N\operatorname{inf}_Q^G\Phi^N(M)
  \simeq
  L_N\bigl(S^{-A^N}\otimes\Sigma^\infty_G(X^N)_+\bigr).
\]
The comparison on $M$ is the composite
\[
\begin{aligned}
  L_N\bigl(S^{-A^N}\otimes\Sigma^\infty_G(X^N)_+\bigr)
  &\xrightarrow{\simeq}
  L_N\bigl(S^{-A}\otimes\Sigma^\infty_G(X^N)_+\bigr)\\
  &\xrightarrow{\simeq}
  L_N\bigl(S^{-A}\otimes\Sigma^\infty_GX_+\bigr)
  =L_NM.
\end{aligned}
\]
The first arrow uses the equivalence of inverse spheres.
The second is obtained by tensoring
\eqref{eq:normal-inclusion} with $S^{-A}$.

Therefore this extends the
localized comparison on pointed motivic spaces to a natural
equivalence
\[
  L_N\operatorname{inf}_Q^G\Phi^N\simeq L_N.
\]

Conversely \(\Phi^N\operatorname{inf}_Q^G=\operatorname{id}\).
Every generator used to construct \(e_{\F[N]}\) has empty
\(N\)-fixed scheme, so \(\Phi^N(e_{\F[N]})=0\).
It follows that \(\Phi^NL_N=\Phi^N\).
These two comparisons prove the theorem.
\end{proof}

\begin{corollary}\label{cor:geometric-categorical-fixed}
Let \((-)^N\) be the right adjoint of genuine inflation.
Then
\[
 \Phi^N M\simeq(L_NM)^N.
\]
The localization map induces a natural comparison
\(c_N:M^N\to\Phi^N M\).
\end{corollary}
\begin{proof}
For an \(L_N\)-local object \(Z\) and \(V\in\D_Q\), adjunction gives
\[
 \Map(V,Z^N)
 \simeq\Map(\operatorname{inf}_Q^G V,Z)
 \simeq\Map(L_N\operatorname{inf}_Q^G V,Z)
 \simeq\Map(V,\Phi^N Z).
\]
The last equivalence is Theorem~\ref{thm:normal-inflation}.
Set \(Z=L_NM\) and use \(\Phi^N L_N=\Phi^N\).
\end{proof}

\begin{remark}
Theorem~\ref{thm:normal-inflation} uses linear reductivity, smooth
fixed schemes, and purity. These are the geometric inputs that allow
normal subgroups of positive dimension.
For finite groups, compare \cite[Propositions 4.8 and 4.10]{GH}.
\end{remark}

\begin{proposition}[A normal family]\label{prop:normal-family}
For a normal subgroup \(H\triangleleft G\), put \(Q=G/H\).
Let \(EQ_+\) denote the free idempotent in \(\D_Q\).
Then the augmentations to the unit identify
\[
 e_{\F_{\leq H}}\simeq\operatorname{inf}_Q^G EQ_+.
\]
\end{proposition}

\begin{proof}
Put $A=\operatorname{inf}_Q^G EQ_+$ and
$J=J_{\F_{\leq H}}$.
We prove that the two colocalizations have the same image.

First, an inflated free $Q$-scheme has stabilizer $H$.
Since $J$ is a localizing tensor ideal, this gives
\[
  A\otimes\D_G\subseteq J.
\]

For the reverse inclusion, let $X$ be a smooth $G$-scheme
with Levi classes in $\F_{\leq H}$.
The affine replacements described at the beginning of this section
allow us to assume that $X$ is affine over $k$.

We construct local equivariant maps from $X$ to an inflated free
$Q$-scheme.
Choose a faithful representation
$\rho:Q\hookrightarrow\mathrm{GL}(V)$ and write
\[
  q:X\longrightarrow T=X/\!/G.
\]
The stabilizer $K$ of any geometric closed orbit in $X$ is
linearly reductive.
The family condition and normality of $H$ imply $K\subseteq H$.
Thus $K$ acts trivially on the fiber of the equivariant vector
bundle $X\times V$.

By equivariant vector bundle descent
\cite[Theorem~10.3]{AlperGMS}, there is a vector bundle $E$ on $T$
and a $G$-equivariant isomorphism
\[
  q^*E\simeq X\times V.
\]
Choose a Zariski cover $\{T_i\}$ trivializing $E$ and put
$X_i=q^{-1}(T_i)$.
A frame of $E|_{T_i}$ pulls back to a $G$-equivariant map
\[
  X_i\longrightarrow F:=\mathrm{GL}(V)\times_k S,
\]
where $G$ acts on $\mathrm{GL}(V)$ by left multiplication through
$\rho$.
This action factors through a free $Q$-action, so
$\Sigma^\infty_G F_+$ belongs to $A\otimes\D_G$.

The graph and projection exhibit $X_i$ as a retract of
$F\times_S X_i$.
Closure under tensor products and retracts therefore gives
\[
  \Sigma^\infty_G(X_i)_+\in A\otimes\D_G.
\]
The same argument applies to intersections of these opens.
Zariski descent then gives
$\Sigma^\infty_GX_+\in A\otimes\D_G$.

These suspension spectra generate $J$ as a localizing tensor ideal.
Hence $J=A\otimes\D_G$.
Both augmentations are the colocalization onto this full subcategory,
so they canonically identify
\[
  e_{\F_{\leq H}}\simeq\operatorname{inf}_Q^G EQ_+.
\]
\end{proof}

\section{Reconstruction of one isotropy layer}
\label{sec:one-layer}

Fix a linearly reductive \(k\)-subgroup \(H\subset G\), and put
\[
 N_H=N_G(H),\qquad W_H=N_H/H.
\]
Both groups are linearly reductive in characteristic zero.
We may check this after extending \(k\) to an algebraic closure.
Richardson's theorem implies that \(C_G(H)^\circ\) is reductive
\cite[Proposition~10.1.5]{RichardsonOrbits}.
Let \(U=R_u(N_H)\).
The intersection \(U\cap H\) is a normal unipotent subgroup of \(H\),
so it is trivial. Since \(U\) and \(H\) are normal in \(N_H\),
\[
 [U,H]\subset U\cap H=1.
\]
Thus \(U\) is a connected normal unipotent subgroup of
\(C_G(H)^\circ\), hence \(U=1\).
This proves that \(N_H\) is linearly reductive.
Its quotient \(W_H\) is also linearly reductive, since taking
\(W_H\)-invariants is taking \(N_H\)-invariants after inflation.

Induction \(\operatorname{Ind}_{N_H}^G\) means genuine induction.
On suspension spectra it sends \(Y\) to \(G\times_{N_H}Y\).
The homogeneous space \(G/N_H\) is affine because \(N_H\)
is reductive.
This construction has the usual induction--restriction
adjunction and projection formula.

Let \(\F_{\leq H}\) consist of geometric Levi classes subconjugate
to \(H\), and let \(\F_{<H}\) omit the class of \(H\).
Set
\[
 e_H=e_{\F_{\leq H}}\otimes\widetilde e_{\F_{<H}},
 \qquad
 \D_{G,H}=e_H\otimes\D_G.
\]
Equivalently,
\(\D_{G,H}=J_{\F_{\leq H}}/J_{\F_{<H}}\).
Note that this is a layer of a filtration, not a direct factor of \(\D_G\).

\subsection{The geometric input}

\begin{lemma}\label{lem:affine-maximal-locus}
Let \(X\) be an affine smooth \(G\)-scheme in
\(\Sm^G_S[\F_{\leq H}]\).
The natural morphism
\[
 j:G\times_{N_H}X^H\longrightarrow X
\]
is a closed immersion.
Its source is smooth over \(S\).
Its geometric stabilizers are conjugate to \(H\).
The complement has Levi classes in \(\F_{<H}\).
The normal representation at every point of the source has no
fixed vector under its stabilizer.
\end{lemma}

\begin{proof}
We prove that \(j\) is a closed immersion 
by checking this on Luna slice charts.
We first work over an algebraically closed field.

Let \(z\in X\) have closed orbit, and put \(K=G_z\).
Matsushima's criterion implies that \(K\) is reductive
\cite[I.2, Proposition]{Luna}.
The family hypothesis therefore makes \(K\) subconjugate to \(H\).
After conjugation, we may assume \(K\subset H\).

By Luna's slice theorem
\cite[III.1, Theorem, p.~97]{Luna},
there is an invariant affine open neighborhood \(U\subset X\)
of the orbit and a surjective etale morphism
\[
 p:X'=G\times_KV\longrightarrow U.
\]
Write \(Q=U/\!/G\) and \(Q'=X'/\!/G\).
The theorem also gives a \(G\)-equivariant isomorphism
\[
 X'\simeq U\times_QQ',
\]
where \(G\) acts trivially on \(Q\) and \(Q'\).

We first check that the comparison map on \(X'\) is the
pullback of \(j\).
Taking fixed schemes in the preceding identification gives
\[
 (X')^H\simeq U^H\times_QQ'.
\]
Taking associated bundles then gives
\[
 \begin{aligned}
 G\times_{N_H}(X')^H
 &\simeq (G\times_{N_H}U^H)\times_QQ'\\
 &\simeq (G\times_{N_H}U^H)\times_UX'.
 \end{aligned}
\]
Thus the square
\[
 \begin{array}{ccc}
 G\times_{N_H}(X')^H
   &\longrightarrow&G\times_{N_H}U^H\\
 \downarrow j_{X'}&&\downarrow j_U\\
 X'&\xrightarrow{\ p\ }&U
 \end{array}
\]
is Cartesian.
Since \(U\) is \(G\)-invariant, \(j_U\) is the restriction of
\(j\) to \(U\).
It therefore suffices to prove that \(j_{X'}\) is a closed
immersion.

If \(K\subsetneq H\), an \(H\)-fixed geometric point
\([g,v]\in G\times_KV\) would give
\[
 g^{-1}Hg\subset K_v\subset K\subsetneq H.
\]
This is impossible: an inclusion between conjugate smooth
algebraic groups is an equality, since their dimensions and
numbers of connected components agree.
Hence \((X')^H=\varnothing\), and \(j_{X'}\) is a closed
immersion.

Suppose now that \(K=H\).
We will identify \(j_{X'}\) with
\[
 G\times_HV^H\longrightarrow G\times_HV.
\]
For this, we compute \((X')^H\) scheme-theoretically.
Let \(q:G\times V\to X'\) be the quotient torsor.
For every \(k\)-scheme \(T\) and every
\((g,v)\in G(T)\times V(T)\), the quotient relation gives
\[
 q(g,v)\in (X')^H(T)
 \quad\Longleftrightarrow\quad
 g^{-1}H_Tg\subset H_v\subset H_T,
\]
where \(H_v\) is the subgroup scheme fixing \(v\).

On every geometric fiber, the inclusion
\(g^{-1}H_Tg\hookrightarrow H_T\) is an isomorphism.
Consequently,
\[
 g^{-1}H_Tg=H_T.
\]
The fixed-point condition is thus equivalent to
\[
 g\in N_H(T),\qquad v\in V^H(T).
\]
This identifies \(q^{-1}((X')^H)\) with \(N_H\times V^H\).
Taking the quotient by \(H\) yields
\[
 (X')^H\simeq N_H\times_HV^H.
\]

We can now identify the comparison map.
The canonical isomorphisms
\[
 G\times_{N_H}(X')^H
 \simeq G\times_{N_H}(N_H\times_HV^H)
 \simeq G\times_HV^H
\]
send \([g,[n,v]]\) to \([gn,v]\).
Under these identifications, \(j_{X'}\) is the map induced
by \(V^H\hookrightarrow V\).
Therefore this shows that \(j_{X'}\)
is a closed immersion.

For a general characteristic-zero field \(k\), apply the
argument after extension to an algebraic closure.
Fixed schemes and associated bundles commute with this
extension.
Faithfully flat descent then proves that \(j\) is a closed
immersion over \(k\).
\end{proof}

\begin{corollary}\label{cor:affine-layer}
For \(X\) as in Lemma~\ref{lem:affine-maximal-locus}, the morphism
\(j\) induces an equivalence
\[
 e_H\otimes\Sigma^\infty_G(G\times_{N_H}X^H)_+
 \xrightarrow{\ \simeq\ }
 e_H\otimes\Sigma^\infty_GX_+.
\]
\end{corollary}

\begin{proof}
The complement of \(j\) is killed by \(e_H\).
The punctured normal bundle is also killed since its stabilizers
are proper subgroups of the stabilizers along \(j\).
Compare the purity cofiber sequence with the zero-section
cofiber sequence, exactly as in
\eqref{eq:normal-inclusion}.
\end{proof}

\begin{remark}
The affine hypothesis is essential to this geometric lemma.
For \(G=\mathrm{SL}_2\), a maximal torus \(H\), and \(X=G/B\),
the corresponding morphism \(G/H\to G/B\) is an affine-line
bundle, not a closed immersion.
To apply Corollary~\ref{cor:affine-layer} to a general smooth
quasiprojective generator, first use the equivariant Jouanolou
replacement of \cite[Proposition 2.20]{Hoyois}.
Lemma~\ref{lem:family-motivic} preserves its family condition.
\end{remark}

\subsection{Free spectra and the layer equivalence}

\begin{lemma}\label{lem:unipotent-free}
The localizing subcategory generated by smooth schemes with
trivial Levi class equals the subcategory generated by strictly
free smooth schemes.
\end{lemma}

\begin{proof}
Only one inclusion needs proof.
Use an affine equivariant Jouanolou replacement of a generator.
Every geometric stabilizer then has trivial reductive quotient.
A closed orbit in this affine scheme has reductive stabilizer,
so its stabilizer is trivial and its dimension is \(\dim G\).
Every orbit closure contains a closed orbit.
A nonclosed orbit would have strictly greater dimension than
a closed orbit in its boundary.
That would exceed \(\dim G\), which is impossible.
Thus every orbit is closed and every stabilizer is trivial.
In characteristic zero these stabilizer schemes are smooth,
so the action is scheme-theoretically free.
\end{proof}

Let
\[
 i_{W!}:\SH(BW_S)\longrightarrow\SH^W(S)
\]
be the free embedding.
Its image is the localizing subcategory generated by suspension
spectra of strictly free smooth \(W\)-schemes and their ordinary
Tate desuspensions.
For \(H\subset G\), define
\[
 \Phi_g^H=\Phi^{H}\operatorname{res}_{N_H}^G:
 \D_G\longrightarrow\D_{W_H},
 \qquad
 \Phi_{\lis}^H=i_{W_H}^*\Phi_g^H.
\]
Thus restriction to the normalizer is part of the definition.

\begin{theorem}[One isotropy layer]\label{thm:isotropy-layer}
There is an equivalence
\[
 \Phi_{\lis}^H:\D_{G,H}
   \xrightarrow{\ \simeq\ }\SH(BW_{H,S}).
\]
An inverse is
\begin{equation}\label{eq:layer-inverse}
 \Xi_H(M)=
 e_H\otimes
 \operatorname{Ind}_{N_H}^G
 \operatorname{inf}_{W_H}^{N_H} i_{W_H!}M.
\end{equation}
The free embedding in this formula is taken for \(W_H\), before
genuine inflation to \(N_H\).
\end{theorem}

\begin{proof}

We first prove that \(\Phi_g^H\) sends
\(J_{\F_{\leq H}}\) into the free \(W_H\)-subcategory.
Let \(X\in\Sm^G_S[\F_{\leq H}]\).
Choose an equivariant Jouanolou replacement \(X'\to X\)
with \(X'\) affine over \(k\), using the successive replacements
described at the beginning of Section~\ref{sec:isotropy}
if necessary.
Lemma~\ref{lem:family-motivic} shows that \(X'\) still belongs
to \(\Sm^G_S[\F_{\leq H}]\).

By Lemma~\ref{lem:affine-maximal-locus}, the stabilizer in \(G\)
of every geometric point of \((X')^H\) is exactly \(H\).
Thus \(W_H\) acts freely on \((X')^H\), and
\[
 \Phi_g^H(\Sigma^\infty_GX_+)
 \simeq\Sigma^\infty_{W_H}(X^H)_+
 \simeq\Sigma^\infty_{W_H}((X')^H)_+
\]
belongs to the free subcategory.
Hence
\[
 \Phi_g^H(J_{\F_{\leq H}})
 \subset\SH^{W_H}(S)_{\free}.
\]

We next prove that \(\Phi_g^H\) kills \(J_{\F_{<H}}\).
For \(X\in\Sm^G_S[\F_{<H}]\), an \(H\)-fixed geometric point
would have a stabilizer containing \(H\).
Since \(H\) is linearly reductive, it would be conjugate into
a Levi subgroup of that stabilizer.
This contradicts the assumption that the Levi class is
strictly below \(H\).
Therefore \(X^H=\varnothing\), and
\[
 \Phi_g^H(\Sigma^\infty_GX_+)=0.
\]
The two conclusions show that \(\Phi_g^H\) factors through
\[
 J_{\F_{\leq H}}/J_{\F_{<H}}
 \longrightarrow\SH^{W_H}(S)_{\free}.
\]
Identifying the source with \(\D_{G,H}\) and the target with
\(\SH(BW_{H,S})\) gives the functor
\[
 \Phi_{\lis}^H=i_{W_H}^*\Phi_g^H:
 \D_{G,H}\longrightarrow\SH(BW_{H,S}).
\]

We now compute geometric fixed points of induced free generators.
Let \(Y\) be a strictly free smooth \(W_H\)-scheme, viewed as
an \(N_H\)-scheme by inflation.
Its stabilizers in \(N_H\) are exactly \(H\).
Consequently,
\[
 Z_Y=G\times_{N_H}Y
\]
has stabilizers conjugate to \(H\).
The map
\[
 Y\longrightarrow Z_Y^H,\qquad y\longmapsto[1,y],
\]
is a canonical \(W_H\)-equivariant isomorphism.
Indeed, an \(H\)-fixed geometric point \([g,y]\) satisfies
\[
 H\subset gHg^{-1}.
\]
The same arguments as in Lemma~\ref{lem:affine-maximal-locus} 
show that the inclusion is an equality and \(g\in N_H\).
It identifies
\[
 Z_Y^H\simeq N_H\times_{N_H}Y\simeq Y.
\]

We specify the natural map underlying this fixed-point calculation.
Write
\[
 \mathcal I(E)=
 \operatorname{Ind}_{N_H}^G
 \operatorname{inf}_{W_H}^{N_H}E.
\]
Applying \(\Phi^H\) to the induction--restriction unit gives
\[
 \begin{aligned}
 u_E:E
 &\simeq\Phi^H\operatorname{inf}_{W_H}^{N_H}E\\
 &\longrightarrow
 \Phi^H\operatorname{res}_{N_H}^G
 \operatorname{Ind}_{N_H}^G
 \operatorname{inf}_{W_H}^{N_H}E\\
 &=\Phi_g^H\mathcal I(E).
 \end{aligned}
\]
For \(E=\Sigma^\infty_{W_H}Y_+\), this is the suspension
of the map \(y\mapsto[1,y]\).
The preceding calculation proves that \(u_E\) is an equivalence
on free suspension generators.

We next identify the images of the family idempotents.
Put
\[
 P=\Phi_g^H(e_{\F_{\leq H}}).
\]
The first part of the proof shows that \(P\) is free.
Since \(Z_Y\) has stabilizers conjugate to \(H\), the family
augmentation gives an equivalence
\[
 e_{\F_{\leq H}}\otimes\Sigma^\infty_G(Z_Y)_+
 \xrightarrow{\sim}
 \Sigma^\infty_G(Z_Y)_+.
\]
Applying \(\Phi_g^H\) and using \(Z_Y^H\simeq Y\) gives
\[
 P\otimes\Sigma^\infty_{W_H}Y_+
 \xrightarrow{\sim}
 \Sigma^\infty_{W_H}Y_+.
\]
This map is induced by the augmentation \(P\to\one\).
This characterizes the free idempotent and gives
\[
 \Phi_g^H(e_{\F_{\leq H}})\simeq EW_{H,+},
\]
compatibly with the augmentations.

The vanishing on \(J_{\F_{<H}}\) gives
\[
 \Phi_g^H(e_{\F_{<H}})=0.
\]
Applying \(\Phi_g^H\) to the defining cofiber sequence for
\(\widetilde e_{\F_{<H}}\) therefore yields
\[
 \Phi_g^H(\widetilde e_{\F_{<H}})\simeq\one.
\]
Since
\(e_H=e_{\F_{\leq H}}\otimes\widetilde e_{\F_{<H}}\),
we obtain
\[
 \Phi_g^H(e_H)\simeq EW_{H,+}.
\]

We finally prove that
\(\Phi_{\lis}^H\Xi_H\simeq\operatorname{id}\).
Let \(M\in\SH(BW_{H,S})\), and put \(E=i_{W_H!}M\).
Since \(E\) is free, the free augmentation and full faithfulness
of \(i_{W_H!}\) identify
\[
 i_{W_H}^*(EW_{H,+}\otimes E)
 \simeq i_{W_H}^*E
 \simeq M.
\]
Using \(u_E\) and the calculation of \(\Phi_g^H(e_H)\),
we define the natural comparison
\[
 \begin{aligned}
 M
 &\simeq i_{W_H}^*(EW_{H,+}\otimes E)\\
 &\xrightarrow{\ i_{W_H}^*(1\otimes u_E)\ }
 i_{W_H}^*
 \bigl(EW_{H,+}\otimes\Phi_g^H\mathcal I(E)\bigr)\\
 &\simeq
 i_{W_H}^*\Phi_g^H
 \bigl(e_H\otimes\mathcal I(E)\bigr)\\
 &=\Phi_{\lis}^H\Xi_H(M).
 \end{aligned}
\]

The comparison is an equivalence when \(E\) is a free suspension
generator.
It is therefore an equivalence on \(\SH(BW_{H,S})\).

We have proved the natural equivalence
\[
 \Phi_{\lis}^H\Xi_H\simeq\operatorname{id}.
\]

For the other composite, use the natural maps
\[
 G\times_{N_H}X^H\longrightarrow X.
\]
Corollary~\ref{cor:affine-layer} says that they become
equivalences on the \(H\)-layer for scheme generators.

We spell out the stabilization issue.
On the free \(W_H\)-subcategory put
\[
 \Psi(Y)=e_H\otimes
 \operatorname{Ind}_{N_H}^G\operatorname{inf}_{W_H}^{N_H}Y.
\]
For a \(G\)-representation \(A\), the inclusion \(A^H\hookrightarrow A\)
and the induction projection formula give a natural map
\begin{equation}\label{eq:isotropy-representation-compatibility}
 \Psi(S^{A^H}\otimes Y)\longrightarrow S^A\otimes\Psi(Y).
\end{equation}
This map is an equivalence.
To check it, first take \(Y=\Sigma^\infty_{W_H}Z_+\) for a strictly
free smooth \(W_H\)-scheme \(Z\).
Linear reductivity gives an \(N_H\)-equivariant splitting
\[
 A|_{N_H}=A^H\oplus A_{\mathrm{mov}}.
\]
The cofiber of the sphere inclusion is, up to the invertible
factor \(S^{A^H}\) and a suspension, the motive of
\((A_{\mathrm{mov}}-0)\times Z\).
The inflated action on \(Z\) has stabilizer \(H\).
Since \(A_{\mathrm{mov}}^H=0\), every stabilizer in this product
is a proper subgroup of \(H\).
After induction, its motive is killed by \(e_H\).

The map \eqref{eq:isotropy-representation-compatibility} therefore extends to representation desuspensions:
on \(S^{-A}\otimes\Sigma^\infty_GX_+\).
The stated compatibility makes this extension well defined in
the genuine stabilization.
The objects
\[
 e_H\otimes S^{-A}\otimes\Sigma^\infty_GX_+,
 \qquad X\in\Sm^G_S[\F_{\leq H}],
\]
generate the layer.
The comparisons thus extend through colimits to
\[
 \Xi_H\Phi_{\lis}^H\simeq\operatorname{id}_{\D_{G,H}}.
\]
Together with the first comparison, this proves the theorem.
\end{proof}

\begin{remark}\label{rem:arithmetic-isotropy}
Theorem~\ref{thm:isotropy-layer} applies over any characteristic zero
field to a subgroup \(H\) defined over that field.
Over an algebraically closed field, the family generation result
below shows that the principal categories for subgroups defined
over \(k\) suffice for an exhaustive categorical filtration.
This argument does not require a separate assertion that every
abstract subgroup over a larger algebraically closed field
descends to \(k\).
Over a general field, a complete indexing theorem requires
additional descent data for geometric classes, since the desired subgroups
may not exist over the base field.
\end{remark}

\subsection{Exhaustivity over an algebraically closed field}

\begin{proposition}\label{prop:family-generation}
Assume that \(k\) is algebraically closed.
For every geometric family \(\F\), the category \(J_\F\) is
generated as a localizing subcategory by the principal
subcategories \(J_{\F_{\leq H}}\), for \((H)\in\F\).
The family filtration is exhaustive.
\end{proposition}

\begin{proof}


Let
\[
 \mathcal C=
 \operatorname{Loc}_{\D_G}
 \bigl\{J_{\F_{\leq H}}:[H]\in\F,\ H\subset G
 \text{ a linearly reductive }k\text{-subgroup}\bigr\}.
\]
We must prove that \(\mathcal C=J_\F\).
The family axioms give \(\F_{\leq H}\subseteq\F\) whenever
\([H]\in\F\), so \(\mathcal C\subseteq J_\F\).
To prove the reverse inclusion, it suffices to show that
\[
 \Sigma^\infty_G X_+\in\mathcal C
 \qquad\text{for every }X\in\Sm^G_S[\F].
\]

We may reduce to the case where \(X\) is affine over \(k\)
by the arguments given in the beginning of Section~\ref{sec:isotropy}.
Choose an ordinary affine Jouanolou bundle \(S'\to S\).

We construct invariant open neighborhoods belonging to
principal families.
Let \(x\in X(k)\) have closed orbit, and put \(H=G_x\).
Such a point is defined over \(k\) because \(k\) is algebraically closed.
The stabilizer \(H\) is reductive by
\cite[I.2, Proposition]{Luna}.
Its Levi class is therefore \([H]\), so \([H]\in\F\).

Luna's slice theorem gives an invariant open neighborhood
\(U_x\subset X\) of the orbit such that every geometric
stabilizer on \(U_x\) is conjugate to a subgroup of \(H_x\),
after extension to the corresponding geometric field;
see \cite[III.1, Theorem and Remark~4, pp.~97--98]{Luna}.
Every Levi subgroup of such a stabilizer is therefore
subconjugate to \(H_x\).
Hence
\[
 U_x\in\Sm^G_S[\F_{\leq H_x}].
\]

These neighborhoods cover \(X\).
Indeed, a nonempty invariant closed complement would contain
a closed orbit, obtained by choosing an orbit of minimal
dimension in that complement.
This is impossible, since every closed orbit lies in one
of the chosen neighborhoods.
By quasicompactness, we obtain a finite cover
\[
 X=U_1\cup\cdots\cup U_m,
 \qquad
 U_i\in\Sm^G_S[\F_{\leq H_i}],
 \quad [H_i]\in\F.
\]

Nisnevich excision therefore implies that the filtration is exhaustive.
\end{proof}

The partial order of \(k\)-defined reductive subgroups, modulo
geometric conjugacy, is well founded.
In a strict descending chain, either dimension decreases or,
while the dimension is constant, the number of connected
components decreases.
The dimension can decrease only finitely many times, and
each fixed-dimension segment is finite.
The following corollary is a motivic
instance of the conservativity principle in
\cite[Remark~1.3.10]{AMGR}.

\begin{corollary}[Joint conservativity of the isotropy layers]
\label{cor:isotropy-joint-conservativity}
Under the standing hypotheses, assume that \(k\) is
algebraically closed.
For every geometric family \(\F\), the functors
\[
 \left\{
 \Phi_{\lis}^H:J_\F\longrightarrow\SH(BW_{H,S})
 \right\}_{(H)\in\F}
\]
are jointly conservative.
Here \(H\) runs through the linearly reductive
\(k\)-subgroups of \(G\), up to conjugacy, whose classes
belong to \(\F\).
\end{corollary}

\begin{proof}
Write
\[
 \Gamma_{\leq H}M=e_{\F_{\leq H}}\otimes M,
 \qquad
 \Gamma_{<H}M=e_{\F_{<H}}\otimes M,
 \qquad
 \Delta_HM=e_H\otimes M.
\]
The family localization sequence gives
\[
 \Gamma_{<H}M\longrightarrow
 \Gamma_{\leq H}M\longrightarrow\Delta_HM.
\]

First, we show that principal colocalizations detect zero
objects in every family.
Let \(\mathcal F\) be a geometric family and
\(N\in J_{\mathcal F}\).
Suppose that \(\Gamma_{\leq K}N=0\) for every
\((K)\in\mathcal F\).
For \(A\in J_{\F_{\leq K}}\), the colocalization adjunction gives
\[
 \Map(A,N)\simeq\Map(A,\Gamma_{\leq K}N)=0.
\]
The objects \(A\) satisfying \(\Map(A,N)=0\) form a
localizing subcategory.
By Proposition~\ref{prop:family-generation}, this subcategory
contains \(J_{\mathcal F}\).
Taking \(A=N\) gives \(N=0\).

Next, let \(0\neq M\in J_\F\).
The preceding argument gives a subgroup \(H\) with
\((H)\in\F\) and \(\Gamma_{\leq H}M\neq0\).
Choose such a class minimal in the well-founded
subconjugacy order.
For every \(K<H\), the family intersection identity gives
\[
 \Gamma_{\leq K}(\Gamma_{<H}M)
 \simeq\Gamma_{\leq K}M=0.
\]
Applying the preceding argument to
\(\Gamma_{<H}M\in J_{\F_{<H}}\) shows that
\(\Gamma_{<H}M=0\).
Hence
\[
 \Delta_HM\simeq\Gamma_{\leq H}M\neq0.
\]

Finally, we compare this layer with the stated functor.
The proof of Theorem~\ref{thm:isotropy-layer} gives
\[
 \Phi_g^H(e_{\F_{\leq H}})\simeq EW_{H,+},
 \qquad
 \Phi_g^H(e_{\F_{<H}})=0.
\]
Lisse restriction sends \(EW_{H,+}\) to \(\one\).
Exactness and symmetric monoidality therefore give
\[
 \Phi_{\lis}^H(e_H)\simeq\one,
 \qquad
 \Phi_{\lis}^H(\Delta_HM)
 \simeq\Phi_{\lis}^H(M).
\]
By Theorem~\ref{thm:isotropy-layer},
\(\Phi_{\lis}^H\) is an equivalence on \(\D_{G,H}\).
Since \(\Delta_HM\neq0\), it follows that
\(\Phi_{\lis}^H(M)\neq0\).

Thus the stated functors jointly detect zero objects.
They are exact, so applying this conclusion to the cofiber
of a morphism proves joint conservativity.
\end{proof}


\section{Splitting a filtration indexed by a poset}
\label{sec:splitting}

We now give a criterion for splitting a filtered object as a sum of its
layers. The criterion uses sections of the layer quotients.

The categorical setting comes from stratifications of stable categories.
We use the downset description of such a stratification. This makes the
maps in the filtration explicit. It also separates the splitting argument
from the geometric problem of constructing a particular stratification.
Our conventions agree with \cite[Definitions 2.4.1 and 2.4.3]{AMGR}.

\subsection{The filtration and its layers}

A \emph{downset} in a poset \(P\) is a subset \(D\) such that
\(q\leq p\) and \(p\in D\) imply \(q\in D\).
We write
\[
 P_{\leq p}=\{q:q\leq p\},
 \qquad
 P_{<p}=\{q:q<p\}.
\]
The poset is \emph{well-founded} if every nonempty subset has a minimal
element. In this section all posets are small.

Let \(\mathcal X\) be a presentable stable \(\infty\)-category.
A full subcategory is \emph{localizing} if it is stable and closed under
all small colimits. Given a collection of full subcategories, its
localizing closure is the smallest localizing subcategory containing
them.

\begin{definition}
\label{def:downset-filtration}
A \emph{compatible downset filtration} of \(\mathcal X\) consists of a
full presentable localizing subcategory \(\mathcal X_D\) for each downset
\(D\subseteq P\), with the following properties.
\begin{enumerate}
 \item Inclusions of downsets give inclusions of subcategories.
 We have \(\mathcal X_\varnothing=0\) and \(\mathcal X_P=\mathcal X\).
 \item The inclusion \(i_D:\mathcal X_D\to\mathcal X\) has a right
 adjoint \(c_D\). Set \(\Gamma_D=i_Dc_D\), with its counit
 \(\epsilon_D:\Gamma_D\to\mathrm{id}_{\mathcal X}\).
 Each \(\Gamma_{P_{\leq p}}\) preserves colimits.
 \item The colocalizations satisfy
 \[
  \Gamma_D\Gamma_E\simeq\Gamma_{D\cap E}
 \]
 by natural equivalences compatible with their counits.
 \item The subcategory \(\mathcal X_D\) is the localizing closure of
 the subcategories \(\mathcal X_{P_{\leq p}}\) for \(p\in D\).
\end{enumerate}
We abbreviate \(\Gamma_{P_{\leq p}}\) to \(\Gamma_{\leq p}\), and use
\(\Gamma_{<p}\) in the same way.
\end{definition}

The object \(\Gamma_DX\) is the part of \(X\) supported on \(D\).
This terminology refers to its universal property, written using
mapping spectra:
\[
 \Map_{\mathcal X}(Y,\Gamma_DX)
 \simeq
 \Map_{\mathcal X}(Y,X),
 \qquad Y\in\mathcal X_D.
\]
For \(D\subseteq E\), this property gives the natural map
\(\Gamma_DX\to\Gamma_EX\). Its composite with
\(\Gamma_EX\to X\) is the counit for \(D\).
These maps form a diagram over the poset of downsets.

For \(p\in P\), define the \emph{supported layer}
\begin{equation}
\label{eq:supported-layer}
 \Delta_pX
 =
 \cofib\bigl(\Gamma_{<p}X\longrightarrow\Gamma_{\leq p}X\bigr).
\end{equation}
This definition places the layer in \(\mathcal X_{\leq p}\).
Applying \(\Gamma_{<p}\) to
\eqref{eq:supported-layer} gives zero. Thus \(\Delta_pX\) lies in
the right orthogonal to \(\mathcal X_{<p}\): its mapping spectrum from
any object of \(\mathcal X_{<p}\) is zero.

\begin{lemma}
\label{lem:filtration-excision}
For downsets \(D,E\subseteq P\), the square
\[
\begin{tikzcd}
 \Gamma_{D\cap E}X \arrow[r] \arrow[d]
   & \Gamma_DX \arrow[d]\\
 \Gamma_EX \arrow[r]
   & \Gamma_{D\cup E}X
\end{tikzcd}
\]
is a pushout.
If \((D_\alpha)_{\alpha<\lambda}\) is an increasing family of downsets
indexed by a nonzero ordinal, then the natural map
\begin{equation}
\label{eq:filtration-continuity}
 \colim_{\alpha<\lambda}\Gamma_{D_\alpha}X
 \longrightarrow
 \Gamma_{\bigcup_{\alpha<\lambda}D_\alpha}X
\end{equation}
is an equivalence.
\end{lemma}

\begin{proof}
We first record a consequence of the generation condition.
The functors \(\Gamma_{\leq p}\), for \(p\in D\), are jointly
conservative on \(\mathcal X_D\).
To prove this, let \(C\in\mathcal X_D\) satisfy
\(\Gamma_{\leq p}C=0\) for every \(p\in D\).
Adjunction gives \(\Map(Y,C)=0\) for every
\(Y\in\mathcal X_{\leq p}\).
The class of objects \(Y\) with \(\Map(Y,C)=0\) is localizing.
It therefore contains \(\mathcal X_D\).
Taking \(Y=C\) shows that \(C=0\).
Applying this argument to the cofiber of a map proves joint
conservativity.

Let \(Q\) be the pushout of the upper and left maps in the displayed
square. All its terms belong to \(\mathcal X_{D\cup E}\).
Hence the natural map \(Q\to\Gamma_{D\cup E}X\) can be tested by
\(\Gamma_{\leq p}\) for \(p\in D\cup E\).
These functors are exact. The right adjoints \(c_D\) preserve finite
limits, which are also finite colimits in a stable category.
The inclusions \(i_D\) are exact because the subcategories are stable.
If \(p\in D\), applying \(\Gamma_{\leq p}\) makes both vertical maps
equivalences, by the intersection condition.
The resulting square is a pushout.
For \(p\in E\), both horizontal maps become equivalences instead.
This proves the first claim.
It is the colocalization form of the Mayer--Vietoris statement
\cite[Lemma 3.3.5(1)]{AMGR}.

For the second claim, put \(D=\bigcup_{\alpha<\lambda}D_\alpha\).
The source and target of \eqref{eq:filtration-continuity} lie in
\(\mathcal X_D\).
Fix \(p\in D\), and choose \(\alpha_0\) with \(p\in D_{\alpha_0}\).
For \(\alpha\geq\alpha_0\), the intersection condition identifies
\[
 \Gamma_{\leq p}\Gamma_{D_\alpha}X
 \simeq
 \Gamma_{\leq p}X.
\]
The transition maps become identities under these identifications.
Since \(\Gamma_{\leq p}\) preserves colimits, applying it to
\eqref{eq:filtration-continuity} gives an equivalence.
Joint conservativity now proves the claim.
\end{proof}

\subsection{The splitting criterion}

\begin{definition}
\label{def:layer-sections}
Let \(F:\mathcal X\to\mathcal C\) preserve colimits, where
\(\mathcal C\) is presentable and stable.
A \emph{choice of layer sections} for \(X\) consists of maps
\[
 s_p:F(\Delta_pX)\longrightarrow F(\Gamma_{\leq p}X),
 \qquad p\in P,
\]
together with identifications \(q_ps_p\simeq\mathrm{id}\).
Here \(q_p:F(\Gamma_{\leq p}X)\to F(\Delta_pX)\) is the image
of the quotient map in \eqref{eq:supported-layer}.
\end{definition}

Each section gives a map
\[
 \omega_p:
 F(\Delta_pX)\xrightarrow{s_p}F(\Gamma_{\leq p}X)
 \longrightarrow F(X).
\]

\begin{theorem}
\label{thm:layer-splitting}
Let \(P\) be well-founded, and give \(\mathcal X\) a compatible
downset filtration. Let \(F:\mathcal X\to\mathcal C\) preserve
colimits, with \(\mathcal C\) presentable and stable.
For every object \(X\) with a choice of layer sections, the map
\begin{equation}
\label{eq:layer-splitting}
 \bigoplus_{p\in P}F(\Delta_pX)
 \xrightarrow{\ \bigoplus_p\omega_p\ }F(X)
\end{equation}
is an equivalence.
\end{theorem}

\begin{proof}
We build the filtration one layer at a time and show that the
chosen section splits each new step.

Since \(P\) is small and well-founded, choose a well-ordering
\((p_\alpha)_{\alpha<\kappa}\) in which every strict predecessor
of an element appears before that element. For \(\alpha\leq\kappa\), put
\[
 D_\alpha=\{p_\beta:\beta<\alpha\}.
\]
Each \(D_\alpha\) is a downset, with
\(D_0=\varnothing\) and \(D_\kappa=P\).

For \(p\in D_\alpha\), we have \(P_{\leq p}\subseteq D_\alpha\).
Thus the chosen sections define a map
\[
 \omega_\alpha:
 \bigoplus_{p\in D_\alpha}F(\Delta_pX)
 \longrightarrow F(\Gamma_{D_\alpha}X),
\]
whose restriction to the \(p\)-summand is
\[
 F(\Delta_pX)\xrightarrow{s_p}
 F(\Gamma_{\leq p}X)\longrightarrow F(\Gamma_{D_\alpha}X).
\]
These maps are compatible as \(\alpha\) increases. We prove by
transfinite induction that each \(\omega_\alpha\) is an equivalence.

For \(\alpha=0\), both sides are zero.

Suppose that \(\omega_\alpha\) is an equivalence. Write
\(D=D_\alpha\), \(p=p_\alpha\), and \(D'=D_{\alpha+1}\).
Every strict predecessor of \(p\) already lies in \(D\), while
\(p\) does not. Hence
\[
 D\cap P_{\leq p}=P_{<p},
 \qquad
 D\cup P_{\leq p}=D'.
\]
Lemma~\ref{lem:filtration-excision} gives a pushout square
\[
\begin{tikzcd}
 \Gamma_{<p}X \arrow[r] \arrow[d]
   & \Gamma_{\leq p}X \arrow[d,"u"]\\
 \Gamma_DX \arrow[r]
   & \Gamma_{D'}X .
\end{tikzcd}
\]
The cofibers of the horizontal maps are therefore canonically
equivalent. The upper cofiber is \(\Delta_pX\), so applying
the colimit-preserving functor \(F\) gives a cofiber sequence
\[
 F(\Gamma_DX)\xrightarrow{i}
 F(\Gamma_{D'}X)\xrightarrow{\rho}F(\Delta_pX).
\]
Moreover, \(\rho\circ F(u)=q_p\), where \(q_p\) is the original
layer quotient. Consequently
\[
 t=F(u)\circ s_p:
 F(\Delta_pX)\longrightarrow F(\Gamma_{D'}X)
\]
is a section of \(\rho\), since
\[
 \rho\circ t=q_p\circ s_p\simeq\mathrm{id}.
\]

It follows that the inclusion and this section give an equivalence
\[
 (i,t):
 F(\Gamma_DX)\oplus F(\Delta_pX)
 \xrightarrow{\;\sim\;}F(\Gamma_{D'}X).
\]

Under the decomposition \(D'=D\sqcup\{p\}\), the map
\(\omega_{\alpha+1}\) is precisely
\[
 \biggl(\bigoplus_{q\in D}F(\Delta_qX)\biggr)
       \oplus F(\Delta_pX)
 \xrightarrow{\omega_\alpha\oplus\mathrm{id}}
 F(\Gamma_DX)\oplus F(\Delta_pX)
 \xrightarrow{(i,t)}
 F(\Gamma_{D'}X).
\]
Both arrows are equivalences. This proves the successor step.

Now let \(\lambda\leq\kappa\) be a nonzero limit ordinal.
Since \(D_\lambda=\bigcup_{\alpha<\lambda}D_\alpha\),
the continuity statement in
Lemma~\ref{lem:filtration-excision}, together with preservation
of colimits by \(F\), gives
\[
 F(\Gamma_{D_\lambda}X)
 \simeq
 \colim_{\alpha<\lambda}F(\Gamma_{D_\alpha}X).
\]
Likewise, the direct sum over \(D_\lambda\) is the colimit of
the direct sums over the \(D_\alpha\). Compatibility of the
maps \(\omega_\alpha\) therefore identifies
\[
 \omega_\lambda
 \simeq\colim_{\alpha<\lambda}\omega_\alpha.
\]
It is an equivalence because every \(\omega_\alpha\) in this
colimit is an equivalence.

Finally, \(D_\kappa=P\) and \(\Gamma_PX\simeq X\).
The map \(\omega_\kappa\) is exactly the map
\(\bigoplus_p\omega_p\) in \eqref{eq:layer-splitting},
which proves the theorem.
\end{proof}

\begin{remark}
The theorem concerns sections of
\(F(\Gamma_{\leq p}X)\to F(\Delta_pX)\).
These are maps between a supported part and its layer.
A section of a different map, such as a localization
\(F(X)\to F(L_{\geq p}X)\), is not part of the hypothesis.
There is no formal passage from the latter section to the former.
This distinction matters when \(P\) has incomparable elements.
\end{remark}

\subsection{Tensor filtrations}

Retain the hypotheses of Theorem~\ref{thm:layer-splitting}.
Suppose further that \(\mathcal X\) is presentably symmetric monoidal.
Thus its tensor product preserves colimits in each variable.
Assume that the colocalizations in
Definition~\ref{def:downset-filtration} have the form
\[
 \Gamma_DX=e_D\otimes X,
 \qquad e_D=\Gamma_D\one.
\]
Write
\[
 \widetilde e_D=\cofib(e_D\longrightarrow\one).
\]
The intersection condition gives
\(e_D\otimes e_E\simeq e_{D\cap E}\).
In particular, tensoring \(e_{<p}\to\one\) with \(e_{\leq p}\)
identifies the layers as follows:
\begin{equation}
\label{eq:tensor-layer}
 \Delta_pX
 \simeq
 e_{\leq p}\otimes\widetilde e_{<p}\otimes X.
\end{equation}

\begin{corollary}
\label{cor:tensor-layer-splitting}
Under these hypotheses, suppose that each quotient
\[
 F(e_{\leq p}\otimes X)
 \longrightarrow
 F(e_{\leq p}\otimes\widetilde e_{<p}\otimes X)
\]
has a chosen section.
Then these sections give an equivalence
\[
 F(X)\simeq
 \bigoplus_{p\in P}
 F(e_{\leq p}\otimes\widetilde e_{<p}\otimes X).
\]
\end{corollary}

\begin{proof}
Use \eqref{eq:tensor-layer} in
Theorem~\ref{thm:layer-splitting}.
\end{proof}

The role of geometry in an application is now explicit.
One must construct the compatible filtration, identify its layers,
and construct sections of these particular quotient maps.
For finite groups, the corresponding isotropy filtration and the
motivic tom Dieck splitting are developed in
\cite[Sections 3 and 7, especially Theorem 7.4]{GH}.
In our setting, Sections~\ref{sec:isotropy} and~\ref{sec:one-layer}
provide the filtration and identify its layers.
The next section constructs the supported sections needed to apply
Corollary~\ref{cor:tensor-layer-splitting}.


\section{Normalizer transfer and tom Dieck splitting}
\label{sec:td-splitting}

Let \(k\) be algebraically closed of characteristic zero, let \(G\) be a
linearly reductive affine algebraic group, and let \(S\) be smooth and
quasiprojective over \(k\), with trivial \(G\)-action.
Write \(\D_G=\SH^G(S)\).
We use the family filtration by Levi classes and the normalizer
layer equivalence established above.
Proposition~\ref{prop:family-generation} gives its exhaustive indexing
by the \(k\)-defined stabilizers of closed orbits.

For a linearly reductive subgroup \(H\), put
\[
 N=N_G(H),\qquad W=N/H,\qquad
 F=\operatorname{Ind}_N^G,\quad U=\operatorname{res}_N^G,\quad R=\operatorname{coind}_N^G.
\]
Thus \(F\dashv U\dashv R\). These are functors of genuine categories.
The functor \(U\) is symmetric monoidal, and \(F\) satisfies the
projection formula. Superscripts on family objects distinguish the
ambient groups. In particular,
\[
 e_H^G=e_{\le H}^G\otimes\widetilde e_{<H}^G,\qquad
 e_H^N=e_{\le H}^N\otimes\widetilde e_{<H}^N.
\]
The symbols \(\le H\) and \(<H\) refer to subconjugacy of embedded
Levi subgroups in the indicated ambient group.

We use \(Q=\Phi_N^H:\D_N\to\D_W\) and \(P=QU\).
Let \(q^*:\D_W\to\D_N\) be genuine inflation, and let
\((-)^H\) be its right adjoint. Write
\[
 L_H=\widetilde e_{\mathcal F[H]}^N\otimes-,
 \qquad \mathcal F[H]=\{K:H\nsubseteq K\}.
\]
Normal geometric fixed points identify \(QZ=(L_HZ)^H\).
The pure \(H\)-subcategory of \(\D_N\) is the intersection
of the \(\le H\)-supported subcategory with the \(L_H\)-local subcategory.

\subsection{The induction unit on a family}

Proposition~\ref{prop:normal-family} gives, compatibly with augmentations,
\begin{equation}\label{eq:td-normal-inflation}
 q^*e_W\simeq e_{\le H}^N.
\end{equation}
Here \(e_W\) is the universal free \(W\)-object.

\begin{lemma}\label{lem:td-induction-unit}
For \(C\in\mathcal D_{N,\le H}\), the induction unit induces an
equivalence
\[
 Q\eta_C:QC\xrightarrow{\ \sim\ }PFC.
\]
\end{lemma}

\begin{proof}
Both \(Q\) and \(PF\) preserve colimits, so it suffices to check
the induction unit on family generators.

Let \(Z\) be a smooth affine \(N\)-scheme in the family \(\le H\).
The quotient argument in the proof of
Proposition~\ref{prop:normal-family} shows that its geometric
stabilizers are contained in \(H\).
The transporter calculation in the proof of
Lemma~\ref{lem:affine-maximal-locus} then identifies
\[
 Z^H\xrightarrow{\ \sim\ }(G\times_N Z)^H,
 \qquad z\longmapsto[1,z].
\]
By Lemma~\ref{lem:fixed-scheme-functor}, suspension of this
isomorphism is precisely \(Q\eta_{\Sigma_N^\infty Z_+}\).

Compatibility with Thom twists and extension through Jouanolou
replacement and genuine stabilization follow by the same arguments
as in the proof of Theorem~\ref{thm:isotropy-layer}.
Thus \(Q\eta\) is an equivalence on the family generators,
and preservation of colimits proves the result.
\end{proof}

\subsection{A transfer normalized at the identity coset}

Set \(A=F\one\). Geometric duality makes \(A\) dualizable.
The projection formula and duality give
\[
 RU(M)\simeq DA\otimes M,\qquad
 FU(M)\simeq A\otimes M.
\]
For \(C\in\D_N\), define a natural transfer
\(\tau_C:RC\to FC\) by
\begin{equation}\label{eq:td-trace-transfer}
\begin{split}
 RC&\xrightarrow{R\eta_C}RUFC
       \simeq DA\otimes FC\\
   &\xrightarrow{1\otimes F\eta_C}
       DA\otimes FUF C
       \simeq DA\otimes A\otimes FC
     \xrightarrow{\mathrm{ev}\otimes1}FC.
\end{split}
\end{equation}
The last arrow is evaluation. This construction does not assert
that induction and coinduction are equivalent on the whole category.

\begin{lemma}\label{lem:td-transfer-normalization}
For every \(Z\in\D_G\), the transfer
\[
 P\tau_{e_{\le H}^N\otimes UZ}
\]
is an equivalence. For each dualizable stage \(C_\alpha\) of
Lemma~\ref{lem:new-family-approx}, identify \(PFC_\alpha\) with
\(QC_\alpha\) by the induction unit, and identify \(PRC_\alpha\)
with \(QC_\alpha\) by dualizing the induction unit for \(DC_\alpha\).
Under these identifications, \(P\tau_{C_\alpha}\) is the identity.
\end{lemma}

\begin{proof}
We first prove the assertion on the dualizable family stages of
Lemma~\ref{lem:new-family-approx}. On each stage, we compute the
two maps entering the transfer and show that evaluation gives the
identity. We then allow arbitrary coefficients and pass to the
filtered colimit.

Write
\[
 e_{\le H}^N\simeq\colim_\alpha C_\alpha.
\]
The objects \(C_\alpha\), \(FC_\alpha\), and \(FDC_\alpha\)
are dualizable. For the moment, fix one stage and write \(C=C_\alpha\).

\smallskip
\noindent
\emph{1. Duality and the unit identifications.}
The object \(DC\) belongs to the same family ideal as \(C\).
Indeed, the duality triangle identities exhibit \(DC\) as a
retract of \(DC\otimes C\otimes DC\), and the family ideal is
closed under tensor products with arbitrary objects and retracts.

We will use the following coinduction formula, valid for every
\(Z\in\D_G\):
\[
 R(C\otimes UZ)\simeq D(FDC)\otimes Z.
\]
To verify it, let \(T\in\D_G\). There are natural equivalences
\begin{align*}
 \Map_G(T,D(FDC)\otimes Z)
 &\simeq \Map_G(T\otimes FDC,Z)\\
 &\simeq \Map_G(F(UT\otimes DC),Z)\\
 &\simeq \Map_N(UT\otimes DC,UZ)\\
 &\simeq \Map_N(UT,C\otimes UZ).
\end{align*}
These use duality of \(FDC\), the induction projection formula,
the adjunction \(F\dashv U\), and duality of \(C\), respectively.
Taking \(Z=\one\) gives
\[
 RC\simeq D(FDC).
\]
In particular,
\[
 R(C\otimes UZ)\simeq RC\otimes Z,
 \qquad
 F(C\otimes UZ)\simeq FC\otimes Z.
\]
No dualizability assumption on \(Z\) is needed.

By Lemma~\ref{lem:td-induction-unit}, both maps
\[
 \alpha_C:=Q\eta_C:QC\xrightarrow{\sim}PFC,
 \qquad
 \alpha_{DC}:=Q\eta_{DC}:QDC\xrightarrow{\sim}PFDC
\]
are equivalences.
The symmetric monoidal functors \(P\) and \(Q\) preserve duals.
Thus the second map gives an equivalence
\[
 \beta_C:
 PRC\simeq D(PFDC)
 \xrightarrow{D(\alpha_{DC})}D(QDC)
 \simeq QC.
\]
Our goal on this stage is to prove
\[
 P\tau_C=\alpha_C\circ\beta_C.
\]
This also specifies what it means for the transfer to be the
identity under the unit identifications.

Put \(B=PA\). By Lemma~\ref{lem:fixed-scheme-functor},
\(B\) is the suspension spectrum of
\(((G/N)^H\times S)_+\).
The proof of Lemma~\ref{lem:td-induction-unit} shows that the
identity coset is open and closed in \((G/N)^H\).
It therefore gives decompositions
\[
 B\simeq\one\oplus B',
 \qquad
 DB\simeq\one\oplus DB'.
\]
Write
\[
 j:\one\longrightarrow B
\]
for the inclusion of this summand. Its dual
\(Dj:DB\to\one\) is the projection onto the corresponding
dual summand. The map \(j\) is \(Q\eta_{\one}\).

\smallskip
\noindent
\emph{2. The orbit coaction inserts the unit summand.}
We claim that, under the projection formula,
\[
 P(F\eta_C)=j\otimes\mathrm{id}_{PFC}:
 PFC\longrightarrow B\otimes PFC.
\]
Naturality of the induction unit, applied to
\(\eta_C:C\to UFC\), gives
\[
 UF(\eta_C)\circ\eta_C
 =
 \eta_{UFC}\circ\eta_C.
\]
After applying \(Q\), this becomes
\[
 P(F\eta_C)\circ\alpha_C
 =
 Q(\eta_{UFC})\circ\alpha_C.
\]
Since \(\alpha_C\) is an equivalence, we can cancel it.

For any \(M\in\D_G\), the projection formula identifies
\(\eta_{UM}\) with
\[
 UM\simeq\one\otimes UM
 \xrightarrow{\eta_{\one}\otimes\mathrm{id}}
 UA\otimes UM.
\]
This is the compatibility of the induction unit with the
projection formula. Taking \(M=FC\) and applying \(Q\) gives
\[
 Q(\eta_{UFC})=j\otimes\mathrm{id}_{PFC},
\]
which proves the claim.

\smallskip
\noindent
\emph{3. The coinduced unit has the required unit coordinate.}
We next prove
\[
 (Dj\otimes\mathrm{id}_{PFC})\circ P(R\eta_C)
 =
 \alpha_C\circ\beta_C.
\]
Only this coordinate of \(P(R\eta_C)\) will enter the transfer.

Dualizing \(\eta_C\) gives
\[
 D\eta_C:U(D(FC))\longrightarrow DC.
\]
Naturality of the induction unit for this map says
\[
 UF(D\eta_C)\circ\eta_{U(D(FC))}
 =
 \eta_{DC}\circ D\eta_C.
\]
Apply \(Q\) and use the projection formula. If we write
\[
 \theta_C:=PF(D\eta_C):
 B\otimes D(PFC)\longrightarrow PFDC,
\]
the resulting identity is
\[
 \theta_C\circ
 (j\otimes\mathrm{id}_{D(PFC)})
 =
 \alpha_{DC}\circ D(\alpha_C).
\]

We now dualize this identity.
The dual of
\[
 F(D\eta_C):FU(D(FC))\longrightarrow FDC
\]
is
\[
 R\eta_C:RC\longrightarrow RUFC.
\]
This follows from the natural duality between induction and
coinduction on the dualizable objects in question.
Since \(P\) preserves duals, it follows that
\[
 D\theta_C=P(R\eta_C):
 PRC\longrightarrow DB\otimes PFC,
\]
where we use the symmetry of the tensor product.
Duality reverses composition, so the preceding identity becomes
\[
 (Dj\otimes\mathrm{id}_{PFC})\circ P(R\eta_C)
 =
 \alpha_C\circ D(\alpha_{DC})
 =
 \alpha_C\circ\beta_C.
\]
The last equality uses the definition of \(\beta_C\).
This proves the required coordinate formula.

\smallskip
\noindent
\emph{4. Evaluation gives the normalized identity.}
Apply \(P\) to the definition of the transfer in
\eqref{eq:td-trace-transfer}.
The calculation in Step~2 identifies the resulting composite as
\[
 PRC
 \xrightarrow{P(R\eta_C)}
 DB\otimes PFC
 \xrightarrow{\mathrm{id}_{DB}\otimes j\otimes\mathrm{id}}
 DB\otimes B\otimes PFC
 \xrightarrow{\mathrm{ev}_B\otimes\mathrm{id}}
 PFC.
\]
By the definition of the dual map \(Dj\),
\[
 \mathrm{ev}_B\circ(\mathrm{id}_{DB}\otimes j)=Dj.
\]
Thus evaluation retains precisely the coordinate computed in
Step~3. We obtain
\[
 P\tau_C
 =
 (Dj\otimes\mathrm{id}_{PFC})\circ P(R\eta_C)
 =
 \alpha_C\circ\beta_C.
\]
Both \(\alpha_C\) and \(\beta_C\) are equivalences. Hence
\(P\tau_C\) is an equivalence, and its normalization satisfies
\[
 \alpha_C^{-1}\circ P\tau_C\circ\beta_C^{-1}
 =
 \mathrm{id}_{QC}.
\]

\smallskip
\noindent
\emph{5. Arbitrary coefficients and the filtered colimit.}
Let \(Z\in\D_G\).
We check that, under the formulas from Step~1,
\[
 \tau_{C\otimes UZ}
 =
 \tau_C\otimes\mathrm{id}_Z.
\]

First, the induction unit is compatible with the projection
formula: for any \(K\in\D_N\), the identification
\[
 UF(K\otimes UZ)\simeq UFK\otimes UZ
\]
carries \(\eta_{K\otimes UZ}\) to
\(\eta_K\otimes\mathrm{id}_{UZ}\).
This can be checked by adjunction: both maps correspond to the
identity of \(F(K\otimes UZ)\).

The coinduction formula is compatible with these maps as well.
More explicitly, the canonical map
\[
 RK\otimes Z\longrightarrow R(K\otimes UZ)
\]
is adjoint to
\[
 URK\otimes UZ\longrightarrow K\otimes UZ,
\]
induced by the counit of \(U\dashv R\).
It is natural in \(K\), and Step~1 shows that it is an
equivalence for \(K=C\).
For objects of the form \(UM\), its compatibility is also
expressed by \(RU(M)\simeq DA\otimes M\).
Consequently, the first arrow in the transfer for \(C\otimes UZ\)
is the first arrow for \(C\) tensored with \(Z\).
The same statement for the second arrow follows from the
compatibility of the induction unit above.
The last arrow is evaluation tensored with the identity.
This proves the asserted formula for \(\tau_{C\otimes UZ}\).

Since \(P\) is symmetric monoidal, we conclude that
\[
 P\tau_{C_\alpha\otimes UZ}
 \simeq P\tau_{C_\alpha}\otimes\mathrm{id}_{PZ}
\]
is an equivalence for every \(\alpha\).

Finally,
\[
 e_{\le H}^N\otimes UZ
 \simeq\colim_\alpha(C_\alpha\otimes UZ).
\]
The functors \(F\), \(U\), and \(Q\) preserve colimits.
The functor \(R\) preserves filtered colimits because its left
adjoint \(U\) preserves compact generators.
Naturality of the transfer therefore identifies
\[
 P\tau_{e_{\le H}^N\otimes UZ}
 \simeq
 \colim_\alpha P\tau_{C_\alpha\otimes UZ}.
\]
Every map in this filtered diagram is an equivalence, so its
colimit is an equivalence. This proves the lemma.
\end{proof}

\begin{lemma}\label{lem:td-pure-transfer}
Induction and coinduction carry the pure \(H\)-subcategory of
\(\D_N\) to the pure \(H\)-subcategory of \(\D_G\).
If \(C\) is pure \(H\) and has the form
\(e_{\le H}^N\otimes UZ\), then \(\tau_C:RC\to FC\) is an equivalence.
\end{lemma}

\begin{proof}
Induction carries the \(\le H\)-family to the corresponding
\(G\)-family. Coinduction does also. To see this, write
\(e_{\le H}^G\) as a filtered colimit of dualizable family stages.
For each such stage the dualizable projection formula applies to
\(R\). Passage to the colimit gives
\[
 e_{\le H}^G\otimes RC
 \simeq R(Ue_{\le H}^G\otimes C)=RC
\]
when \(C\) is supported on the \(N\)-family \(\le H\).

Let \(C\) be pure \(H\). It is \(L_H\)-local.
Since \(Q(Ue_{<H}^G)=0\), localization gives
\(Ue_{<H}^G\otimes C=0\). Projection formula consequently gives
\(e_{<H}^G\otimes FC=0\), so \(FC\) is pure \(H\).

If \(T\) is supported on the strict lower \(G\)-family, then
\(QUT=0\). Adjunction and the \(L_H\)-locality of \(C\) give
\[
 \Map_G(T,RC)\simeq\Map_N(UT,C)=0.
\]
Together with the support assertion, this proves that \(RC\)
is pure \(H\). The normalizer stratum equivalence says that
\(P=\Phi_N^HU\) is conservative on the pure \(H\)-subcategory
of \(\D_G\). For \(C=e_{\le H}^N\otimes UZ\), apply
Lemma~\ref{lem:td-transfer-normalization}.
\end{proof}

\subsection{Sections of the actual supported quotients}

A \(G\)-splitting of \(X\in\D_G\) consists, for each \(H\),
of a specified section in the genuine \(W\)-category,
\[
 s_H:PX\longrightarrow(UX)^H
\]
of the canonical map \(c_H:(UX)^H\to PX\), where
\(PX=\Phi_N^H\operatorname{res}_N^G X\).
The comparison is the one in
Corollary~\ref{cor:geometric-categorical-fixed}.

Set
\[
 B_X=e_{\le H}^N\otimes UX,\qquad
 C_X=e_H^N\otimes UX.
\]
There are canonical equivalences
\[
 B_X^N\simeq(e_W\otimes(UX)^H)^W,\qquad
 C_X^N\simeq(e_W\otimes PX)^W.
\]
Indeed, use \eqref{eq:td-normal-inflation}, the dualizable-stage
projection formula for \(H\)-fixed points, and preservation of
colimits by genuine fixed points. The intersection
\((\le H)\cap\mathcal F[H]=(<H)\) identifies the quotient
\(B_X^N\to C_X^N\) with \((1\otimes c_H)^W\).
Thus \(s_H\) gives a section.
\[
 t_H^N:C_X^N\longrightarrow B_X^N.
\]

We now compare \(B_X\to C_X\) with the supported quotient
\(\Gamma_{\le H}^G X\to\Delta_H^G X\), in order to transport
the section \(t_H^N\) to this quotient after taking fixed points.

The augmentation \(B_X\to UX\), followed by the induction counit,
gives
\[
 \alpha_X:FB_X\longrightarrow FUX
              \xrightarrow{\epsilon_X}X.
\]
Induction carries the \(N\)-family \(\le H\) into the
\(G\)-family \(\le H\), so \(FB_X\) is supported on that family.
The universal property of colocalization therefore gives a
canonical factorization
\[
 FB_X\xrightarrow{\bar\alpha_X}
 \Gamma_{\le H}^G X=e_{\le H}^G\otimes X
 \longrightarrow X.
\]

The restricted strict lower \(G\)-family meets the
\(N\)-family \(\le H\) in exactly the strict lower \(N\)-family.
The family tensor identities give
\[
 C_X\simeq e_{\le H}^N\otimes
            U(\widetilde e_{<H}^G\otimes X).
\]
Since \(C_X\) is pure \(H\), Lemma~\ref{lem:td-pure-transfer}
shows that \(FC_X\) and \(RC_X\) are pure \(H\), and that
\(\tau_{C_X}:RC_X\to FC_X\) is an equivalence.

To identify \(FC_X\) with the supported layer of \(X\), consider
the cofiber sequence
\[
 F(e_{<H}^N\otimes UX)\longrightarrow FB_X\longrightarrow FC_X.
\]
Its first term is supported on the strict lower \(G\)-family,
so \(\Delta_H^G\) kills it; its last term is pure \(H\), so
\(\Delta_H^G\) fixes it. Hence the quotient map gives a canonical
equivalence \(\Delta_H^G(FB_X)\simeq FC_X\). Define
\begin{equation}\label{eq:td-induced-layer}
 \beta_X:FC_X
 \simeq\Delta_H^G(FB_X)
 \xrightarrow{\Delta_H^G(\alpha_X)}\Delta_H^G X.
\end{equation}
By Lemma~\ref{lem:td-induction-unit}, the induction unit identifies
\(PFC_X\) with \(QC_X\simeq e_W\otimes PX\).
Under this identification and
\(P\Delta_H^G X\simeq e_W\otimes PX\), the unit--counit identity
makes \(P\beta_X\) the identity.
Since \(P\) is conservative on the pure \(H\)-subcategory,
\(\beta_X\) is an equivalence.

Naturality of \(\tau\) gives the left square below, and the
definition of \(\beta_X\) gives the right square:
\[
\begin{tikzcd}[column sep=large]
 RB_X \arrow[r,"\tau_{B_X}"] \arrow[d]
   & FB_X \arrow[r,"\bar\alpha_X"] \arrow[d]
   & \Gamma_{\le H}^G X \arrow[d]\\
 RC_X \arrow[r,"\tau_{C_X}","\sim"']
   & FC_X \arrow[r,"\beta_X","\sim"']
   & \Delta_H^G X .
\end{tikzcd}
\]
The first two vertical arrows come from \(B_X\to C_X\), and
the last is the supported-layer quotient.
Both bottom arrows are equivalences. Apply \(G\)-fixed points and
use \((RZ)^G=Z^N\). The composite
\[
 C_X^N\xrightarrow{t_H^N}B_X^N
 \xrightarrow{(\tau_{B_X})^G}(FB_X)^G
 \longrightarrow(\Gamma_{\le H}^G X)^G
\]
is a section of the supported-layer quotient after identifying its
source with \((\Delta_H^G X)^G\) by
\((\beta_X\tau_{C_X})^G\).

\begin{theorem}\label{thm:td-general}
Let \(k\) be algebraically closed of characteristic zero,
\(G\) linearly reductive, and \(S\) smooth and quasiprojective
with trivial \(G\)-action. A \(G\)-split spectrum \(X\) has a
natural equivalence
\[
 \bigoplus_{(H)}
 \bigl(i_W^*\Phi_N^H\operatorname{res}_N^G X
             \otimes\mu_W\bigr)_{hW}
 \xrightarrow{\ \sim\ }X^G,
 \qquad N=N_G(H),\quad W=N/H .
\]
The sum is over conjugacy classes of linearly reductive subgroups.
Here \(i_W^*\) is restriction to the free/lisse category,
\(\mu_W\) is the diagonal Adams twist, and
\(M_{hW}=f_{W\sharp}M\) for \(f_W:BW_S\to S\).
Naturality refers to maps respecting the specified splittings.
\end{theorem}

\begin{proof}
The preceding construction gives a section of each actual supported
quotient after \(G\)-fixed points. The subgroup poset is well-founded.
Genuine fixed points preserve colimits, since inflation preserves
compact generators. Theorem~\ref{thm:layer-splitting},
including its limit-ordinal argument, therefore gives
\[
 X^G\simeq\bigoplus_{(H)}(\Delta_H^G X)^G.
\]
The bottom equivalence in the displayed square identifies the
\(H\)-summand with \(C_X^N=(e_W\otimes PX)^W\).
Corollary~\ref{cor:genuine-adams-repaired}, applied to \(W\), identifies this with
\((i_W^*PX\otimes\mu_W)_{hW}\).
Every summand map is the supported map constructed above followed
by the family augmentation to \(X^G\).
\end{proof}

\begin{corollary}\label{cor:td-suspensions}
For a based motivic \(G\)-space \(Y\), there is a canonical equivalence
\[
 (\Sigma_G^\infty Y)^G\simeq
 \bigoplus_{(H)}
 \bigl(i_W^*\Sigma_W^\infty Y^H\otimes\mu_W\bigr)_{hW}.
\]
\end{corollary}

\begin{proof}
The unstable counit \(q^*Y^H\to UY\), after suspension and
adjunction, gives
\(\Sigma_W^\infty Y^H\to(U\Sigma_G^\infty Y)^H\).
Geometric fixed points identify its source with
\(P\Sigma_G^\infty Y\). The composite with \(c_H\) is the identity,
by the counit identity. Apply Theorem~\ref{thm:td-general}.
\end{proof}

\begin{remark}
When \(G\) is split diagonalizable, every geometric subgroup descends
to the ground field: a subgroup of \(D(M)\) is specified by a quotient
of its character group \(M\). The same proof therefore works over
an arbitrary characteristic-zero field for split diagonalizable
\(G\). All normalizers equal \(G\), so the transfer reduces to
the identity and the supported section is directly
\((1\otimes s_H)^{G/H}\).
To justify family generation over this field, take the invariant open
images of Luna slices after algebraic closure.
Each such open, and its subgroup \(H\), is defined after a finite
extension. The subgroup is already defined over \(k\).
Over a finite Galois extension, take the union of the conjugates of
the open. All these opens still have stabilizers in \(\le H\).
Their union descends to an invariant Zariski open over \(k\).
These descended opens cover, so the same Nisnevich gluing argument
applies. This descends open subsets before using motivic categories.
\end{remark}


\subsection{Explicit splittings in positive dimension}

We now combine Theorem~\ref{thm:td-general} with the twist
calculations of Section~\ref{sec:new-twist}.
In particular,
Lemma~\ref{lem:explicit-central-torus-twist} identifies
the twists for groups whose identity components are
central split tori.

\begin{corollary}[Split diagonalizable groups and split tori]
\label{cor:explicit-diagonalizable-splitting}

Let \(k\) be a field of characteristic zero (not necessarily algebraically closed), 
let \(G\) be a split diagonalizable group of positive dimension, and let
\(S\) be smooth and quasiprojective over \(k\), with trivial
\(G\)-action. For each closed subgroup scheme \(H\subseteq G\),
put
\[
 W_H=G/H,\qquad d_H=\dim(G/H).
\]
For every \(G\)-split spectrum \(X\in\SH^G(S)\), there is
an equivalence
\[
 X^G\simeq
 \bigoplus_{H\subseteq G}
 S^{d_H,d_H}\otimes
 \bigl(i_{W_H}^*\Phi_G^H X\bigr)_{hW_H}.
\]
In particular, for a based motivic \(G\)-space \(Y\),
\[
 (\Sigma_G^\infty Y)^G\simeq
 \bigoplus_{H\subseteq G}
 S^{d_H,d_H}\otimes
 \bigl(i_{W_H}^*\Sigma_{W_H}^\infty Y^H\bigr)_{hW_H}.
\]
The sums include all closed subgroup schemes, not only the
connected ones. If \(G=T=\Gm^r\), then \(W_H\) is a split
torus and \(d_H=r-\dim H\).
\end{corollary}

\begin{proof}
We will identify the Weyl groups and their twists,
then substitute them into the tom Dieck splitting.
The character-group description of diagonalizable groups
\cite[Theorem~12.9]{MilneAG} shows that every \(H\) and
\(G/H\) is split diagonalizable. In characteristic zero,
\((G/H)^\circ\) is a split torus of rank \(d_H\).
Since \(G/H\) is commutative,
Lemma~\ref{lem:explicit-central-torus-twist} gives
\[
 \mu_{W_H}\simeq S^{d_H,d_H}_{BW_{H,S}}.
\]
Moreover, \(N_G(H)=G\), and conjugacy does not identify
distinct subgroups.

Apply Theorem~\ref{thm:td-general}, using its extension to
split diagonalizable groups over arbitrary
characteristic-zero fields explained after
Corollary~\ref{cor:td-suspensions}. The smooth projection
formula moves \(S^{d_H,d_H}\) outside \((-)_{hW_H}\),
giving the first equivalence. The suspension-spectrum
identification in Corollary~\ref{cor:td-suspensions}
gives the second.
\end{proof}

\begin{example}[The multiplicative group]
\label{ex:explicit-gm-splitting}
Under the base hypotheses of
Corollary~\ref{cor:explicit-diagonalizable-splitting}, let
\(Y\) be a based motivic \(\Gm\)-space. Then
\[
 \begin{aligned}
 (\Sigma_{\Gm}^\infty Y)^{\Gm}
 \simeq{}&\Sigma^\infty Y^{\Gm}\\
 &\oplus\bigoplus_{n\geq1}S^{1,1}\otimes
 \bigl(i_{\Gm/\mu_n}^*
       \Sigma_{\Gm/\mu_n}^\infty Y^{\mu_n}\bigr)
       _{h(\Gm/\mu_n)}.
 \end{aligned}
\]
Indeed, the character-group classification
\cite[Theorem~12.9]{MilneAG} identifies the proper closed
subgroups of \(\Gm\) as \(\mu_n\). Each quotient \(\Gm/\mu_n\) is a split
one-dimensional torus; the remaining subgroup is
\(H=\Gm\), whose quotient is trivial.
The formula therefore follows from
Corollary~\ref{cor:explicit-diagonalizable-splitting}.
The quotient \(\Gm/\mu_n\) may be identified with \(\Gm\)
through \(t\mapsto t^n\); the action on \(Y^{\mu_n}\)
is the residual quotient action.
\end{example}

\begin{corollary}[Groups with central torus identity component]
\label{cor:explicit-central-torus-splitting}
Let \(k\) be algebraically closed of characteristic zero.
Let \(G\) be a smooth affine linearly reductive group of
positive dimension whose identity component \(T=G^\circ\)
is a central torus. Let \(S\) be smooth and quasiprojective
over \(k\), with trivial \(G\)-action. Put
\[
 N_H=N_G(H),\qquad W_H=N_H/H,
 \qquad d_H=\dim G-\dim H.
\]
For every \(G\)-split spectrum \(X\), there is an equivalence
\[
 X^G\simeq
 \bigoplus_{(H)} S^{d_H,d_H}\otimes
 \bigl(i_{W_H}^*\Phi_{N_H}^H
             \operatorname{res}_{N_H}^G X\bigr)_{hW_H}.
\]
For a based motivic \(G\)-space \(Y\), this specializes to
\[
 (\Sigma_G^\infty Y)^G\simeq
 \bigoplus_{(H)}S^{d_H,d_H}\otimes
 \bigl(i_{W_H}^*\Sigma_{W_H}^\infty Y^H\bigr)_{hW_H}.
\]
The sums run over conjugacy classes of closed linearly
reductive subgroups. This applies, in particular, to
\(G=\Gm^r\times F\), with \(r>0\) and \(F\) a finite
group, without assuming that \(F\) is abelian.
\end{corollary}

\begin{proof}
Our goal is to verify that each Weyl group satisfies
Lemma~\ref{lem:explicit-central-torus-twist}.
Since \(T\) is central in \(G\), it normalizes \(H\), so
\(T\subseteq N_H\). The connectedness of \(T\) then gives
\(T\subseteq N_H^\circ\). Conversely, the inclusion
\(N_H\subseteq G\) gives \(N_H^\circ\subseteq G^\circ=T\).
Thus \(N_H^\circ=T\).
The image \(\overline T\) of \(T\) in \(W_H\) is
\[
 \overline T\simeq T/(T\cap H).
\]
It is a central split torus by the character-group
description \cite[Theorem~12.9]{MilneAG}.
Since \(\overline T\) is connected, it lies in
\(W_H^\circ\). Conversely, the quotient
\(W_H/\overline T\) is a quotient of the finite \'etale
group \(N_H/T=N_H/N_H^\circ\), and hence is finite \'etale.
The connected group \(W_H^\circ\) therefore maps to the
identity point of \(W_H/\overline T\), so it is contained
in the kernel \(\overline T\). Thus
\(W_H^\circ=\overline T\).
Also \(H^\circ\subseteq T\), so
\[
 \dim W_H^\circ
 =\dim T-\dim(T\cap H)
 =\dim G-\dim H=d_H.
\]
Lemma~\ref{lem:explicit-central-torus-twist} now gives
\(\mu_{W_H}\simeq S^{d_H,d_H}\).
Substitute this in Theorem~\ref{thm:td-general}
and Corollary~\ref{cor:td-suspensions}, and apply the
smooth projection formula to each summand.
\end{proof}


\appendix
\section{A mixed base-change theorem}\label{sec:mixed-bc}

We record a base change theorem for a fourth adjoint separately.
It applies to the classifying stack projections.

Let \(g:U\to S\) be smooth and of finite presentation.
Write
\begin{equation}\label{eq:mixed-square}
\begin{array}{ccc}
 BG_U&\xrightarrow{p_2}&BG_S\\
 {\scriptstyle p_1}\downarrow&&\downarrow{\scriptstyle f}\\
 U&\xrightarrow{g}&S.
\end{array}
\end{equation}
The canonical mixed exchange
\begin{equation}\label{eq:mixed-exchange}
 \delta_G:f_\sharp p_{2*}\longrightarrow g_*p_{1\sharp}
\end{equation}
is adjoint to
\[
 g^*f_\sharp p_{2*}M
 \simeq p_{1\sharp}p_2^*p_{2*}M
 \longrightarrow p_{1\sharp}M.
\]
The final arrow is the counit for \(p_2\).

The functors \(g_*\) and \(p_{2*}\) preserve colimits.
So they have further right adjoints, denoted \(g^\sharp\) and \(p_2^\sharp\).
For any map \(a\) with this adjoint, put
\[
 \sigma_a=a^\sharp\one.
\]
Note that this notation refers to a right adjoint of \(a_*\), 
not of \(a_!\).

\begin{theorem}\label{thm:mixed-exchange}
The map \(\delta_G\) is an equivalence.
Its right adjoint transformation is the equivalence
\[
 p_1^*g^\sharp Y\xrightarrow{\ \sim\ }p_2^\sharp f^*Y
 \qquad(Y\in\SH(S)).
\]
In particular,
\[
 p_1^*\sigma_g\simeq\sigma_{p_2}.
\]
\end{theorem}
\begin{proof}
We follow the comparison-square argument of
\cite[Proposition~6.13]{GH}, using Theorem~\ref{thm:lisse-adams}.
That theorem identifies the continuous direct image. We first give
an argument with compact dualizable coefficients that allows us to
use the ordinary direct image; this plays the role of the argument
in \cite[Lemma~6.14]{GH}.

Let \(A\in\SH(BG_S)\) be compact and dualizable.
Its dual \(DA\) is also compact: the duality identities exhibit it
as a retract of \(DA\otimes A\otimes DA\), and tensoring with a
dualizable object preserves compact objects.
We claim that
\[
 f_*(A\otimes-):\SH(BG_S)\longrightarrow\SH(S)
\]
preserves colimits.
Its left adjoint is \(DA\otimes f^*(-)\).
For a compact generator \(v_\sharp\one_V\) of \(\SH(S)\), with
\(v:V\to S\) smooth, let \(v_G:BG_V\to BG_S\) be its base change.
Smooth base change and the projection formula give
\[
 DA\otimes f^*(v_\sharp\one_V)
 \simeq (v_G)_\sharp v_G^*DA.
\]
This is compact by Lemma~\ref{lem:compact-pullback}.
Thus \(DA\otimes f^*(-)\) preserves compact objects. 
So \(f_*(A\otimes-)\), as its right adjoint, preserves all colimits.

For compact \(N\), the object \(A\otimes N\) is compact, so
Theorem~\ref{thm:lisse-adams} and Lemma~\ref{lem:continuous} give
\begin{equation}\label{eq:mixed-supported-adams}
 \nu_f(A\otimes N):
 f_\sharp(A\otimes N\otimes\mu_f)
 \xrightarrow{\ \sim\ }f_*(A\otimes N).
\end{equation}
Both sides preserve colimits in \(N\), so this equivalence holds
for every \(N\in\SH(BG_S)\).
The same argument applies to \(p_1:BG_U\to U\).

We may express arbitrary coefficients using such objects.
Choose a faithful representation \(G\hookrightarrow GL_m\), and put
\[
 j:(GL_m/G)_S\longrightarrow BG_S,\qquad A=j_\sharp\one.
\]
As explained in Section~\ref{sec:formalism}, this smooth atlas has
Nisnevich local sections. The object \(A\) is compact.
It is also dualizable by Proposition~\ref{prop:geometric-duality}.
The motives of the terms of its \v{C}ech nerve are
\[
 A_n=A^{\otimes(n+1)}\qquad(n\geq0).
\]
They are compact and dualizable. Descent for this atlas gives
\begin{equation}\label{eq:mixed-cech}
 \colim_{[n]\in\Delta^{\mathrm{op}}}A_n\simeq\one_{BG_S},
 \qquad
 M\simeq\colim_{[n]\in\Delta^{\mathrm{op}}}
       (p_2^*A_n\otimes M)
 \quad(M\in\SH(BG_U)).
\end{equation}
Here \(\Delta\) is the simplex category.
The objects \(p_2^*A_n\) are again compact and dualizable.

We now compare the mixed exchange with the Adams maps.
Theorem~\ref{thm:new-mu-invertible} makes \(\mu_f\) invertible, and
smooth base change in its diagonal definition gives
\(\mu_{p_1}\simeq p_2^*\mu_f\).
For \(a=f,p_1\), write the ordinary norm in the form
\[
 t_a(N)=\nu_a(N\otimes\mu_a^{-1}):
 a_\sharp N\longrightarrow a_*(N\otimes\mu_a^{-1}).
\]
For every \(M\in\SH(BG_U)\), there is a commutative square
\begin{equation}\label{eq:mixed-adams-square}
\begin{tikzcd}[column sep=large,row sep=large]
 f_\sharp p_{2*}M
   \arrow[r,"\delta_G(M)"]
   \arrow[d,"t_f(p_{2*}M)"']
 & g_*p_{1\sharp}M
   \arrow[d,"g_*t_{p_1}(M)"] \\
 f_*(p_{2*}M\otimes\mu_f^{-1})
   \arrow[r,"\sim"']
 & g_*p_{1*}(M\otimes\mu_{p_1}^{-1}).
\end{tikzcd}
\end{equation}
The bottom equivalence is the projection formula for the invertible
object \(\mu_f^{-1}\), followed by
\(f_*p_{2*}\simeq g_*p_{1*}\).
Apply the square to \(M=p_2^*A_n\otimes N\).
The dualizable projection formula gives
\[
 p_{2*}M\simeq A_n\otimes p_{2*}N.
\]
Consequently both vertical arrows are equivalences by
\eqref{eq:mixed-supported-adams}, applied over \(S\) and \(U\)
and with the inverse twists included in the coefficient.
Thus \(\delta_G\) is an equivalence on every \(p_2^*A_n\otimes N\).
Note that the functors \(g_*\) and \(p_{2*}\) preserve colimits.
Hence both functors in \eqref{eq:mixed-exchange} preserve colimits.
The \v{C}ech resolution \eqref{eq:mixed-cech} proves that
\(\delta_G\) is an equivalence for every coefficient.

For the final statements, the right adjoints of \(f_\sharp p_{2*}\) 
and \(g_*p_{1\sharp}\) are
\(p_2^\sharp f^*\) and \(p_1^*g^\sharp\), respectively.
The right adjoint transformation is therefore an equivalence
\[
 p_1^*g^\sharp Y\xrightarrow{\ \sim\ }p_2^\sharp f^*Y.
\]
Setting \(Y=\one_S\) gives
\(p_1^*\sigma_g\simeq\sigma_{p_2}\).
\end{proof}



\section{Invertibility of the twist in positive characteristic}
\label{subsec:new-charp-twist}

We now remove the characteristic zero hypothesis from twist invertibility.
This argument uses a compactification of the group.

\begin{lemma}[A compactification criterion]\label{lem:new-boundary-duality}
Let \(J\) be a smooth affine linearly reductive group over a field.
Let \(U\) be a smooth \(J\)-scheme.
Suppose \(U\) has a smooth projective \(J\)-equivariant compactification \(X\).
Suppose the reduced boundary is a finite union
\[
 D=X\setminus U=D_1\cup\cdots\cup D_m
\]
of \(J\)-stable closed subschemes.
Assume every nonempty intersection
\(D_I=\bigcap_{i\in I}D_i\), for \(\varnothing\ne I\subseteq\{1,\ldots,m\}\),
is smooth and proper over the field.
Assume these schemes are \(J\)-quasiprojective.
Then the suspension spectrum of \(U_+\) is dualizable in the genuine category.
The representable motive
\[
 [U/J]\longrightarrow BJ
\]
is also dualizable in the lisse category.
\end{lemma}
\begin{proof}
Write \(q:X\to\Spec k\), \(u:U\to\Spec k\), and \(d:D\to\Spec k\).
Keep the actual equivariant Thom coefficient
\(\Theta=\Th(T_{X/k})\).
Smooth proper ambidexterity implies that \(q_*\Theta\) is dualizable.
It also implies that
\[
 (D_I\to\Spec k)_*(\Theta|_{D_I})
\]
is dualizable.
Indeed, for a smooth proper map \(a\) and a dualizable coefficient \(C\),
closed adjunction and the proper projection formula give
\[
 \Hom(a_\sharp C,Z)
 \simeq a_*(DC\otimes a^*Z)
 \simeq a_*DC\otimes Z.
\]
Thus \(a_\sharp C\) is dualizable.
The formula \(a_*=a_\sharp(-\otimes\Th(-T_a))\) proves the assertion for
\(a_*\) as well.

Proper excision expresses \(d_*(\Theta|_D)\) by finitely many fibers and
finite sums of these objects.
More explicitly, for a union \(A\cup B\) of two closed subschemes, use the
proper cover \(A\amalg B\to A\cup B\).
It is an isomorphism away from \(A\cap B\).
The resulting proper-excision square gives the fiber sequence
\[
 (A\cup B)_*(\Theta|_{A\cup B})
 \longrightarrow
 A_*(\Theta|_A)\oplus B_*(\Theta|_B)
 \longrightarrow
 (A\cap B)_*(\Theta|_{A\cap B}),
\]
where each subscript denotes its structure map to \(\Spec k\).
Induction on the number of components proves the assertion.
Dualizable objects are closed under these finite operations.

Localization on \(X\), followed by \(q_*=q_!\), now gives
\[
 u_\sharp\one
 \longrightarrow q_*\Theta
 \longrightarrow d_*(\Theta|_D).
\]
The first term has this form because
\(u_!\Th(T_{U/k})=u_\sharp\one\).
It is dualizable since the other two terms are dualizable.
The lisse case is similar, 
whose proof uses the representable six operations over \(BJ\).
\end{proof}

\begin{lemma}[A boundary preserved by the finite symmetries]
\label{lem:new-strict-toric-boundary}
Let \(T\) be a torus over \(k\), and let \(F\) be a finite \'{e}tale group
acting on \(T\) by group automorphisms.
There is a smooth projective compactification \(T\subset X\) on which
\(T\rtimes F\) acts.
Its boundary has a finite decomposition into \(T\rtimes F\)-stable smooth
proper divisors whose nonempty intersections are smooth and proper.
Moreover, \(X\) has an ample \(T\rtimes F\)-linearized line bundle.
\end{lemma}
\begin{proof}
We first construct the compactification over a splitting field,
keeping track of both the finite symmetries and the Galois action.
We then arrange the boundary into invariant smooth divisors and
descend the construction to \(k\).

\smallskip
\noindent\emph{Step 1: combine the finite symmetries and the Galois action.}
Choose a finite Galois extension \(L/k\) splitting both \(T\) and \(F\).
Put
\[
 \Gamma=\operatorname{Gal}(L/k),\qquad
 F_0=F(L),\qquad
 N=X_*(T_L).
\]
The action of \(F\) on \(T\) gives an action of \(F_0\) on \(N\).
Because the original action is defined over \(k\), this action
is compatible with the Galois actions on \(F_0\) and \(N\).
Thus \(F_0\rtimes\Gamma\) acts on \(N\).
Let
\[
 A=\operatorname{im}\bigl(F_0\rtimes\Gamma
                 \longrightarrow\operatorname{GL}(N)\bigr).
\]
This is a finite group. We will choose the toric construction
to be invariant under \(A\), so that it respects all the
symmetries needed for descent.

\smallskip
\noindent\emph{Step 2: choose a smooth projective toric compactification
with separated boundary components.}
Recall that a fan in \(N_{\mathbb R}=N\otimes_{\mathbb Z}\mathbb R\)
is a finite collection of rational polyhedral cones, closed under
taking faces, such that two cones meet along a common face.
A fan specifies how affine toric charts are glued together.

The construction of Colliot-Th\'el\`ene, Harari, and Skorobogatov
gives a smooth complete projective \(A\)-invariant fan \(\Sigma\)
with the following property:
\begin{equation}\label{eq:new-separated-fan}
 \rho,\ a\rho\subseteq\sigma
 \quad\Longrightarrow\quad
 a\rho=\rho
 \qquad
 \bigl(a\in A,\ \rho\in\Sigma(1),\ \sigma\in\Sigma\bigr).
\end{equation}
Here \(\Sigma(1)\) denotes the set of one-dimensional cones,
called rays.
Thus two distinct rays in the same \(A\)-orbit never lie in
a common cone.
For the construction, see
\cite[Theorem~1, Proposition~2, and the proof of Proposition~3]{CTHS}.
Proposition~2 supplies the separation property, and the proof
of Proposition~3 verifies that it is preserved while making
the fan smooth.

Let \(Y=X_{\Sigma,L}\) be the associated toric variety.
It is a smooth projective \(L\)-scheme containing \(T_L\) as
a dense open subscheme.
The action of \(T_L\) by translations extends to \(Y\).
Functoriality of the toric construction also extends the
\(F_0\)-action and the semilinear \(\Gamma\)-action.
Their compatibility is inherited from their actions on \(T_L\).

We recall explicitly the boundary geometry that will be used.
A smooth cone of dimension \(d\) is generated by \(d\) members
of a basis of \(N\).
Its affine chart consequently has coordinate ring
\[
 L[x_1,\ldots,x_d,x_{d+1}^{\pm1},\ldots,x_r^{\pm1}],
\]
and its boundary is the union of the coordinate hyperplanes
\(x_1=0,\ldots,x_d=0\).
Globally, the boundary prime divisors \(E_\rho\) are indexed
by the rays \(\rho\in\Sigma(1)\).
A collection of these divisors can meet only when their rays
belong to a common cone.
On the corresponding charts, their intersections are obtained
by setting distinct coordinates equal to zero.
In particular, every such nonempty intersection is smooth.
These coordinate descriptions work over \(L\) in any
characteristic.

\smallskip
\noindent\emph{Step 3: group the boundary components into invariant
smooth divisors.}
For each \(A\)-orbit \(O\subseteq\Sigma(1)\), define
\[
 D_{O,L}=\bigcup_{\rho\in O}E_\rho.
\]
The divisors occurring in this union are pairwise disjoint.
Indeed, if \(E_\rho\) met \(E_{a\rho}\), their rays would belong
to a common cone, and \eqref{eq:new-separated-fan} would force
\(a\rho=\rho\).
Thus \(D_{O,L}\) is a disjoint union of smooth divisors.
It is therefore smooth, and it is proper because it is closed
in \(Y\).

These \(D_{O,L}\) are stable under translations by \(T_L\),
under \(F_0\), and under the semilinear Galois action.
They cover \(Y\setminus T_L\).

We also check the intersections required in
Lemma~\ref{lem:new-boundary-duality}.
For distinct orbits \(O_1,\ldots,O_s\), there is a decomposition
\[
 D_{O_1,L}\cap\cdots\cap D_{O_s,L}
 =
 \coprod_{\substack{\rho_i\in O_i\\1\leq i\leq s}}
 \bigl(E_{\rho_1}\cap\cdots\cap E_{\rho_s}\bigr),
\]
with empty terms omitted.
The union is disjoint since two different choices differ in some
orbit \(O_i\).
Each term is smooth by the coordinate description above.
Hence every nonempty intersection of the \(D_{O,L}\) is
smooth and proper.

\smallskip
\noindent\emph{Step 4: construct a linearized ample line bundle
and descend to \(k\).}
Choose an ample line bundle \(\mathcal L_0\) on \(Y\).
Since \(T_L\) is a split torus, its coordinate ring is a
Laurent polynomial ring and \(\operatorname{Pic}(T_L)=0\).
A trivialization of \(\mathcal L_0|_{T_L}\) therefore gives
a rational section of \(\mathcal L_0\) whose divisor \(B_0\)
is supported on \(Y\setminus T_L\).
Thus
\[
 \mathcal L_0\simeq\mathcal O_Y(B_0),
 \qquad
 B_0=\sum_{\rho\in\Sigma(1)}b_\rho E_\rho.
\]
In particular, \(B_0\) is a torus-invariant Cartier divisor.

The lattice action of \(A\) gives \(L\)-automorphisms of \(Y\).
Set
\[
 B=\sum_{a\in A}a^*B_0.
\]
Then \(B\) is invariant under \(A\) and \(T_L\), and
\[
 \mathcal O_Y(B)
 \simeq\bigotimes_{a\in A}a^*\mathcal L_0
\]
is ample.
\(B\) is also invariant under the semilinear
\(\Gamma\)-action since its permutation of the boundary divisors
factors through \(A\).

Local sections of \(\mathcal O_Y(B)\) are rational functions \(f\)
with \(\operatorname{div}(f)+B\geq0\), together with zero.
Invariance of \(B\) ensures that pullback preserves this condition.
For the action and projection
\(\mu,p_2:(T_L\rtimes F_0)\times_LY\to Y\), the equality
\(\mu^*B=p_2^*B\) therefore gives
\[
 \mu^*\mathcal O_Y(B)\simeq p_2^*\mathcal O_Y(B).
\]
Composition of pullbacks gives the cocycle identity, so this is a
\(T_L\rtimes F_0\)-linearization.
The same construction for \(\Gamma\) gives Galois descent data
compatible with this linearization, again by functoriality of pullback.
See \cite[Tags~03LE and~0CDQ]{Stacks}.

Galois descent for the projective variety \(Y\), as in
\cite[proof of Corollary~1]{CTHS}, now gives a compactification
\(T\subset X\) over \(k\).
The group action, the invariant divisors \(D_{O,L}\), and
the linearized line bundle descend as well.
Write \(D_O\subset X\) and \(\mathcal L\) for their descents.
Smoothness, properness, and ampleness can be checked after
the field extension \(L/k\).
Thus \(X\) is smooth and projective,
\(\mathcal L\) is ample and \(T\rtimes F\)-linearized, and
\[
 X\setminus T=\bigcup_O D_O
\]
is a union of \(T\rtimes F\)-stable smooth proper divisors
whose nonempty intersections are smooth and proper.
This proves all the assertions.
\end{proof}

\begin{proposition}[A compactification with both group actions]
\label{prop:new-torus-finite-compactification}
Suppose \(G\) is a smooth affine linearly reductive group fitting into
an exact sequence
\[
 1\longrightarrow T\longrightarrow G\longrightarrow F\longrightarrow1,
\]
where \(T\) is a torus and \(F\) is finite \'{e}tale.
Then \(G\), with the action \((a,b)\cdot g=agb^{-1}\), satisfies
Lemma~\ref{lem:new-boundary-duality} for \(J=G\times G\).
\end{proposition}
\begin{proof}
Since \(T\) is commutative, conjugation on \(T\) factors through \(F\).
Apply Lemma~\ref{lem:new-strict-toric-boundary} to this action.
Let \(X\) be the resulting compactification and define
\[
 \overline G=G\times^T X.
\]
It is a bundle over the finite \'{e}tale scheme \(G/T=F\).
Its open subscheme \(G\times^T T\) identifies with \(G\) by
\([g,t]\mapsto gt\).

Left multiplication acts on the first coordinate.
Write \(c_h\) for the extension to \(X\) of conjugation by \(h\).
Right multiplication extends by
\[
 [g,x]\cdot h=[gh,c_{h^{-1}}(x)].
\]
To check well-definedness, recall the relation \([gt,x]=[g,tx]\).
Since \(h^{-1}th\in T\) and
\(c_{h^{-1}}(tx)=(h^{-1}th)c_{h^{-1}}(x)\), we have
\[
 \begin{aligned}
 {[gt,x]}\cdot h
 &=[gth,c_{h^{-1}}(x)]\\
 &=[gh,(h^{-1}th)c_{h^{-1}}(x)]\\
 &=[gh,c_{h^{-1}}(tx)]=[g,tx]\cdot h.
 \end{aligned}
\]

Using inverse right multiplication for the second factor therefore gives
\[
 (a,b)\cdot[g,x]=[agb^{-1},c_b(x)],
\]
extending the required action \((a,b)\cdot g=agb^{-1}\) on \(G\).

Pulling back \(\overline G\to F\) along the \(T\)-torsor \(G\to F\) gives
\[
 G\times_F\overline G\simeq G\times_kX.
\]
The projection to \(G\) is smooth and proper because \(X/k\) is.
These properties descend to \(\overline G\to F\)
by \cite[Tag~02YJ]{Stacks}.
Since \(F/k\) is finite \'{e}tale, \(\overline G\) is smooth and proper over \(k\).

Let \(\mathcal L\) be the ample linearized line bundle from
Lemma~\ref{lem:new-strict-toric-boundary}.
Its \(T\)-linearization gives the associated line bundle
\(\overline{\mathcal L}=G\times^T\mathcal L\) on \(\overline G\).
The compatible lifts \(\widetilde c_b\) of conjugation give a
\(G\times G\)-linearization by
\[
 (a,b)\cdot[g,v]=[agb^{-1},\widetilde c_b(v)].
\]
Well-definedness follows from the same calculation as for points.

Over the cover \(G\to F\), this line bundle becomes
\(\operatorname{pr}_X^*\mathcal L\), which is relatively ample over \(G\).
Relative ampleness descends by \cite[Tag~0D2P]{Stacks}.
Since \(F\) is affine, \(\overline{\mathcal L}\) is therefore ample on
\(\overline G\).
Properness now implies projectivity \cite[Tag~0B44]{Stacks}.
A sufficiently high tensor power of \(\overline{\mathcal L}\) gives an
equivariant projective embedding, so \(\overline G\) is also
\(G\times G\)-quasiprojective.

The boundary pieces are \(G\times^T D_O\).
They are invariant under both actions.
Their nonempty intersections are \(G\times^T(\bigcap_O D_O)\).
They are smooth and proper by torsor descent and the finite \'{e}tale
base \(F\).
The restricted ample linearization proves their equivariant
quasi-projectivity as well.
These are all the hypotheses of Lemma~\ref{lem:new-boundary-duality}.
\end{proof}

\begin{theorem}[Positive-characteristic twist invertibility]
\label{thm:new-charp-invertibility}
Let \(k\) have characteristic \(p>0\), and let \(G/k\) be smooth,
affine, and linearly reductive.
Let \(S\) be smooth and quasiprojective over \(k\), with trivial
\(G\)-action. For \(f:BG_S\to S\), the object
\(\Delta_\sharp\one\) is dualizable and
\[
 \mu_f=q_{1*}\Delta_\sharp\one
 \quad\text{is invertible in }\SH(BG_S).
\]
Its inverse is \(q_{1\sharp}D(\Delta_\sharp\one)\).
No splitting assumption and no inversion of \(p\) are required.
\end{theorem}
\begin{proof}
Nagata's classification says that \(G^\circ\) is of multiplicative type
and that \(|\pi_0(G)(\overline{k})|\) is prime to \(p\).
For the statement over an arbitrary field, see
\cite[Definition~1.1 and Theorem~1.2]{HallRydhGroups}.
Since \(G\) is smooth, \(G^\circ\) is a torus.
Its quotient is finite \'{e}tale.
Proposition~\ref{prop:new-torus-finite-compactification} and
Lemma~\ref{lem:new-boundary-duality} therefore apply.

The quotient of \(G\), with its left-right action, by \(G\times G\)
is \(BG\).
The resulting morphism to \(B(G\times G)\) is the diagonal.
Thus the lisse conclusion of Lemma~\ref{lem:new-boundary-duality}
is precisely dualizability of \(\Delta_\sharp\one\).
Pullback gives the same statement over \(S\).
Theorem~\ref{thm:new-mu-invertible} now proves the stated formula
for the inverse.
\end{proof}

The theorem proves invertibility of the diagonal twist.
An Adams norm equivalence in that generality requires a separate argument.

\enlargethispage{2\baselineskip}
\bibliographystyle{amsalpha}
\bibliography{references}
\end{document}
